\documentclass[11pt]{article}
\usepackage[margin=1in]{geometry}
\usepackage{amsmath,amssymb,amsthm,mathtools,microtype}
\usepackage{xcolor,tikz,caption,enumitem,needspace}
\usetikzlibrary{arrows.meta,positioning,calc,fit,backgrounds}
\usepackage[colorlinks=true,linkcolor=blue!40!black,urlcolor=blue!40!black,citecolor=blue!40!black]{hyperref}
\hypersetup{pdftitle={From Packing Stars to Global Coloring},pdfauthor={Rishikesh Gajjala and Arindam Khan}}
\setlength{\parskip}{2pt}
\setlength{\emergencystretch}{2em}
\setlist{itemsep=3pt,topsep=5pt}
\captionsetup{font=small,labelfont=bf,margin=8pt}
\newtheorem{theorem}{Theorem}[section]
\newtheorem{lemma}[theorem]{Lemma}
\newtheorem{proposition}[theorem]{Proposition}
\theoremstyle{definition}\newtheorem{algorithm}{Algorithm}[section]
\numberwithin{equation}{section}
\theoremstyle{remark}\newtheorem*{remark}{Remark}
\newcommand{\E}{\mathbb E}
\newcommand{\Prb}{\mathbb P}
\newcommand{\bp}{\operatorname{bp}}
\newcommand{\one}{\mathbf 1}
\newcommand{\cC}{\mathcal C}
\newcommand{\cP}{\mathcal P}
\newcommand{\eps}{\varepsilon}
\definecolor{inkblue}{RGB}{34,87,128}
\definecolor{warmred}{RGB}{172,70,46}
\definecolor{softblue}{RGB}{228,238,246}
\definecolor{softgold}{RGB}{249,239,213}
\definecolor{softgreen}{RGB}{224,239,228}
\tikzset{vertex/.style={circle,draw=black,fill=white,inner sep=2.5pt},
 lab/.style={font=\small},
 panel/.style={draw=inkblue,rounded corners=3pt,fill=softblue,align=center,inner sep=7pt}}
\title{From Packing Stars to Global Coloring\\\large Asymptotically Optimal Edge Coloring of Weighted Graphs}
\author{Rishikesh Gajjala\\\small New York University Abu Dhabi
 \and Arindam Khan\\\small Indian Institute of Science (IISc)}
\date{}


%%%%%%%%%%%%%%%%%%%%%%%%%%%%%% User specified LaTeX commands.
\newcommand{\dw}{\Delta_w}
\newcommand{\cw}{\chi'_w}

\newif\ifdraft
\drafttrue %comment out this line to remove color

\ifdraft
\usepackage[colorinlistoftodos,prependcaption,textsize=tiny]{todonotes}
\usepackage{xargs}
\newcommand{\ymr}[1]{\todo[color=blue!100!black!50]{AW: #1}}
\newcommand{\akr}[1]{\todo[color=red!100!black!50]{AK: #1}}
\newcommand{\rgr}[1]{\todo[color=green,size=\tiny]{FG: #1}}
\newcommand{\ak}[1]{\textcolor{red}{#1}}
\newcommand{\rg}[1]{\textcolor{green!60!black}{#1}}

\else
    \newcommand{\rgr}[1]{}
    \newcommand{\akr}[1]{}
    \newcommand{\ak}[1]{\textcolor{black}{#1}}
    \newcommand{\rg}[1]{\textcolor{black}{#1}}
\fi


%comment block ends

\makeatletter
\newcommand*\l@result{\@dottedtocline{1}{0em}{9.2em}}
\newcommand*\l@algorithm{\@dottedtocline{1}{0em}{9.2em}}
\newcommand{\listofresults}{%
  \section*{Index of results}%
  \pdfbookmark[0]{Index of results}{result-index}%
  Every entry links to its statement; the right-hand number is its page.
  The two results in Section~\ref{sec:inputs} concern local choices.
  Section~\ref{sec:general-bp} also states the ordinary edge-coloring
  results used there. Algorithm descriptions are indexed below.
  \par\medskip
  \begingroup\footnotesize\setlength{\parskip}{0pt}\@starttoc{res}\endgroup\par\medskip
  \noindent\textbf{Algorithms}\par\smallskip
  \begingroup\footnotesize\setlength{\parskip}{0pt}\@starttoc{alg}\endgroup}
\makeatother
\let\baseTheorem\theorem
\renewcommand{\theorem}[1][]{%
  \baseTheorem[#1]%
  \addcontentsline{res}{result}{\protect\numberline{Theorem~\thetheorem}#1}}
\let\baseLemma\lemma
\renewcommand{\lemma}[1][]{%
  \baseLemma[#1]%
  \addcontentsline{res}{result}{\protect\numberline{Lemma~\thelemma}#1}}
\let\baseProposition\proposition
\renewcommand{\proposition}[1][]{%
  \baseProposition[#1]%
  \addcontentsline{res}{result}{\protect\numberline{Proposition~\theproposition}#1}}

\let\baseAlgorithm\algorithm
\renewcommand{\algorithm}[1][]{%
  \baseAlgorithm[#1]%
  \addcontentsline{alg}{algorithm}{\protect\numberline{Algorithm~\thealgorithm}#1}}
\begin{document}
\maketitle\vspace{-12pt}
\begin{abstract}
We study weighted edge coloring, where we are given an edge-weighted graph $G$ with edge weights in $[0,1]$. We want a coloring of edges using minimum number of colors $\chi'_w(G)$ such that the sum of edge weights of edges of a particular color incident to any vertex of the graph is at most one.  
This is a generalization of two classical problems: edge coloring (all edge weight one) and bin packing (only two vertices). Let $\mathrm{bp}(G)$ be the maximum number of unit bins to pack all edge weights incidents at any vertex and let $\Delta_w(G)$ be the maximum incident weight.  
%If $\chi'_w(G)$ is the edge chromatic number of $G$, then  $\chi'_w(G) \le  \Delta_w(G) \le \mathrm{bp}(G)$. 

For unit-weight bipartite graphs, Kőnig's theorem implies $\Delta_w(G)$ colors are sufficient. However, for arbitrary weighted bipartite graph it is easy to see that $2\Delta_w(G)-1$ colors are sometimes necessary. 
Chung and Ross [1991] conjectured that for a bipartite graph $2\Delta_w(G)-1$ (and a weaker conjecture that $2\mathrm{bp}(G)-1$ ) colors are sufficient. 
After a sequence of work, Feige and Singh [2008] showed $\lceil 9/4 \cdot \Delta_w(G) \rceil$ colors are sufficient, and Khan and Singh [2015] showed $20/9 \cdot \mathrm{bp}(G) +o(\mathrm{bp}(G))$ colors are sufficient. This remained the best-known bounds even for simple bipartite graphs. 
Also, it remains an open problem if there are similar bounds for the weighted case in general (simple) graphs or multigraphs (a generalization of Vizing's theorem or Shannons theorem). 

Our results settle several major open problems in the area.
Our main result produces a surprising analogue of Vizing's theorem for weighted simple graphs:
$\chi'_w(G)\le \mathrm{bp}(G)+o(\mathrm{bp}(G))$.
Informally, this shows that packing of local stars (at each vertex) can be extended to the global coloring. 
For weighted multigraphs, in terms of total weight we prove
$2\Delta_w+o(\Delta_w)$ colors for bipartite multigraphs -- settling the Chung-Ross conjecture asymptotically. 
Furthermore, we show the exact bound $3\lceil\Delta_w\rceil$ for general multigraphs, mirroring the classical Kőnig--Shannon 3/2 gap in bipartite vs general unweighted graphs.

We further show that parallel edges create genuine obstructions: we construct bipartite multigraphs whose underlying simple graph is a degree-three tree with
$\chi'_w(G)=\tfrac43\,\mathrm{bp}(G).$
For bipartite multigraphs we prove
$\chi'_w(G)\le \tfrac{81}{50}\mathrm{bp}(G)+o(\mathrm{bp}(G))$,
improving beyond the harmonic bin-packing constant  $h_\infty\approx 1.691$ and asymptotically going beyong the Chung--Ross bound.  For general multigraphs we obtain
$\chi'_w(G)\le 1.999\,\mathrm{bp}(G)+o(\mathrm{bp}(G))$.
Together with the bipartite lower bound $4/3$, this shows that the
ratio between the optimal general and bipartite asymptotic constants
for $\mathrm{bp}(G)$ is strictly smaller than $3/2$.

Technically, we develop a local-to-global rounding framework that couples bin configurations at the two endpoints of an edge without splitting it, uses hypergraph matching to realize whole-edge choices, and reinterprets Kempe switches as a mixing mechanism for missing colors.  The results expose a common structure behind bin packing, classical edge coloring, matching theory, and multirate network routing. We also provide polynomial-time algorithms for the simple-graph, bipartite, and total-weight bounds. The new $1.999$ bound has a finite construction; its running-time guarantee is distinguished in Section~\ref{sec:algorithm-scope}. 
\akr{wrote a longish abstract. Will shorten later. last technical para needs to be changed.}

\end{abstract}
%\vspace{-6pt}

\newpage
\tableofcontents
\newpage

\begingroup\small
\section*{Main results and direct links}
\pdfbookmark[0]{Main results and direct links}{main-results}
All graphs are finite and have no loops; edge weights lie in $(0,1]$.
Write $\chi'_w(G)$ for the fewest feasible colors, $\Delta_w(G)$ for
the largest total weight at a vertex, and $\bp(G)$ for the largest
number of unit bins needed to pack the weights at one vertex.

\begin{enumerate}[leftmargin=*,itemsep=3pt,topsep=3pt]
\item \textbf{Simple graphs, including graphs that are not bipartite.}
$\chi'_w(G)\le\bp(G)+o(\bp(G))$.
The coefficient one is best possible.
\hyperref[sec:simple]{Proof: Section~\ref*{sec:simple}.}

\item \textbf{Bipartite multigraphs: a bound using local bin packings.}
$\chi'_w(G)\le \frac{81}{50}\bp(G)+o(\bp(G))$.
In particular, $(13/8)\bp(G)$ colors suffice for sufficiently large
$\bp(G)$.
\hyperref[thm:bp-harmonic-first]{First: a simpler $h_\infty$ proof.}
\hyperref[thm:bp-unrestricted]{Then: Theorem~\ref*{thm:bp-unrestricted} and the three-choice improvement.}
\hyperref[alg:bp]{Polynomial-time construction.}

\item \textbf{Bipartite multigraphs: a sharp total-weight bound.}
$\chi'_w(G)\le2\Delta_w(G)+o(\Delta_w(G))$.
The coefficient two is best possible.
\hyperref[thm:total-weight]{Theorem~\ref*{thm:total-weight} and proof.}

\item \textbf{General multigraphs: a bound using local bin packings.}
$\chi'_w(G)\le\frac{1999}{1000}\bp(G)+o(\bp(G))$.
The bound allows arbitrary parallel multiplicity and arbitrarily small
positive weights.
\hyperref[thm:general-bp-subtwo]{Theorem~\ref*{thm:general-bp-subtwo}.}
\hyperref[sec:subtwo-local]{Five-vertex argument.}
\hyperref[sec:general-subtwo]{Full proof and construction.}

\item \textbf{General multigraphs: a sharp total-weight bound.}
$\chi'_w(G)\le3\lceil\Delta_w(G)\rceil$.
There is no lower-order error term.
The coefficient three is best possible.
\hyperref[thm:total-weight]{Theorem~\ref*{thm:total-weight} and proof.}

\item \textbf{An exact lower-bound family.}
For every integer $N\ge1$, a bipartite multigraph whose underlying
simple graph is a tree of maximum degree three has $\bp(G_N)=3N$ and
$\chi'_w(G_N)=4N$. Its ratio is exactly $4/3$.
\hyperref[prop:lower-four-thirds]{Proposition~\ref*{prop:lower-four-thirds} and proof.}

\item \textbf{A counterexample to the exact total-weight conjecture.}
A simple bipartite graph has $\Delta_w(G)=39/20$, $\bp(G)=3$, and
$\chi'_w(G)=4$. This disproves the Feige--Singh bound $2N-1$,
where $N$ is the least integer strictly greater than $\Delta_w(G)$.
\hyperref[prop:counterexample]{Proposition~\ref*{prop:counterexample} and proof.}
\end{enumerate}

Here $o(n)/n\to0$ as $n\to\infty$. All errors are independent
of graph size and minimum edge weight. For unrestricted bipartite
multigraphs, the remaining asymptotic gap is from $4/3$ to $81/50$.
An unrestricted $3/2$ bipartite bound is not established here. For
general multigraphs the gap is from $3/2$ to $1.999$.
In particular, the general-to-bipartite gap between the best asymptotic
$\bp$ coefficients is below $3/2$; see
\hyperref[prop:bp-konig-shannon-gap]{Proposition~\ref*{prop:bp-konig-shannon-gap}}. The construction and running-time
guarantees are stated in Section~\ref{sec:algorithm-scope}.

\endgroup
\clearpage
\listofresults
\clearpage


\section{Introduction}

Edge coloring is a classical and fundamental topic in graph theory and algorithms. We study {\em weighted edge coloring} problem, where we are given an edge-weighted graph $G:=(V,E)$, possibly with parallel edges but without loops,  with edge weights $w \rightarrow [0,1]$. 
A coloring is a \emph{proper weighted coloring} if, at every vertex, the sum
of the weights of incident edges of each color is at most one. The objective
is to minimize the number of colors, denoted by $\chi'(G)$.

This problem is a generalization of classical edge coloring in unweighted graphs (when all edge weights are one) and bin packing (when the graph is a star). Apart from its structural importance, the problem also finds many application in switching networks, scheduling, frequency assignment, and graph decomposition. 
A principal application is multirate routing in three-stage
Clos networks~\cite{Clos53,MT89,CR91}. Here, vertices represent input and
output switches, edges represent unsplittable connection requests, and
colors represent middle-stage switches. Minimizing the number of colors
reduces the interconnection resources needed to support simultaneous
requests in communication and parallel-computing systems. We refer to
Correa and Goemans~\cite{CG07} for a detailed account of this connection.


In (unweighted) edge coloring, a color class forms a matching.  Three classical theorems describe much of the landscape. 
K\H{o}nig's theorem gives $\chi'(G)=\Delta(G)$ for bipartite
multigraphs~\cite{Konig16}; Vizing's theorem gives
$\chi'(G)\le\Delta(G)+1$ for simple graphs~\cite{Vizing64}; and Shannon's
theorem gives $\chi'(G)\le\lfloor3\Delta(G)/2\rfloor$ for arbitrary
multigraphs~\cite{Shannon49}. Here $\Delta(G)$ is the maximum vertex degree.


In the weighted setting, two natural local lower bounds are
$\dw(G):=\max_{v\in V}\sum_{e\ni v}w_e$, and $\bp(G):=\max_{v\in V}\bp_v(G)$,
where $\bp_v(G)$ is the minimum number of unit-capacity bins needed to pack
the weights of the edges incident with $v$. The weighted degree $\dw(G)$
is a volume bound, whereas $\bp(G)$ also captures local incompatibilities
between item sizes. Every proper weighted coloring induces a feasible
bin packing at each vertex, so
$\lceil\dw(G)\rceil\le\bp(G)\le\cw(G)$.
We call the ratio $\cw(G)/\bp(G)$, when $\bp(G)>0$, the
\emph{global-to-local packing gap}: it measures the additional cost of
reconciling locally feasible packings into a single coloring. Even on a
star, the distinction between volume and local packability matters.
A star with $k$ edges of weight $1/2+\eps$ has $\cw(G)=\bp(G)=k$ and
$\dw(G)=k(1/2+\eps)$. In particular, $\cw(G)>2\dw(G)-1$ whenever
$0<\eps<1/(2k)$.



In 1991, Chung and Ross \cite{CR91} conjectured that 
every  bipartite
multigraph admits a proper weighted coloring
 using $2 \bp(G) -1$ colors. 
 \footnote{Their original conjecture was formulated
in the context of Clos networks; this graph-theoretic formulation is
explicit in later work~\cite{FS,KS15}.}Feige and Singh~\cite{FS} had mentioned even a stronger conjecture of proper weighted coloring with $2N-1$ colors, where $N$ is the smallest integer strictly greater
than $\dw(G)$, namely $N=\lfloor\dw(G)\rfloor+1$.
These conjectures ask how efficiently locally feasible
packings can be reconciled into a globally consistent coloring. 
The conjectures motivated a sequence of works in approximation algorithms. 

Du, Gao, Hwang, and Kim~\cite{DGHK98} obtained the bound
$\cw(G) \le \frac{41}{16}\bp(G)+O(1)$. Correa and Goemans~\cite{CG07} improved this to $2.5480\bp(G)+o(\bp(G))$ and separately established
$2.557\dw(G)+o(\dw(G))$. Their approach combines balanced decompositions into matchings with the analysis of an associated continuous bin-packing instance. Feige and Singh~\cite{FS} subsequently gave a combinatorial algorithm using $\lceil\frac94\dw(G)\rceil$ colors. Khan and Singh~\cite{KS15} obtained the explicit guarantee $\lceil2.2223\bp(G)\rceil$. Their algorithm first selects a degree-bounded subgraph favoring heavy edges, colors it using K\H{o}nig's theorem, and then
inserts remaining edges in decreasing weight order. Its analysis uses bin-packing configuration linear programs to show that a blocked edge would contradict local packability. 
%A constructive bound in terms of $\bp(G)$ also gives an approximation guarantee relative to $\cw(G)$, but an extremal gap between these parameters is not approximation hardness.

Ngo and Vu~\cite{NV} constructed bipartite instances exhibiting a
$5/4$ asymptotic gap between $\cw(G)$ and $\bp(G)$, improving the earlier $11/9$ lower bound. They also proved
$\cw(G)\le2\bp(G)-1+\lceil(r-1)/2\rceil$ when each part has $r$ vertices.
Restricted weights admit stronger guarantees. The bound
$2\bp(G)-1$ holds when all weights are at most $1/2$~\cite{MT89}
and when there are only two distinct rates~\cite{CR91}. If every weight
exceeds $1/4$, Khan and Singh~\cite{KS15} obtain
$\lceil2.2\bp(G)\rceil$ colors.


For arbitrary multigraphs, Feige and Singh~\cite{FS} extend their
algorithm to use at most $\lceil\frac94\dw(G)\rceil+\mu(G)$ colors,
where $\mu(G)$ is the maximum edge multiplicity.
Odd cycles and multiplicities create genuine
obstructions. A triangle with $t$ parallel edges on each side, all of weight
$1/2+\varepsilon$, has $\cw(G)=3t$, $\bp(G)=2t$, and
$\dw(G)=t(1+2\varepsilon)$: its global-to-local packing gap is $3/2$, while its ratio to weighted degree approaches three.
Improved bounds are also known
for certain special graph classes; see Huc~\cite{Huc11} and
Sannyasi~\cite{Sannyasi}. 

Despite this progress, prior work had left both conjectures unresolved even for simple bipartite graphs,
and sharp bounds for general multigraphs were also unknown. The following question captures the difficulty:
\begin{quote}
\emph{When locally feasible packings share objects, how many additional
bins are needed to make their assignments globally consistent?}
\end{quote}

\subsection{Our results}

We settle both the conjectures and resolve many fundamental questions in weighted edge coloring. 
\begin{theorem}[Main upper bounds]\label{thm:main}
For all simple weighted graphs,
\[
 \chi'_w(G)\le\bp(G)+o(\bp(G)).
\]
For arbitrary bipartite weighted multigraphs,
\[
 \chi'_w(G)\le2\Delta_w(G)+o(\Delta_w(G)),
\qquad
 \chi'_w(G)\le \frac{81}{50}\bp(G)+o(\bp(G)).
\]
For arbitrary weighted multigraphs without loops,
\[
 \chi'_w(G)\le3\lceil\Delta_w(G)\rceil,
 \qquad
 \chi'_w(G)\le\frac{1999}{1000}\bp(G)+o(\bp(G)).
\]
Each lower-order error that appears here depends only on its displayed
parameter, not on the number of vertices or the smallest edge weight.
The general total-weight bound is exact. The coefficient one for
simple graphs and the coefficients two and three in the total-weight
bounds are best possible.
\end{theorem}
\footnote{
Here,  $h_\infty$ is the classical constant from harmonic bin packing \cite{Epstein}:
\begin{equation}\label{eq:harmonic-constant}
 r_1=1,\qquad r_{j+1}=r_j(r_j+1),\qquad
 h_\infty=\sum_{j=1}^{\infty}\frac1{r_j}
 =1.691 \ldots.
\end{equation}}
%Here $f(n)=o(n)$ means $f(n)/n\to0$ as $n\to\infty$. For example, the
%first conclusion says that for every $\eps>0$ there is a number
%$n_0(\eps)$ such that at most $(1+\eps)n$ colors suffice whenever
%$n=\bp(G)\ge n_0(\eps)$. The same $n_0(\eps)$ works for every number
%of vertices. This point matters: the graph may have far more vertices
%than any fixed function of $n$.

In the following we explain the results in more details. 


\paragraph{Simple graphs: an asymptotic weighted analogue of Vizing's theorem.}
Our first result shows that
$\cw(G)\le\bp(G)+o(\bp(G))$ for every simple graph, whether or not
it is bipartite. This extends the asymptotic coefficient-one phenomenon
of Vizing's theorem to arbitrary edge weights.


The local packing problems here are considerably richer than in the
unweighted setting. When every edge has weight one, a local bin can
contain only one edge. With arbitrary weights, an optimal local packing
may rely on intricate combinations of different item sizes. Moreover,
these packings cannot be chosen independently: an edge belongs to the
packing problems at both of its endpoints and must receive the
\emph{same} bin label in both. Our result surprisingly shows that locally feasible
packings can nevertheless be synchronized with an asymptotically
vanishing relative loss. The result is also algorithmic.

  
\paragraph{Bipartite multigraphs: 
weighted K\H{o}nig-type bounds and 
 resolution of Chung-Ross and Feige-Singh conjectures.}
Parallel edges can destroy this coefficient-one phenomenon even in
bipartite graphs. We give a family of bipartite multigraphs satisfying
$\cw(G)=\frac43\bp(G)$, improving the previous asymptotic lower bound
of $5/4$~\cite{NV} after more than two decades. This reveals a structural
contrast: \emph{nonbipartite structure is asymptotically harmless for
simple graphs, whereas parallel edges alone can create a constant
global-to-local packing gap.}

We disprove the stronger conjecture in Feige and Singh~\cite[Conjecture~1.2]{FS} that predicted $2N-1$ colors, where $N$ is the least integer strictly greater than $\Delta_w(G)$.  In
Appendix~\ref{sec:counterexample}, we give a simple bipartite example with
 $\Delta_w(G)=39/20<2,\qquad \bp(G)=3,\qquad \chi'_w(G)=4$. 
It disproves the total-weight conjecture at $N=2$, while satisfying the
Chung--Ross bound of five colors.
Although this exact weighted-degree conjecture is false, we obtain the
asymptotically optimal bound
$\cw(G)\le2\dw(G)+o(\dw(G))$. 

Sharp weighted-degree bounds do not determine what is possible relative
to $\bp(G)$. The examples establishing the optimality of the coefficients
two and three can have local bin-packing numbers substantially larger
than their weighted degrees. Using $\bp(G)$ as the benchmark therefore
requires accounting for packing structure rather than total weight alone.

For bipartite multigraphs, we first prove
$\cw(G)\le h_\infty\bp(G)+o(\bp(G))$, resolving the Chung-Ross conjecture asymptotically. The appearance of the harmonic
constant is natural. Assign each positive edge weight $w$ the charge
$a(w)=\frac1{\lfloor1/w\rfloor}$,
with $a(0)=0$. A unit-capacity bin has total harmonic charge at most
$h_\infty$. We show how to convert a bound on the total charge at every
vertex into essentially the same bound on the number of global colors.
The improvement to $81/50$ comes from exploiting compatibility between
different size classes before applying the harmonic argument.
  
Together with the lower-bound construction above, the bipartite multigraph problem currently lies in the interval
$\frac43
    \le
    \limsup\frac{\chi'_w(G)}{\mathrm{bp}(G)}
    \le
    \frac{81}{50}$.
Further optimization may improve the bounds slightly, but we do not see how to obtain beyond 1.57 using strategies adapted from harmonic algorithms.

\paragraph{General multigraphs: a weighted analogue of Shannon's theorem.}
For arbitrary multigraphs, classical odd-set obstructions arise in
addition to the local packing obstruction. We prove the nonasymptotic
bound $\cw(G)\le3\lceil\dw(G)\rceil$,
whose leading coefficient three is best possible. The increase from the
bipartite coefficient two to the general coefficient three reproduces
the classical factor-$3/2$ distinction between K\H{o}nig's and Shannon's
theorems.
 

%In the bipartite proof, parallel edges are grouped and Kőnig's theorem colors the groups heavier than $1/2$ with about $2\Delta_w$ colors.  The exact general proof works directly with the original edges.  It scans them from heaviest to lightest and builds a skeleton of maximum degree $2\lceil\Delta_w\rceil$.  Every rejected edge has $2\lceil\Delta_w\rceil$ no-lighter skeleton edges at one endpoint.  Shannon's theorem colors the skeleton with $3\lceil\Delta_w\rceil$ colors, and those witnesses guarantee a common color with enough room at both endpoints for each rejected edge.  Thus the same factor $3/2$ appears in
%$$   \frac32\cdot 2\Delta_w=3\Delta_w. $$
%Thus the classical Kőnig--Shannon distinction survives inside the weighted problem after total incident weight is converted into a degree bound for a carefully chosen skeleton. 

The proof makes this connection explicit. We process edges in decreasing
order of weight and greedily build a subgraph, called the
\emph{skeleton}, of maximum degree at most $2\lceil\dw(G)\rceil$.
Every rejected edge has, at one endpoint, $2\lceil\dw(G)\rceil$
skeleton edges of at least its own weight. Shannon's theorem colors the
skeleton using at most $3\lceil\dw(G)\rceil$ colors, and an insertion
lemma uses these heavier-edge witnesses to extend the coloring to all
remaining edges. Thus weighted load is first converted into a degree
bound with leading coefficient two; the factor $3/2$ from ordinary
multigraph edge coloring then produces the coefficient three.

For the local bin-packing bound, we prove the stronger statement
\[
 \chi'_w(G)\le\frac{1999}{1000}\bp(G)+o(\bp(G)).
\]
The proof packs parallel edges into groups, treats dense sets of three
or five vertices separately, and then colors the remaining graph.
When a set is put back, its internal coloring respects the colors
already used on its boundary. The main local step combines two
feasible patterns that contain the same heavy item. A short list of
exact inequalities, given in Appendix~\ref{sec:subtwo-potential-tables},
controls the resulting five-vertex problem. The full argument is in
Sections~\ref{sec:subtwo-local} and~\ref{sec:general-subtwo}.

The heavy parallel-edge triangle gives the lower coefficient $3/2$.
Thus the general-multigraph interval is now $[3/2,1.999]$.

\begin{proposition}[No Shannon--K\H{o}nig factor-$3/2$ gap for local bin packing]
\label{prop:bp-konig-shannon-gap}
Let $c_{\rm bip}$ and $c_{\rm gen}$ be the smallest asymptotic
coefficients in bounds of the form
$\chi'_w(G)\le c\,\bp(G)+o(\bp(G))$, for bipartite and general
loopless multigraphs, respectively. Then
\[
 1\le\frac{c_{\rm gen}}{c_{\rm bip}}
 \le\frac{1999/1000}{4/3}
 =\frac{5997}{4000}=1.49925<\frac32.
\]
\end{proposition}
\begin{proof}
More precisely, define each coefficient as the infimum of the constants
$c$ for which, for every $\varepsilon>0$, all graphs in the given
class with sufficiently large $\bp(G)$ satisfy
$\chi'_w(G)\le(c+\varepsilon)\bp(G)$.
Bipartite multigraphs are a subclass of general multigraphs, so
$c_{\rm bip}\le c_{\rm gen}$. Proposition~\ref{prop:lower-four-thirds}
gives $c_{\rm bip}\ge4/3$, with arbitrarily large local bin-packing
number. Theorem~\ref{thm:general-bp-subtwo} gives
$c_{\rm gen}\le1999/1000$. Dividing proves the assertion.
\end{proof}

Thus the classical factor-$3/2$ gap does not persist for the optimal
asymptotic constants measured against $\bp(G)$. The exact $4/3$
bipartite lower bound is already enough for this conclusion. The
sharp total-weight coefficients two and three still have ratio $3/2$.

\subsection{Our Techniques}
\akr{need to rewrite}
\paragraph{Gluing local configurations without splitting edges.}
The coefficient-one theorem for simple graphs requires a way to coordinate optimal local packings without independently rounding the two endpoints of an edge.  Such independent rounding would immediately lose a constant factor.

Instead, we regard a local feasible packing as a distribution over knapsack configurations.  We compress these configurations into a bounded collection of \emph{profiles}: a profile records explicit slots for large edges and a reserved amount of capacity for small edges.  For an edge $uv$, the local choices at $u$ and $v$ are then coupled so that their joint mass sums to the original marginal probability, rather than its square.

The remaining difficulty is to choose whole edges.  Large-edge choices become a matching problem in a hypergraph: an option contains the edge itself and one available slot at each endpoint.  A matching selects exactly one compatible option for each edge while using every slot at most once.  Small edges are rounded by a separate capacity-preserving argument.

Conceptually, the result is a rounding theorem of the following form:
\emph{local configuration solutions subject to two overlapping packing constraints can be globally glued with asymptotically vanishing loss.}
This point of view may be useful beyond edge coloring in packing and scheduling problems where each object participates in two local resource constraints.

\paragraph{Kempe switches as a mixing operation.}
Parallel edges create another difficulty.  After heavy packets are properly colored, different vertices may have very different sets of missing colors.  It is not enough for both endpoints of an uncolored edge to have many free colors; they need sufficiently many \emph{common} free colors.

We use classical Kempe switches for this.  For a pair of colors, the corresponding edges form paths and even cycles.  Exchanging the two colors along a path preserves the proper coloring and the number of missing colors at every vertex, but changes \emph{which} colors are missing at its endpoints.

Repeated carefully chosen path switches therefore act as a mixing process on the missing-color sets.  We control not only the frequency of individual missing colors but also pairwise correlations.  After mixing, the intersections of available-color sets behave regularly enough to place almost all remaining weighted edges.

Thus Kempe chains, usually used as a local recoloring or repair device, become here a mechanism for \textit{pseudorandomizing local color availability}.
This mixing viewpoint may be useful in other coloring-completion problems where the main obstruction is not the number of available colors at each object but the overlap between the available sets of neighboring objects.

\paragraph{Why the errors are independent of graph size.}
A recurring technical issue is that the graph can have arbitrarily many vertices.  A union bound over vertices would therefore destroy any guarantee depending only on the local parameter.

Our rounding arguments instead control many local quantities simultaneously with errors that depend on the magnitude of those quantities rather than on how many of them exist.  We combine local-lemma/resampling arguments, concentration with bounded influence, and hypergraph matching results whose hypotheses depend on local degrees and codegrees.  This is what makes the $o(n)$ terms independent of graph order and allows arbitrarily tiny edge weights.

%%%%%%%%%%

%%%%%%%%%%%




%The stated sharpness claims have short examples. A star with $n$ unit-weight edges has $\bp(G)=\chi'_w(G)=n$. A star with $2n-1$ edges, each of weight $n/(2n-1)>1/2$, has $\Delta_w(G)=n$ and needs $2n-1$ colors. These examples do not conflict: the second star also has $\bp(G)=2n-1$. For the general multigraph bound, take a triangle with $t$ parallel edges on each side, all of weight $(t+1)/(2t)$. It has $\Delta_w=t+1$ and needs
%$3t$ colors. Section~\ref{sec:general-multi} gives the full argument
%and proves the upper bound. Its ratio to $\bp=2t$ is only $3/2$;
%Section~\ref{sec:general-bp} proves the upper coefficient





\subsection*{Overview of the paper.}
Section~\ref{sec:warmup} gives a shorter, independent factor-two proof
for simple bipartite graphs. It constructs a random feasible color with
a uniform edge-inclusion probability and controls the number of remaining
edges above every weight threshold, allowing the argument to track local
bin packability rather than just total weight.

%Section~\ref{sec:warmup} gives a shorter independent factor-two proof for simple bipartite graphs. One random feasible color keeps each edge with probability $1/(2m)$. Controlling the number of remaining edges above every weight threshold makes this useful for bin packing.

Section~\ref{sec:round-overview} coordinates all proposed colors at once.
It records slots for large edges and capacity for small edges, couples
these choices across endpoints, and finally selects whole edges. This
stronger rounding argument yields coefficient one for all simple graphs.

Section~\ref{sec:total-weight} proves the weighted-degree bounds. The
bipartite argument groups parallel edges and uses color exchanges along
two-color paths. The nonasymptotic general-multigraph bound uses a different
and shorter argument: a decreasing-weight greedy skeleton, Shannon's
ordinary edge-coloring theorem, and an insertion lemma.

Section~\ref{sec:bp-bounds} develops the bipartite bounds in terms of
$\bp(G)$: first $h_\infty$, then the improvements to $79/48$ and
$81/50$, together with the lower-bound construction.
Section~\ref{sec:general-bp} gives the packet construction and a
simpler algorithmic general-multigraph bound. Sections~\ref{sec:subtwo-local}
and~\ref{sec:general-subtwo} strengthen the asymptotic coefficient to
$1.999$ by treating dense sets separately. The exact local calculations
are collected in the first appendix. Section~\ref{sec:algorithm-scope}
specifies which constructions have a proved polynomial running time.

The main text presents the algorithms and explains their guarantees.
The appendices contain the detailed rounding, mixing, and charge
calculations, including the counterexample to the exact weighted-degree
conjecture in Appendix~\ref{sec:counterexample}.



%Section~\ref{sec:round-overview} then coordinates all proposed colors at once. It records places for large edges and space for small edges, agrees on these choices at both endpoints, and finally chooses whole edges. This stronger step gives coefficient one for all simple graphs. For the bipartite total-weight bound, we group parallel edges and change colors along two-color paths.  The exact general total-weight bound uses a different and shorter argument: a decreasing-weight greedy skeleton, Shannon's ordinary edge coloring theorem, and an insertion lemma. Both arguments are in Section~\ref{sec:total-weight}. Section~\ref{sec:bp-bounds} gives the bipartite bounds in terms of $\bp(G)$: first $h_\infty$, then the improvements to $79/48$ and $81/50$, together with the lower-bound example. Section~\ref{sec:general-bp} uses a two-collection packing inequality and ordinary multigraph edge coloring, followed by a compatible-range precoloring, to obtain $h_\infty+\frac25$. For rational weights, the upper-bound proofs give randomized algorithms with expected polynomial running time. Section~\ref{sec:algorithm-scope} states the input model and algorithms.

\subsection*{Notation}

All graphs in this paper are finite and have no loops. A multigraph may have
several edges between the same two vertices; a simple graph has at most
one. The \emph{load} of a color
at a vertex is the sum of the weights of incident edges with that color.
A coloring is \emph{feasible} if every such load is at most one.
Let  $W_G(v)=\sum_{e\ni v}w_e$.

%We write $\chi'_w(G)$ for the fewest colors in a feasible coloring.

%There are two natural lower bounds. First, define
%\[
%W_G(v)=\sum_{e\ni v}w_e,\qquad
%\Delta_w(G)=\max_v W_G(v).
%\]
%Since one color carries at most one unit of weight at $v$, at least
%$\lceil W_G(v)\rceil$ colors are needed. Second, view the weights incident
%with $v$ as items to be packed into bins of capacity one. Let $\bp_v(G)$
%be the fewest bins needed, and put $\bp(G)=\max_v\bp_v(G)$. Every feasible
%coloring gives such a packing at each vertex, so
%\[
% \Delta_w(G)\le\bp(G)\le\chi'_w(G).
%\]
%Figure~\ref{fig:local-bins} illustrates the second bound. The difficulty
%is that an edge belongs to two local packing problems, and its two ends
%must choose the same color.


\begin{figure}[htbp]
\centering
\begin{tikzpicture}[x=1cm,y=1cm,font=\small]
\node[vertex,label=left:$v$] (v) at (0,0) {};
\foreach \yy/\ww/\cc [count=\idx] in {1.2/0.60/inkblue,0.4/0.40/inkblue,-0.4/0.55/warmred,-1.2/0.45/warmred}{
 \node[vertex] (a\idx) at (2.4,\yy) {};
 \draw[thick,\cc] (v)--node[pos=.72,fill=white,inner sep=1pt,sloped,above] {$\ww$} (a\idx);
}
\node at (1.2,1.85) {Edges at $v$};
\draw[-{Latex},thick] (3.1,0)--(4.1,0);
\node at (6.7,1.85) {Two bins of capacity one};
\fill[softblue] (4.7,.35) rectangle (7.1,1.05);
\fill[softblue!55] (7.1,.35) rectangle (8.7,1.05);
\draw[thick] (4.7,.35) rectangle (8.7,1.05);
\draw (7.1,.35)--(7.1,1.05);
\node at (5.9,.7) {$0.60$};\node at (7.9,.7) {$0.40$};
\node[right] at (8.85,.7) {color 1};
\fill[warmred!15] (4.7,-1.05) rectangle (6.9,-.35);
\fill[warmred!8] (6.9,-1.05) rectangle (8.7,-.35);
\draw[thick] (4.7,-1.05) rectangle (8.7,-.35);
\draw (6.9,-1.05)--(6.9,-.35);
\node at (5.8,-.7) {$0.55$};\node at (7.8,-.7) {$0.45$};
\node[right] at (8.85,-.7) {color 2};
\end{tikzpicture}
\caption{A feasible coloring gives a bin packing at every vertex. Here
each bin contains weight exactly one, so two colors suffice at $v$.
For the whole graph, each edge must also fit the chosen color at its
other endpoint.}
\label{fig:local-bins}
\end{figure}
\akr{this figure is not adding much value for submission? can be removed or put to appendix?}



The set of edges incident with $v$ is $E(v)$. We write $[k]=\{1,\ldots,k\}$.
An indicator $\one_A$ equals one when statement $A$ holds and zero otherwise;
for a set of edges $S$, the vector $\one_S$ records membership in $S$.
Expectations and probabilities are denoted by $\E$ and $\Prb$.
The constants in $O(n^a)$ may depend on parameters fixed before $n$ is
chosen, but never on the number of vertices. A subscript, as in
$O_\eps(n^a)$, names allowed dependencies when useful.

% Standalone warm-up section. Citation keys needed: GG61 and MT.
% Its threshold-domination and bin-rounding lemmas are also used by balanced cuts.
\section{A short factor-two proof for simple bipartite graphs}
\label{sec:warmup}

\begin{theorem}[Warm-up]\label{thm:warmup}
For every $\varepsilon>0$ and all sufficiently large integers $M$,
every finite simple bipartite graph $H$ with $\operatorname{bp}(H)\le M$
satisfies
\[
 \chi'_w(H)\le(2+\varepsilon)M.
\]
The threshold depends only on $\varepsilon$.
For rational edge weights and fixed $\varepsilon$, the proof gives a
randomized algorithm with expected running time polynomial in the input
length and a bound of $(2+\varepsilon)\operatorname{bp}(H)+O_\varepsilon(1)$.
\end{theorem}

There are two points to understand. We can sample a feasible color
without favoring any edge. We can then measure progress by counting the
remaining edges above every weight threshold. Total weight alone would
not give the same control over the number of bins.

\subsection{Sampling one color}
The following interval observation supplies the factor two.

\begin{lemma}[Selecting half the active intervals]
\label{lem:half-intervals}
Let $I_i=(a_i,b_i]\subseteq(0,1]$ be nonempty intervals, and let
$x_i\in[0,1]$ satisfy
\[
 \sum_{i:t\in I_i}x_i\le1\qquad(t\in(0,1]).
\]
Activate the intervals independently, with probabilities $x_i$.
Using additional independent coins, one can select disjoint active
intervals so that $I_i$ is selected with probability $x_i/2$.
\end{lemma}
\begin{proof}
Order the intervals by nondecreasing right endpoint, breaking ties
arbitrarily, and process them in reverse order. Every already processed
interval meeting $I_i$ contains $b_i$. These intervals pairwise intersect,
so at most one can have been selected. Inductively, their selection
probabilities are $x_j/2$. Thus the probability that $I_i$ is blocked is
\[
 q_i=\frac12\sum_{\substack{j>i\\I_j\cap I_i\ne\varnothing}}x_j
       \le\frac{1-x_i}{2}.
\]
Blocking depends only on later activations and their coins, and hence is
independent of the activation of $I_i$. If $I_i$ is active and unblocked,
select it with probability $1/[2(1-q_i)]$. This number is at most one,
and the resulting selection probability is exactly $x_i/2$.
\end{proof}

\begin{lemma}[One feasible color]\label{lem:one-color}
Let $F=(L,R;E)$ be simple and bipartite, and suppose each star has a
specified packing into $m$ bins, where $m\ge1$ is an integer.
There is a random feasible color $X$ such that
\[
 \mathbb P(e\in X)=\frac1{2m}\qquad(e\in E).
\]
The decision at a right vertex uses only the bin choices at its left
neighbors and private coins at that right vertex.
\end{lemma}
\begin{proof}
Pad the specified packings with empty bins where necessary.
Independently at each left vertex, choose one bin uniformly and
tentatively keep its edges. At a right vertex $v$, these tentative
indicators are independent and each has probability $1/m$: simplicity
gives different left endpoints to different incident edges.

Lay out every bin of $v$'s packing as disjoint intervals in $(0,1]$,
using the edge weights as interval lengths. At a point there is at most
one interval from each bin. Lemma~\ref{lem:half-intervals} therefore
applies with $x_e=1/m$. Use private coins at $v$ to select intervals.
The selected edges fit at $v$ because their intervals are disjoint.
They fit at a left vertex because they form a subset of its chosen bin.
Each edge is selected with probability $1/(2m)$.
\end{proof}

This is exactly where simplicity enters the argument. With parallel
edges, several tentative indicators at a right vertex could depend on
the same left-bin choice.

\subsection{From size counts to bin packing}
For a finite item multiset $I$, put
\[
 N_I(t)=|\{i\in I:w_i\ge t\}|.
\]
We use the Gilmore--Gomory pattern LP~\cite{GG61}. Index the items
individually. A pattern is a nonnegative integer vector $p$ satisfying
$\sum_i p_iw_i\le1$; it may repeat an item type. Give nonnegative mass
to patterns, minimize their total mass, and require each item type to
be covered to extent at least one. Write $\operatorname{LP}(I)$ for
this minimum. An integral pattern cover can be turned into a packing
of the physical items by deleting surplus occurrences, so its integral
optimum is exactly $\operatorname{bp}(I)$.

\begin{lemma}[Threshold domination]\label{lem:threshold-domination}
Suppose $I$ fits in $m$ bins. If $0\le\alpha\le1$, $b\ge0$, and
\[
 N_J(t)\le\alpha N_I(t)+b\qquad(t\in(0,1]),
\]
then $\operatorname{LP}(J)\le\alpha m+b$.
\end{lemma}
\begin{proof}
Give each item of $I$ supply $\alpha$ and add a dummy supplier of
capacity $b$. Each item of $J$ demands one unit. A real supplier may
serve an item no larger than itself; the dummy may serve any item.
These supplies satisfy the fractional Hall inequalities. Indeed, if
$S\subseteq J$ is nonempty and $t$ is its smallest weight, then the
real neighbors of $S$ are all items of $I$ of weight at least $t$.
Their supply, including the dummy, is at least $N_J(t)\ge|S|$.
The max-flow theorem gives assignments $\lambda_{ei},\lambda_{e0}$
covering every $e\in J$, with real-supplier loads at most $\alpha$
and dummy load at most $b$.

For each original bin $B$ and $\alpha>0$, independently replace its
item $i$ by item $e$ with probability $\lambda_{ei}/\alpha$, or by
nothing with the remaining probability. Every outcome is a feasible
pattern, since replacements only decrease weights. Give this distribution
total LP mass $\alpha$. Over the $m$ bins it covers $e$ to extent
$\sum_i\lambda_{ei}$. Add singleton pattern $\{e\}$ with mass
$\lambda_{e0}$. The total mass is at most $\alpha m+b$.
For $\alpha=0$, the singleton patterns suffice.
\end{proof}

Patterns may contain repeated types in this replacement experiment.
Later rounding arguments use actual sets of edges, in which repetitions
are forbidden; the two descriptions must be kept distinct.

\begin{lemma}[Elementary bin-packing rounding]\label{lem:bin-rounding}
Fix $0<\tau\le1$ and $\nu>0$. Every finite item multiset $I$ with
weights in $[\tau,1]$ satisfies
\[
 \operatorname{bp}(I)\le(1+\nu)\operatorname{LP}(I)+O_{\nu,\tau}(1).
\]
For fixed $\nu,\tau$ and rational weights, such a packing can be found
in polynomial time.
\end{lemma}
\begin{proof}
The empty instance is immediate. Otherwise let
$L=\operatorname{LP}(I)>0$. The total item weight is at most $L$, so
$|I|\le L/\tau$. Sort the items from largest to smallest and divide
them into consecutive groups of $g=\lceil\nu L\rceil$ items, with a
possibly smaller final group. Set aside the first group for singleton
bins. In each later group, round every weight up to that group's largest
weight.

Assign the rounded items of a group injectively to the original items
of the preceding group. No rounded item is larger than its assigned
original item, and these assignments are disjoint between groups.
In a fractional packing of $I$, replace every occurrence of an assigned
original type by its rounded type, and delete unassigned types. This
shows that the rounded remainder has LP value at most $L$.

There are at most $t=\lceil1/(\nu\tau)\rceil$ rounded types.
An optimal basic solution of their pattern LP has at most $t$ positive
variables, since there are only $t$ demand constraints. Rounding these
variables upward gives an integral pattern cover of cost at most $L+t$.
Delete surplus occurrences and restore the original, no-larger weights.
Including the first group gives
\[
 \operatorname{bp}(I)\le L+t+g\le(1+\nu)L+O_{\nu,\tau}(1).
\]
For the algorithmic statement, a pattern has at most
$\lfloor1/\tau\rfloor$ items. Enumerating all patterns and solving the
resulting LP is polynomial for fixed $\tau$. The grouping and the
second, bounded-type LP are polynomial as well.
\end{proof}

\subsection{One round reduces every star}
Fix $0<\delta<1/8$ and $0<\tau,a<1$. Temporarily remove all edges
lighter than $\tau$, leaving $H^+$. Its maximum degree is at most
$M/\tau$. Put
\[
 q=1-\delta/2+\delta^2<1.
\]

We use the following standard algorithmic form of the local lemma.
If bad events are determined by independent variables, each has
probability at most $p$, and each shares variables with at most $D$
others, then $ep(D+1)\le1$ guarantees an assignment with no bad event.
Moreover, repeatedly resampling the variables of a currently true bad
event terminates in expected polynomial time when the events and their
variables can be listed and tested in polynomial time. This is the
Moser--Tardos resampling theorem~\cite{MT}. Below the number of events
is at most the number of vertices, and each test is explicit.

\begin{lemma}[One round]\label{lem:warm-round}
For all sufficiently large $M$, suppose $F\subseteq H^+$ has specified
star packings into $m$ bins, where $aM\le m\le M$.
Then $\lceil\delta m\rceil$ feasible colors can be removed so that the
residual stars have packings into at most $\lceil qm\rceil$ bins.
The threshold is independent of the number of vertices.
\end{lemma}
\begin{proof}
Sample $k=\lceil\delta m\rceil$ independent colors using
Lemma~\ref{lem:one-color}. Assign an edge appearing in several samples
only to its first sample. Each resulting color is feasible, and the
union of colored edges is unchanged.

At a fixed star $I$, each edge survives with probability
\[
 \rho=\left(1-\frac1{2m}\right)^k\le e^{-\delta/2}.
\]
Its residual $J$ therefore satisfies $\mathbb E N_J(t)=\rho N_I(t)$.
Changing one complete color sample changes $N_J(t)$ by at most $2/t$,
since each feasible color has at most $1/t$ incident edges of weight
at least $t$. The bounded-differences inequality gives
\[
 \mathbb P\left(N_J(t)>\rho N_I(t)+\frac{\delta^2m}{4}\right)
 \le\exp\left(-\frac{(\delta^2m/4)^2t^2}{2k}\right).
\]
Only the distinct incident weights need to be tested. There are at most
$M/\tau$ of them, and each is at least $\tau$. Since $m\ge aM$, the
probability of any failed threshold at one star is
$O_\tau(M)e^{-cM}$ for a constant $c=c(\delta,\tau,a)>0$.

Use as independent variables the left-bin choices and the private
right-vertex seeds, separately for every sample. A failure at a right
vertex depends only on its own seed and its neighbors' bin choices.
A failure at a left vertex depends on right seeds at distance one and
left-bin choices at distance at most two. Events at distance greater
than four have disjoint variables. Since the maximum degree is at most
$M/\tau$, each event shares variables with only $O_\tau(M^4)$ others.
For large $M$, the local-lemma inequality holds. Its resampling algorithm
therefore finds samples for which no threshold fails at any star.

Apply Lemma~\ref{lem:threshold-domination} with
$\alpha=\rho$ and $b=\delta^2m/4$, and then
Lemma~\ref{lem:bin-rounding} with $\nu=\delta^2/8$. This gives
\[
 \operatorname{bp}(J)
 \le(1+\delta^2/8)(e^{-\delta/2}+\delta^2/4)m+O_{\delta,\tau}(1)
 \le(1-\delta/2+\delta^2/2)m+O_{\delta,\tau}(1).
\]
The second inequality follows from
$e^{-\delta/2}\le1-\delta/2+\delta^2/8$ and $\delta<1/8$.
The displayed coefficient is strictly below $q$. Since $m\ge aM$,
the fixed additive term is absorbed for large $M$. The grouping proof
also constructs the required residual packings.
\end{proof}

\subsection{Two greedy rules and the total cost}
\begin{lemma}[Greedy cleanup]\label{lem:general-cleanup}
Every finite loopless weighted multigraph $J$ with maximum weighted
degree at most $R$ satisfies $\chi'_w(J)\le\lfloor4R\rfloor+1$.
\end{lemma}
\begin{proof}
Process the edges in nonincreasing order of weight. For an edge of
weight $b\le1/2$, a color forbidden at an endpoint has load greater
than $1-b\ge1/2$. For $b>1/2$, a forbidden color is nonempty, and
its already colored edges each have weight at least $b>1/2$.
Thus fewer than $2R$ colors are forbidden at each endpoint. A palette
of $\lfloor4R\rfloor+1$ colors always has a color available at both.
\end{proof}

\begin{lemma}[Inserting tiny edges]\label{lem:greedy}
Suppose a weighted graph has maximum weighted degree at most $W$ and
some edges have been feasibly colored using $K$ colors. If every
uncolored edge has weight below $\tau<1$ and $K>2W/(1-\tau)$, all
uncolored edges can be inserted greedily into the same palette.
\end{lemma}
\begin{proof}
For an uncolored edge of weight $b<\tau$, a forbidden color at one
endpoint has load greater than $1-b>1-\tau$. At most $W/(1-\tau)$
colors are forbidden at each endpoint. The palette therefore always
has a color feasible at both endpoints.
\end{proof}

\begin{proof}[Proof of Theorem~\ref{thm:warmup}]
Begin with $m_0=M$. While $m_i\ge aM$, apply
Lemma~\ref{lem:warm-round} and put $m_{i+1}=\lceil qm_i\rceil$.
The number $T$ of rounds is bounded in terms of $\delta,a$, because
\[
 m_i\le q^iM+\frac1{1-q}.
\]
The colors used in these rounds number at most
\[
 \sum_{i<T}\lceil\delta m_i\rceil
 \le\frac{\delta}{1-q}M+O_{\delta,a}(1)
 =\frac2{1-2\delta}M+O_{\delta,a}(1).
\]
At the end each star fits in fewer than $aM$ bins, and hence has total
weight at most $aM$. Lemma~\ref{lem:general-cleanup} colors the terminal
graph with at most $4aM+1$ fresh colors.

The original graph has maximum weighted degree at most $M$. Enlarge the
palette, if necessary, to $\lfloor2M/(1-\tau)\rfloor+1$, and insert
the removed tiny edges using Lemma~\ref{lem:greedy}. The total is
at most
\[
 \max\left\{\left(\frac2{1-2\delta}+4a\right)M+O_{\delta,a}(1),
                  \frac{2M}{1-\tau}+1\right\}.
\]
Choose $\delta,a,\tau$ sufficiently small in terms of $\varepsilon$,
then take $M$ sufficiently large. Both expressions are at most
$(2+\varepsilon)M$.
\end{proof}

The proof can be implemented without solving bin packing exactly.
Use the following steps, with the fixed constants chosen slightly more
accurately than needed for the target $\varepsilon$.
\begin{algorithm}[The factor-two warm-up]\label{alg:warmup}
\begin{enumerate}
\item Remove edges lighter than $\tau$. Use
Lemma~\ref{lem:bin-rounding}, with a fixed small accuracy $\nu_0$, to
pack every remaining star. Let $M_0$ be the largest number of bins used,
and pad the other packings with empty bins. Set $m=M_0$.
\item While $m\ge aM_0$, sample $\lceil\delta m\rceil$ colors.
If some star fails a threshold-count test from
Lemma~\ref{lem:warm-round}, resample precisely the independent variables
on which that failure depends. Continue until no test fails.
\item Give each sampled edge its first sampled color, remove those
edges, and construct new star packings by
Lemma~\ref{lem:bin-rounding}. Pad to $\lceil qm\rceil$ bins, set
$m=\lceil qm\rceil$, and repeat Step~2.
\item Color the terminal graph by decreasing-weight greedy cleanup.
Increase the palette to more than $2\Delta_w(H)/(1-\tau)$ colors if necessary,
and greedily insert all removed tiny edges.
\end{enumerate}
\end{algorithm}
If $M_0$ is below the fixed threshold needed in the proof, skip the
sampling rounds and use greedy cleanup immediately; its cost is then
$O_\varepsilon(1)$. The empty retained graph is handled in the same way.

Lemma~\ref{lem:bin-rounding} ensures
$M_0\le(1+\nu_0)\operatorname{bp}(H)+O_{\nu_0,\tau}(1)$.
Apply the preceding color count with $M_0$ as the size parameter,
and use $\Delta_w(H)\le\operatorname{bp}(H)$ for the final insertion.
Choosing $\nu_0$ small enough gives the stated
$(2+\varepsilon)\operatorname{bp}(H)+O_\varepsilon(1)$ bound.
For fixed accuracy, all pattern enumerations are polynomial, the number
of rounds is bounded, and the resampling theorem gives expected
polynomial time. Every test and every produced color can be checked
explicitly.

The factor two has only one source: the half-selection in
Lemma~\ref{lem:half-intervals}. The other steps preserve its efficiency
up to an arbitrarily small loss.
\section{Rounding without splitting edges}\label{sec:round-overview}
The warm-up keeps half the tentative choices at one endpoint. To reach
coefficient one, we coordinate all colors together. For a vertex $v$, put
\[
 \mathcal K_v=\operatorname{conv}
 \left\{\one_S:S\subseteq E(v),\ \sum_{e\in S}w_e\le1\right\}.
\]
Here $x\in\mathcal K_v$ means that some random feasible subset of the
incident edges contains each edge $e$ with probability $x_e$.
A fractional load at most one alone is not enough: two edges of weight
$0.6$, each given value $0.8$, have fractional load $0.96$, but their
inclusion probabilities cannot both be $0.8$ when at most one edge fits.

\begin{theorem}[Turning local choices into a partial coloring]\label{thm:round}
Fix $C_0,C_1,\eps>0$. Let $G$ be any simple graph with
$W_G(v)\le n$, and let $\cC$ be a set of at most $C_1n$ colors.
Suppose that numbers $x_{e,c}$ satisfy
\[
 0\le x_{e,c}\le C_0/n,\qquad \sum_cx_{e,c}\le1,
 \qquad (x_{e,c}:e\ni v)\in\mathcal K_v
 \quad\text{for every }v,c.
\]
For all sufficiently large $n$, we can feasibly color some of the edges,
using color $c$ on edge $e$ only if $x_{e,c}>0$, so that the colored
weight at every vertex is at least
\[
 \sum_{e\ni v}w_e\sum_cx_{e,c}
       -\eps n-O_{C_0,C_1,\eps}(n^{3/4}).
\]
The threshold for $n$ and the constant in the error term do not depend
on the number of vertices.
\end{theorem}

\subsection*{The idea of the construction}
Fix a small positive threshold $\beta$. A bin contains only a bounded
number of edges of weight at least $\beta$. Record one slot for each
such edge and an amount $R$ for the smaller edges. This record is a
\emph{profile}; its slot sizes $b_i$ satisfy $\sum_i b_i+R\le1$.
Lemma~\ref{lem:profiles} replaces all local bins by a bounded menu of
profiles, while losing only a small fraction of each proposed probability.

\begin{figure}[htbp]
\centering
\begin{tikzpicture}[x=1cm,y=1cm,font=\small]
\fill[softblue] (0,0) rectangle (4.5,.9);
\fill[softgold] (4.5,0) rectangle (7,.9);
\fill[softgreen] (7,0) rectangle (9,.9);
\fill[black!5] (9,0) rectangle (10,.9);
\draw[thick] (0,0) rectangle (10,.9);
\foreach \xx in {4.5,7,9}{\draw (\xx,0)--(\xx,.9);}
\node at (2.25,.45) {$b_1=0.45$};
\node at (5.75,.45) {$b_2=0.25$};
\node at (8,.45) {$R=0.20$};
\node at (9.5,.45) {$0.10$};
\node[align=center] at (2.25,-.55) {one large item};
\node[align=center] at (5.75,-.55) {one large item};
\node[align=center] at (8,-.65) {small\\items};
\node[align=center] at (9.5,-.55) {unused};
\node at (5,1.4) {Total bin capacity: $1$};
\end{tikzpicture}
\caption{A profile records places for large items and space for
small items. Each large-item slot holds at most one item whose weight
fits. Several small items may share the reserved space. Unused capacity
is allowed.}
\label{fig:bin-description}
\end{figure}
An edge must receive the same color at both endpoints. Suppose its
retained probability is $y>0$, and $a_u,a_v$ are its assigned fractions
in two specified endpoint profiles. We pair the conditional choices with
mass $a_ua_v/y$. Summing over the endpoint profiles gives $y$, rather
than $y^2$. This is the reason we do not lose the edge's probability twice.

Profiles are chosen one vertex group at a time, controlling both the
load on every slot and the total error at every vertex. Large edges then
choose disjoint options: each option consists of the edge label and one
slot at each endpoint. Shared extra slots allow an edge to remain
uncolored, while limiting the leftover weight at each vertex. Small edges
are rounded by linear algebra, using the spare space in each profile.
Appendix~\ref{sec:round-proof} gives the full construction.

\begin{algorithm}[Rounding the proposed colors]\label{alg:round}
Input: the graph, palette, fractions, and local feasible-bin distributions
in Theorem~\ref{thm:round}; fix the desired error.
\begin{enumerate}[leftmargin=*,nosep]
\item Partition the vertices into a fixed number of groups. Set aside
the small total weight of edges lying within groups.
\item Build the profile menus and the paired endpoint fractions.
\item Choose profile lists one group at a time, passing the load and
error tests in Appendix~\ref{sec:round-proof}.
\item Clip an edge's total proposal to one if necessary. Create ordinary
and extra slots, then find disjoint options for the large edges.
\item Round one option for each small edge by Lemma~\ref{lem:two}.
Return the resulting feasible partial coloring.
\end{enumerate}
Procedures that run in expected polynomial time for each choice are specified in
Appendix~\ref{sec:algorithms}.
\end{algorithm}

\subsection{From local packings to all simple graphs}\label{sec:simple}
We can now prove the first part of Theorem~\ref{thm:main}. If the graph
has no edges, there is nothing to color. Otherwise put $n=\bp(G)$.
At each vertex $v$, choose a packing into $n$ bins, allowing empty bins,
and denote the edge sets in these bins by $B_{v,1},\ldots,B_{v,n}$.

Use the colors $1,\ldots,n$. For every edge and color set
\[
 x_{e,c}=1/n.
\]
For a fixed $v$ and $c$, the vector of these values is
\[
 (x_{e,c}:e\ni v)=\frac1n\sum_{j=1}^n\one_{B_{v,j}}.
\]
It is an average of feasible sets, so it belongs to $\mathcal K_v$.
Also $\sum_cx_{e,c}=1$ for every edge, and $W_G(v)\le n$. Thus
Theorem~\ref{thm:round} applies with $C_0=C_1=1$.

Fix $\eps>0$. The theorem gives a partial coloring with the $n$ colors
in which the weight left uncolored at any vertex is at most
\[
 R=\eps n+O_\eps(n^{3/4}).
\]
Lemma~\ref{lem:general-cleanup} colors those remaining edges with
at most $4R+1$ additional colors. Therefore
\[
 \chi'_w(G)\le n+4\eps n+O_\eps(n^{3/4}).
\]
After division by $n$, the last error term tends to zero. Since
$\eps$ may be chosen arbitrarily small, this proves
$\chi'_w(G)\le n+o(n)$. Every threshold used depends only on $\eps$,
not on the number of vertices. For the finitely many smaller values
of $n$, Lemma~\ref{lem:general-cleanup} gives the bound $4n+1$, so they can be
included by increasing the error function at those values.


The bound concerns large $\bp(G)$; equality need not hold for small
values. For example, a triangle with unit-weight edges has $\bp(G)=2$
and needs three colors. Such odd cycles are allowed throughout the
proof above. The graph is simple, but it need not have a bipartition.


\begin{algorithm}[Coloring a simple graph]\label{alg:simple}
Compute the approximate local packings in Appendix~\ref{sec:local-input-algorithms}
and let $n$ be the largest
number of bins used. Set $x_{e,c}=1/n$ for $c\in[n]$, using the local
bins as the feasible distributions. Run Algorithm~\ref{alg:round}.
Color its uncolored edges in decreasing weight order with fresh colors,
as in Lemma~\ref{lem:general-cleanup}. Choose both accuracies small
enough for the desired final error. The extra factor from the local
packings is absorbed in that error.
\end{algorithm}

\section{The bounds in terms of total weight}
\label{sec:total-weight}\label{sec:bip-multi}\label{sec:general-multi}
Let $G$ be any weighted multigraph without loops, with $W_G(v)\le n$
at every vertex.  The bipartite proof begins with the grouping below.
The exact general proof later in this section works directly with the
original edges.

First we put parallel edges into groups. Initially each edge forms its
own group. For each fixed pair of endpoints, repeatedly combine two groups
whose weights are both at most $1/2$. The weight of a group is the sum of
its edge weights. Every group produced this way has weight at most one.
When the process ends, at most one group of weight at most $1/2$ remains
for each pair of endpoints. Regard each group as a single edge. Let $H$
consist of the groups of weight greater than $1/2$, and let $R$ consist of
the other groups. Thus $H$ is a multigraph, whereas $R$ is simple and has
weights at most $1/2$. A color assigned to a group will be assigned to
all the original edges in that group. This keeps the load of every color
at every vertex unchanged.

\begin{figure}[htbp]
\centering
\begin{tikzpicture}[x=1cm,y=1cm,font=\small]
\node[vertex,label=left:$u$] (u) at (0,0) {};
\node[vertex,label=right:$v$] (v) at (3.8,0) {};
\foreach \yy/\ww/\cc in {1.6/0.35/inkblue,.8/0.30/inkblue,0/0.20/warmred,-.8/0.15/warmred,-1.6/0.10/warmred}{
 \draw[thick,\cc] (u)..controls (1.1,\yy) and (2.7,\yy)..node[pos=.5,fill=white,inner sep=1pt] {$\ww$}(v);
}
\node at (1.9,1.8) {Five parallel edges};
\draw[-{Latex},thick] (4.55,0)--(5.5,0);
\node[vertex,label=left:$u$] (uu) at (6.2,0) {};
\node[vertex,label=right:$v$] (vv) at (10,0) {};
\draw[very thick,inkblue] (uu)..controls (7.3,1.1) and (8.9,1.1)..node[pos=.5,above=3pt] {$0.65>1/2$}(vv);
\draw[very thick,warmred,dashed] (uu)..controls (7.3,-1.1) and (8.9,-1.1)..node[pos=.5,below=3pt] {$0.45\le1/2$}(vv);
\node at (8.1,1.8) {Two groups};
\end{tikzpicture}
\caption{Parallel edges may be grouped and then treated as single weighted
edges. Here $0.35+0.30=0.65$ and $0.20+0.15+0.10=0.45$.
All original edges in a group receive the same color. The grouping rule
leaves at most one group of weight at most $1/2$ for each pair of vertices.}
\label{fig:packets}
\end{figure}

\begin{lemma}[A common color with room]\label{lem:three-common-room}
Let $n\ge1$ and $0<a<1/2$.  For each of $3n$ colors $c$, let
$x_c,y_c\ge0$ be its present loads at two vertices, and suppose
\[
 \sum_c x_c\le n,\qquad \sum_c y_c\le n.
\]
Suppose also that a set $W$ of at least $2n$ colors satisfies
$x_c\ge a$ for every $c\in W$.  Then some color $c$ has room for an
edge of weight $a$ at both vertices:
\[
 x_c+a\le1\qquad\hbox{and}\qquad y_c+a\le1.
\]
\end{lemma}
\begin{proof}
Assume that no color has room at both vertices.  Then, for each color
$c$, either $x_c>1-a$ or $y_c>1-a$.

Let
\[
 S=\{c:y_c>1-a\},\qquad s=|S|.
\]
If $S$ is nonempty, then
\[
 s(1-a)<\sum_{c\in S}y_c\le n.
\]
The same conclusion $s<n/(1-a)$ is immediate when $S$ is empty.
Every color outside $S$ has $x_c>1-a$.

First suppose that $s\le n$.  There are at least $2n$ colors outside
$S$, so
\[
 \sum_c x_c>(3n-s)(1-a)\ge2n(1-a)>n,
\]
a contradiction.

It remains to consider $s>n$.  At most $n$ of all $3n$ colors lie
outside $W$, since $|W|\ge2n$.  Hence $|S\cap W|\ge s-n$.  The sets
$[3n]\setminus S$ and $S\cap W$ are disjoint, and therefore
\begin{align*}
 \sum_c x_c
 &>(3n-s)(1-a)+(s-n)a.                     \tag{\(*\)}
\end{align*}
The expression on the right has coefficient $2a-1<0$ in $s$.
Using $s<n/(1-a)$, it is therefore larger than
\[
 \left(3n-\frac{n}{1-a}\right)(1-a)
   +\left(\frac{n}{1-a}-n\right)a
 =n+\frac{n(1-2a)^2}{1-a}>n.
\]
This again contradicts $\sum_c x_c\le n$.
\end{proof}


\begin{theorem}[Total-weight bounds]\label{thm:total-weight}
For bipartite weighted multigraphs,
\[
 \chi'_w(G)\le2\Delta_w(G)+o(\Delta_w(G)).
\]
For arbitrary weighted multigraphs without loops,
\[
 \chi'_w(G)\le3\lceil\Delta_w(G)\rceil.
\]
The error in the bipartite statement is independent of the number of
vertices and parallel edges.  The general statement is exact: it has no
lower-order error term.  Both leading coefficients are best possible.
\end{theorem}
\begin{proof}
The empty graph needs no colors. Otherwise let
$n=\lceil\Delta_w(G)\rceil\ge1$.

For the bipartite statement, make the parallel-edge groups above.
Every edge of $H$ has weight greater than $1/2$, so its ordinary maximum
degree is at most $2n$.  The graph $H$ is bipartite. Fix
$0<\varepsilon<1/2$, and let $K$
be the least even integer at least $(2+\varepsilon)n$. K\H{o}nig's
theorem properly colors $H$ with at most $2n\le K$ colors.
For large $n$, $K\le3n$. Lemma~\ref{lem:finish-packets}, with
$\tau=\varepsilon$ and $C=3$, gives
$\chi'_w(G)\le(2+\varepsilon)n+o_\varepsilon(n)$. Taking
$\varepsilon$ arbitrarily small proves the first bound.

We now prove the exact general bound, directly on the original edges.
Order all edges from largest weight to smallest.  Starting with
$F=\varnothing$, inspect the edges in this order.  Put the current edge
into $F$ if its addition leaves the degree of $F$ at most $2n$ at both
endpoints.  Thus
\[
 \Delta(F)\le2n.                                      \tag{1}
\]
If an edge $e\notin F$ was rejected, then at one of its endpoints,
say $u$, there were already $2n$ edges of $F$.  All of them were
inspected before $e$, so each has weight at least $w_e$.  We call them
the witnesses for $e$ at $u$.  The witnesses remain in $F$.

By Shannon's theorem, in the form of
Lemma~\ref{lem:three-per-two}, the multigraph $F$ has a proper edge
coloring with at most $3n$ colors.  Pad it with unused colors to obtain
a palette of exactly $3n$ colors.  This is already a feasible weighted
coloring of $F$, since at most one edge of each color meets any vertex.

We insert the edges outside $F$ one at a time without changing the
colors on $F$.  Consider an edge $e=uv$ of weight $a$.  First,
$a<1/2$.  Indeed, at a witness endpoint the $2n$ witnesses and $e$
are distinct incident edges.  If $a\ge1/2$, their total weight would
be at least $2n a+a>n$, contrary to $W_G(u)\le n$.

Let $x_c$ and $y_c$ be the current loads of color $c$ at $u$ and $v$.
Only a subset of the incident edges has been colored, so
\[
 \sum_c x_c\le W_G(u)\le n,
 \qquad
 \sum_c y_c\le W_G(v)\le n.                         \tag{2}
\]
The $2n$ witness edges at $u$ have distinct colors because the coloring
of $F$ is proper.  In each of those colors the present load at $u$ is
at least the weight of its witness, and hence at least $a$.
Lemma~\ref{lem:three-common-room} now gives a color with room for $e$
at both endpoints.  Give $e$ that color.  The coloring stays feasible,
and the colors of all skeleton edges stay fixed, so the same argument
applies to the next edge.  After finitely many insertions every edge is
colored with one of the same $3n$ colors.  Therefore
\[
 \chi'_w(G)\le3n=3\lceil\Delta_w(G)\rceil.
\]
The examples below prove that the coefficients two and three cannot be
reduced.
\end{proof}

The increase from two to three is the same factor $3/2$ as in ordinary
edge coloring: K\H{o}nig's bound is $\Delta$ for bipartite multigraphs,
whereas Shannon's bound for general multigraphs is
$\lfloor3\Delta/2\rfloor$~\cite{Shannon49}.  In the weighted proof this
same Shannon factor colors the degree-$2n$ skeleton with $3n$ colors;
the common-room lemma then inserts every remaining edge without opening
any new color.


\begin{algorithm}[Coloring a weighted multigraph]\label{alg:total}
In the bipartite case, group parallel edges as above. Properly color the
groups heavier than $1/2$, add unused colors to obtain the even palette
size $K$ from the proof, and find the switches from
Lemma~\ref{lem:weighted-moment-mix}. Compute proposals for the remaining
simple graph, run Algorithm~\ref{alg:round}, color its leftover edges
with fresh colors, and expand every group back to its original edges.

For a general multigraph, use the exact construction instead.  Sort the
original edges by decreasing weight and build the degree-$2n$ skeleton
$F$.  Apply the constructive proof of Lemma~\ref{lem:three-per-two} to
color $F$ with at most $3n$ colors, and pad the palette with unused
colors until it has exactly $3n$ colors.  Then inspect every edge outside $F$ and
scan the palette until finding a color with room at both endpoints.
Lemma~\ref{lem:three-common-room} proves that this scan always succeeds.
Thus the general algorithm uses at most
$3\lceil\Delta_w(G)\rceil$ colors and needs no rounding or cleanup step.
\end{algorithm}

\subsection{Why the coefficients are best possible}
For the bipartite bound, use the star from Section~1: its $2n-1$ edges
all have weight $n/(2n-1)>1/2$, so it has $\Delta_w=n$ and needs
$2n-1$ colors.

For an integer $t\ge2$, take three vertices and put $t$ parallel edges
between each pair. Give every edge weight $(t+1)/(2t)>1/2$.
Any two edges meet, and their weights add up to more than one.
Thus every color contains at most one edge, and
\begin{equation}
 \chi'_w(G_t)=3t,\qquad
 \Delta_w(G_t)=t+1,\qquad \bp(G_t)=2t.
 \label{eq:general-triangle}
\end{equation}
This example is also recorded by Feige and Singh~\cite[Section~1.2]{FS}.
It shows that $\chi'_w(G_t)/\Delta_w(G_t)\to3$, but its ratio to
$\bp(G_t)$ is only $3/2$.

\begin{figure}[htbp]
\centering
\begin{tikzpicture}[x=1cm,y=1cm,font=\small]
\node[vertex,label=below left:$a$] (a) at (0,0) {};
\node[vertex,label=below right:$b$] (b) at (4,0) {};
\node[vertex,label=above:$c$] (c) at (2,2.6) {};
\foreach \bend in {-9,0,9}{
 \draw[inkblue,thick] (a) to[bend left=\bend] (b);
 \draw[inkblue,thick] (b) to[bend left=\bend] (c);
 \draw[inkblue,thick] (c) to[bend left=\bend] (a);
}
\node[below=9pt] at (2,0) {$t$ parallel edges};
\node[rotate=52,above=7pt] at (1,1.3) {$t$ parallel edges};
\node[rotate=-52,above=7pt] at (3,1.3) {$t$ parallel edges};
\node[align=center] at (2,.85) {Each weight\\$\dfrac{t+1}{2t}>\dfrac12$};
\node[align=left,anchor=west] at (5.1,1.4)
 {Every pair of edges meets.\\[5pt]
  Every color holds one edge.\\[8pt]
  $\chi'_w=3t,\quad\Delta_w=t+1,\quad\bp=2t.$};
\end{tikzpicture}
\caption{The obstruction for general multigraphs. Each set of three
drawn curves represents $t$ parallel edges. The required number of
colors approaches three times the maximum total weight, while it is
exactly $3/2$ times the local bin-packing bound.}
\label{fig:general-triangle}
\end{figure}


These examples concern the coefficient of $\Delta_w(G)$. The triangle
has only $\chi'_w(G)=(3/2)\bp(G)$, so it does not rule out a bound
$2\bp(G)+o(\bp(G))$ for general multigraphs. That question remains
unsettled by the results here.

\section{Bounds using the local bin-packing number}\label{sec:bp-bounds}
The coefficient two in the bipartite total-weight bound is sharp, but
its sharpness example has $\chi'_w(G)=\bp(G)$. Parallel edges can
nevertheless force a cost above the local bin-packing number.
\subsection{An exact four-thirds lower bound}
Ngo and Vu~\cite{NV} gave examples with ratio $5/4$. The next
construction gives ratio $4/3$, using only a path with one extra leaf
at each vertex and parallel copies of its edges.

\begin{proposition}[An exact 4/3 lower-bound family]\label{prop:lower-four-thirds}
For every positive integer $N$, there is a bipartite weighted multigraph
$G_N$ such that
\[
 \bp(G_N)=3N,\qquad \chi'_w(G_N)=4N=\frac43\bp(G_N).
\]
The graph obtained by replacing each set of parallel edges by one edge
is a tree of maximum degree three. All weights are rational and greater
than $1/3$.
\end{proposition}

\begin{proof}
\textbf{The graph.}
Put $m=2N$ and take a path $v_0,v_1,\ldots,v_m$.
For $0\le i<m$, give its edge $v_iv_{i+1}$ weight
\[
 w_i=\frac13+\frac{m-i}{6(m+1)}.
\]
Thus $1/2>w_0>w_1>\cdots>w_{m-1}>1/3$.
At every $v_i$, attach a new leaf $u_i$. Give $v_iu_i$ weight
$1-w_i$ when $i<m$, and give $v_mu_m$ weight $1$.
Finally, replace every edge by $2N$ parallel copies of the same weight.
We call the edges incident with the new leaves \emph{leaf edges}.
Figure~\ref{fig:bp-four-thirds} shows the construction.

\begin{figure}[htbp]
\centering
\begin{tikzpicture}[x=1cm,y=1cm,font=\small]
\foreach \name/\xx/\lab in {a/0/0,b/2.8/1,c/7.2/{m-1},d/10/m}{
 \node[vertex,label=below:$v_{\lab}$] (\name) at (\xx,0) {};
 \node[vertex,label=above:$u_{\lab}$] (u\name) at (\xx,1.65) {};
}
\draw[very thick,inkblue] (a)--node[above] {$w_0$}(b);
\draw[very thick,inkblue] (b)--(3.95,0);
\draw[very thick,inkblue,dashed] (3.95,0)--(6.05,0);
\draw[very thick,inkblue] (6.05,0)--(c);
\draw[very thick,inkblue] (c)--node[above] {$w_{m-1}$}(d);
\draw[very thick,warmred] (a)--node[right=3pt] {$1-w_0$}(ua);
\draw[very thick,warmred] (b)--node[right=3pt] {$1-w_1$}(ub);
\draw[very thick,warmred] (c)--node[right=3pt] {$1-w_{m-1}$}(uc);
\draw[very thick,warmred] (d)--node[right=3pt] {$1$}(ud);
\node at (5,1.65) {$\cdots$};
\node[align=center] at (5,-1) {Every segment represents $2N$ parallel edges.\\
Path weights decrease from left to right; $m=2N$.};
\end{tikzpicture}
\caption{The exact $4/3$ family. A leaf edge fits with one path edge
to its right, but cannot fit with a path edge to its left.}
\label{fig:bp-four-thirds}
\end{figure}

\textbf{Local packings use $3N$ bins.}
At an internal vertex $v_i$, pair each of the $2N$ leaf edges with
one of the $2N$ path edges to its right. Each pair has weight one.
Put the $2N$ path edges to its left into $N$ pairs; these fit since
$2w_{i-1}<1$. The same packing at $v_0$ uses only $2N$ bins.
At $v_m$, put its $2N$ unit-weight leaf edges alone and its $2N$
path edges into $N$ pairs. Every leaf vertex needs $2N$ bins.
At $v_m$ these $3N$ bins are necessary: a unit-weight edge must be
alone, and at most two path edges fit in a bin. Hence $\bp(G_N)=3N$.

\textbf{One color contains at most $m+1$ edges.}
Fix any feasible color. In a path bundle, at most two copies have
this color. Assign a single copy to its right endpoint. If two copies
have this color, assign one to each endpoint.
No vertex receives two assignments: that would require one copy from
the path bundle on its left and two copies from the bundle on its right.
Their total weight exceeds one because every path weight exceeds $1/3$.

Every vertex receiving a path-edge assignment has no leaf edge of this
color. At a right endpoint $v_i$ with $i<m$, the path edge from its
left has weight $w_{i-1}>w_i$, so it does not fit with a leaf edge of
weight $1-w_i$. At $v_m$, the leaf edge has weight one. An assignment
to a left endpoint $v_i$ occurs only when two path copies are present;
together with a leaf edge their weight would be
$2w_i+(1-w_i)=1+w_i>1$.

Finally, assign each leaf edge of this color to its path endpoint.
Each leaf bundle contributes at most one edge, since its weight is
greater than $1/2$. All assignments are therefore distinct. We have
assigned every edge of the color to one of $v_0,\ldots,v_m$, proving
the claim. Figure~\ref{fig:bp-four-thirds-count} illustrates the rule.

\begin{figure}[htbp]
\centering
\begin{tikzpicture}[x=1cm,y=1cm,font=\small,>=Stealth]
\begin{scope}
 \node[vertex,label=below:$v_i$] (a) at (0,0) {};
 \node[vertex,label=below:$v_{i+1}$] (b) at (3,0) {};
 \draw[very thick,inkblue] (a)--(b);
 \draw[->,thick,inkblue] (1.1,.28) to[bend left=22] (2.85,.28);
 \node at (1.5,1) {One copy: assign to the right.};
 \draw[warmred,dashed,thick] (b)--(3,.75);
 \node[warmred,font=\large] at (3,.55) {$\times$};
\end{scope}
\begin{scope}[xshift=6cm]
 \node[vertex,label=below:$v_i$] (a) at (0,0) {};
 \node[vertex,label=below:$v_{i+1}$] (b) at (3,0) {};
 \draw[very thick,inkblue] (a) to[bend left=12] (b);
 \draw[very thick,inkblue] (a) to[bend right=12] (b);
 \draw[->,thick,inkblue] (1.15,.5)--(.2,.25);
 \draw[->,thick,inkblue] (1.85,.5)--(2.8,.25);
 \node at (1.5,1) {Two copies: one to each end.};
 \foreach \xx in {0,3}{
  \draw[warmred,dashed,thick] (\xx,0)--(\xx,.75);
  \node[warmred,font=\large] at (\xx,.55) {$\times$};
 }
\end{scope}
\end{tikzpicture}
\caption{Counting the edges of one color. A crossed leaf bundle
contains no edge of this color. The assignments from different path
bundles cannot collide: a collision would require three incident
path edges of the same color.}
\label{fig:bp-four-thirds-count}
\end{figure}

\textbf{Exactly $4N$ colors are needed.}
The graph has $2N(2m+1)$ edges. Since each color has at most $m+1$ edges,
\[
 \chi'_w(G_N)\ge
 \left\lceil\frac{2N(2m+1)}{m+1}\right\rceil
 =\left\lceil4N-\frac{2N}{2N+1}\right\rceil=4N.
\]
To attain this bound, use $2N$ colors for the leaf edges, putting one
copy from every leaf bundle in each color. Use another $2N$ colors for
the path edges, putting one copy from every path bundle in each color.
At an internal vertex, the two path edges in a color have total weight
less than one. Both palettes are feasible, so $4N$ colors suffice.
\end{proof}

The construction is explicit and its coloring takes linear time in the
number of listed edges. Since $\bp(G_N)=3N\to\infty$, it rules out
every upper bound $c\bp(G)+o(\bp(G))$ with $c<4/3$, even on
multigraphs whose underlying simple graph is a tree. Its weights stay
above $1/3$; the gap does not rely on edges of very small weight.

\subsection{A first upper bound: the harmonic constant}

Assign an edge of weight $w$ the charge $a(w)=1/\lfloor1/w\rfloor$.
The main packing fact has a short proof.

\begin{lemma}[Harmonic charge in one bin]\label{lem:general-bin-charge}
Put $a(w)=1/\lfloor1/w\rfloor$. Every finite list of positive weights
with sum at most one has total $a$-charge at most $h_\infty$.
\end{lemma}
\begin{proof}
For every positive integer $t$,
\[
 a(w)\le\frac{t+1}{t}w\qquad(0<w\le1/t).
\]
Indeed, $k=\lfloor1/w\rfloor\ge t$ and
$1/k<((k+1)/k)w\le((t+1)/t)w$.

Start with remaining total weight at most $1/r_1=1$. At step $j$,
suppose the remaining total is at most $1/r_j$. If an item exceeds
$1/(r_j+1)$, it has charge exactly $1/r_j$. Remove it. The new total
is less than $1/r_j-1/(r_j+1)=1/r_{j+1}$, so the next step applies.
Otherwise all remaining items have weight at most $1/(r_j+1)$.
Their total charge is at most
\[
 \frac{r_j+2}{r_j+1}\frac1{r_j}
 =\frac1{r_j}+\frac1{r_{j+1}}.
\]
Including the earlier removed items, this is at most $h_\infty$.
The procedure must stop, since each step that continues removes an
item from a finite list. An empty remainder also satisfies the bound.
\end{proof}

\begin{lemma}[The harmonic tail below $1/6$]
\label{lem:harmonic-tail-six}
If positive weights have total at most $1/6$, then their total harmonic
charge is at most
\[
 h_\infty-\frac32.
\]
\end{lemma}
\begin{proof}
Use the proof of Lemma~\ref{lem:general-bin-charge}, starting with
$r_3=6$ instead of $r_1=1$.  The terms available from that point are
$1/r_3+1/r_4+\cdots=h_\infty-1-1/2$.
\end{proof}

\begin{theorem}[A first bipartite bound: the harmonic constant]
\label{thm:bp-harmonic-first}
Every bipartite weighted multigraph satisfies
\[
 \chi'_w(G)\le h_\infty\bp(G)+o(\bp(G)).
\]
The error is independent of graph order, multiplicity, and minimum
weight. For rational weights and each fixed accuracy, the coloring
can be found in expected polynomial time.
\end{theorem}
\begin{proof}
Put $n=\bp(G)$. Pack each vertex's incident edges into $n$ bins.
Lemma~\ref{lem:general-bin-charge} gives incident charge at most
$h_\infty n$. The harmonic-charge coloring lemma
(Lemma~\ref{lem:harmonic-load}, proved in Appendix~\ref{sec:bp-details})
colors the graph using at most its largest incident charge plus
$o(n)$ colors. Its error is uniform even for arbitrarily small weights,
and its construction is algorithmic. This proves the result.
\end{proof}

\subsection{Improving the bipartite bound}
The harmonic bound can be improved by coloring some compatible large
and medium edges together before applying the charge lemma. A balanced
tenth from two weight ranges gives the short bound
$79/48=1.645833\ldots$.  We give that argument first.  A second step
uses three compatible choices and improves the coefficient to
$81/50=1.62$.  Appendix~\ref{sec:compatible-ranges} gives both
calculations. Appendix~\ref{sec:short-five-thirds} also retains a short
proof of the weaker $5/3$ bound.

\Needspace{16\baselineskip}
\begin{theorem}[A two-range bound for bipartite multigraphs]
\label{thm:bp-two-ranges}
Every bipartite weighted multigraph satisfies
\[
 \chi'_w(G)\le\frac{79}{48}\bp(G)+o(\bp(G)).
\]
More precisely, for every $\eta>0$ there is $N(\eta)$ such that
$\bp(G)\ge N(\eta)$ implies
$\chi'_w(G)\le(79/48+\eta)\bp(G)$.
The threshold is independent of the number of vertices, the number of
parallel edges, and the smallest edge weight. Since
\[
 \frac{79}{48}=\frac53-\frac1{48},
\]
the exact inequality $\chi'_w(G)\le(5/3)\bp(G)$ holds whenever
$\bp(G)$ is sufficiently large.
\end{theorem}
The upper bound first uses the charge
$a(w)=1/\lfloor1/w\rfloor$.
Lemma~\ref{lem:harmonic-load} colors a bipartite multigraph using at
most the largest incident charge plus $o(\bp(G))$ colors.
Grouping parallel edges and mixing their missing colors proves that
lemma, even when weights are arbitrarily small.

To improve the charge bound, select a balanced tenth of the edges in
$\mathcal A=(1/2,3/5]$ and $\mathcal B=(1/3,2/5]$. A vertex-splitting
construction colors the selected graph with at most
$\bp(G)/10+O(1)$ colors. The remaining graph has charge at most
$(371/240)\bp(G)+O(1)$ at every vertex. Their sum is $79/48$.
The exact charge calculation is in Appendix~\ref{sec:bp-details}.

\begin{algorithm}[The two-range bipartite bound]\label{alg:bp-two-ranges}
Use Lemma~\ref{lem:balanced-fraction} with $k=10$ separately on
$\mathcal A$ and $\mathcal B$. Color their selected union by the
two-copy construction in Lemma~\ref{lem:compatible-ranges}.
On the remaining edges run the grouping, mixing, rounding, and cleanup
construction of Lemma~\ref{lem:harmonic-load}, using fresh colors.
Combine the two colorings.
\end{algorithm}

The next theorem is stronger.  Its proof uses the same idea at three
different pairs of weight ranges.  The three selected parts use disjoint
palettes and together cost $(7/30)\bp(G)+O(1)$ colors.  The remaining
graph has harmonic charge at most $(104/75)\bp(G)+O(1)$ at every vertex.

\begin{theorem}[The bipartite three-choice bound]
\label{thm:bp-unrestricted}
Every bipartite weighted multigraph satisfies
\[
 \chi'_w(G)\le\frac{81}{50}\bp(G)+o(\bp(G)).
\]
More precisely, for every $\eta>0$ there is $N(\eta)$ such that
$\bp(G)\ge N(\eta)$ implies
$\chi'_w(G)\le(81/50+\eta)\bp(G)$.  The threshold is independent of
the number of vertices, the parallel multiplicities, and the smallest
edge weight.  Since
\[
 \frac{13}{8}-\frac{81}{50}=\frac1{200},
\]
the exact inequality $\chi'_w(G)\le(13/8)\bp(G)$ holds whenever
$\bp(G)$ is sufficiently large.
\end{theorem}

\begin{algorithm}[The bipartite three-choice algorithm]\label{alg:bp}
Use the three rows of Table~\ref{tab:three-compatible-modes}.  Give each
edge one of the eligible selected labels, or the remainder label, by
Lemma~\ref{lem:two}.  Color each selected part with the vertex-copy
construction of Lemma~\ref{lem:compatible-r-copies}.  On the remainder,
run the grouping, mixing, rounding, and cleanup construction of
Lemma~\ref{lem:harmonic-load}, using fresh colors.  Combine the four
colorings.
\end{algorithm}

The verified asymptotic gap in this paper is therefore
$4/3\le\chi'_w(G)/\bp(G)\le81/50+o(1)$ in the worst case.
An unrestricted $3/2$ upper bound is not established here.

\section{Preparatory bounds for general multigraphs}
\label{sec:general-bp}

We first give a simpler packet construction. It also supplies the
preparation and rounding tools for the stronger $1.999$ result in
Theorem~\ref{thm:general-bp-subtwo}. All these bounds allow arbitrary
parallel edges and arbitrary weights in $(0,1]$.

\begin{theorem}[The packet bound for general multigraphs]
\label{thm:general-bp-baseline}
Every weighted multigraph without loops satisfies
\[
 \chi'_w(G)\le \frac54h_\infty\bp(G)+o(\bp(G)),
 \qquad \frac54h_\infty=2.113787758446\ldots.
\]
More precisely, for each $\varepsilon>0$ there is $C_\varepsilon$ such
that at most $(\frac54h_\infty+\varepsilon)\bp(G)+C_\varepsilon$
colors suffice. The constant is independent of the number of vertices,
parallel multiplicities, and the smallest edge weight. For rational
weights and explicitly listed edges, each fixed accuracy has a
randomized algorithm with expected polynomial running time.
\end{theorem}

We first explain the coefficient. Pack parallel edges into groups of
weight at most one, and properly color the multigraph of groups. Its
maximum degree is at most about $h_\infty\bp(G)$. A separate argument
bounds the number of groups inside every triangle by $2\bp(G)$.
For an odd set of $s\ge5$ vertices, the degree bound loses only the
factor $s/(s-1)\le5/4$. The remaining edges can use colors missing at
their endpoints. We give the details, including this last step.

\subsection{Packing two collections}



For a finite collection $A$ of item types with prescribed counts, write
$F(A)$ for its \emph{fractional count packing cost}. A pattern is a
list of nonnegative integer type counts of total weight at most one.
Assign nonnegative masses to these patterns so that their average
counts equal the prescribed counts; $F(A)$ is the least total mass.
Patterns may repeat a type. Actual packets constructed later will use
distinct available edges. An ordinary packing is one allowed choice,
so $F(A)$ is at most the ordinary bin-packing number.

\begin{lemma}[Packing two collections]\label{lem:general-two-bags}
For any two collections $A,B$, put $a=F(A)$ and $b=F(B)$. Then
\[
 F(A+B)\ge\min\{a/2+b,\ a+b/2\}.
\]
Consequently, for three collections $A,B,C$,
\[
 F(A)+F(B)+F(C)
 \le2\max\{F(A+B),F(A+C),F(B+C)\}.
\]
\end{lemma}
\begin{proof}
Take a fractional packing of $A+B$ of total mass $T$. In each bin,
keep the $A$ items together as one packet and the $B$ items as another.
An $A$-heavy bin has an $A$ packet heavier than $1/2$; define
$B$-heavy similarly. No bin is both. Write $\ell,m,z$ for the masses
of $A$-heavy, $B$-heavy, and remaining bins, interchanging the names
if necessary so that $\ell\ge m$.

Pair all $B$-heavy mass with the same amount of $A$-heavy mass.
In one such pair let the packet weights be $(a_0,b_0)$ and
$(c_0,d_0)$, where $a_0,d_0>1/2$. If $a_0\le d_0$, then
$a_0+c_0\le1$, so the two $A$ packets fit together. If
$a_0\ge d_0$, the two $B$ packets fit together. Thus at least one
collection costs one bin per unit paired mass. For the other
collection, keep its heavy packet in one bin and put two copies of
its small packet in a bin of mass $1/2$. Its cost is at most $3/2$.

Let $p\in[0,m]$ be the paired mass for which the $A$ packets are put
together. An unpaired $A$-heavy bin costs one for $A$ and $1/2$ for
$B$; a remaining bin costs $1/2$ for each, by doubling its packets.
The separate packing costs are therefore at most
\[
 A'=\ell+\frac m2-\frac p2+\frac z2,
 \qquad B'=\frac\ell2+\frac m2+\frac p2+\frac z2.
\]
Empty packets can be omitted. Since
\[
 A'/2+B'=\ell+3m/4+p/4+3z/4\le T,
\]
and the displayed minimum is increasing in both arguments and
symmetric in them, the first assertion follows. For the second,
rename the three costs so that $a\ge b\ge c$. The first assertion
gives $F(A+B)\ge a/2+b\ge(a+b+c)/2$.
\end{proof}

\subsection{Ordinary edge coloring and the remaining edges}

We use two standard results about ordinary multigraph edge coloring.
A fractional proper coloring assigns nonnegative masses to matchings,
covering each edge by total mass at least one; its minimum total mass
is denoted by $\chi'_f(H)$. Edmonds' matching theorem gives
\begin{equation}\label{eq:general-fractional-index}
 \chi'_f(H)=\max\left\{\Delta(H),\
   \max_{\substack{S\subseteq V(H)\\ |S|\ge3\text{ odd}}}
       \frac{2|E_H(S)|}{|S|-1}\right\}.
\end{equation}
Here $E_H(S)$ consists of edges with both endpoints in $S$.
The formula follows from Edmonds' theorem~\cite{Edmonds65}.
We use the following algorithmic ordinary coloring bound.

\begin{proposition}[Ordinary multigraph coloring, Sanders--Steurer]
\label{prop:general-ordinary-coloring}
A loopless multigraph $H$ has a proper coloring with at most
\[
 \chi'_f(H)+\sqrt{9\chi'_f(H)/2}
\]
colors. For explicitly listed edges, such a coloring can be found
deterministically in time
\[
 O\!\left(|E(H)|\sqrt{\Delta(H)}\bigl(|V(H)|+\Delta(H)\bigr)\right).
\]
\end{proposition}
This is Theorem~2.23 of~\cite{SS}. Its error depends only on the
fractional chromatic index, not on the number of vertices. For each
fixed $\eta>0$ the bound is also
$(1+\eta)\chi'_f(H)+O_\eta(1)$.

\begin{lemma}[Finishing with a fixed charge margin]
\label{lem:general-charge-finish}
Fix $b=1/T\le1/4$, $L\in\mathbb N$, and $\rho,C>0$.
Let $H$ have a proper coloring with an even number $K\le Cn$ of
colors. Let $R$ be a weighted multigraph on the same vertices, of
multiplicity at most $L$. Give an edge of $R$ charge
\[
 q_e=\begin{cases}
       a(w_e),&w_e\ge b,\\
       w_e/(1-b),&w_e<b,
     \end{cases}
 \qquad Q_R(v)=\sum_{e\in R(v)}q_e.
\]
If $K-d_H(v)\ge Q_R(v)+\rho n$ at every vertex, then for each
$\zeta>0$, $R$ can be colored using colors missing in $H$, together
with at most $\zeta n+O_{b,L,\rho,C,\zeta}(1)$ new colors.
For rational inputs and fixed constants, the algorithm has expected
polynomial running time.
\end{lemma}
\begin{proof}
Write $f_v=K-d_H(v)$. The charges dominate weights, their sum at a
vertex is at most $Cn$, and their sum on any pair is at most $L$.
After multiplying its size parameter by a fixed constant,
Lemma~\ref{lem:weighted-moment-mix} gives fractions supported on
common missing colors such that $x_{e,c}=O(1/n)$,
$\sum_cx_{e,c}\le1$, and
\[
 \sum_{e\in R(v)}q_ex_{e,c}
 \le Q_R(v)/f_v+o(1)\le1-\rho/C+o(1).
\]
The total uncovered weight at each vertex is $o(n)$.

For completeness, these fractions are averages of actual feasible
sets at each vertex. Put each edge with $w_e\ge b$ in class
$k=\lfloor1/w_e\rfloor$, where $1\le k\le T$. For a fixed vertex
and color, choose class $k$ with probability
\[
 p_k=\max\left\{\frac1k\sum_{e\in I_k(v)}x_{e,c},
                       \max_{e\in I_k(v)}x_{e,c}\right\}.
\]
Conditioned on that choice, the probabilities $x_{e,c}/p_k$ describe
a random set of at most $k$ distinct edges, as in
Lemma~\ref{lem:bp-classchoice}. The same argument works for $k=1$.
Choose the tiny class with probability
\[
 p_0=\max\left\{\sum_{\substack{e\in R(v)\\w_e<b}}
                 \frac{w_e}{1-b}x_{e,c},
                 \max_{\substack{e\in R(v)\\w_e<b}}x_{e,c}\right\}.
\]
Its conditional expected load is at most $1-b$; rounding as in
Lemma~\ref{lem:algorithmic-local-load} gives a distinct-edge set of
load at most one with these inclusion probabilities. Zero-probability
classes are omitted. Since the number of classes is fixed,
\[
 \sum_{k=0}^T p_k
 \le\sum_{e\in R(v)}q_ex_{e,c}+(T+1)\max_e x_{e,c}\le1
\]
for large $n$. Use the empty set for the remaining probability.

Lemma~\ref{lem:bp-mult-round}, again with its size parameter rescaled
by a fixed constant, leaves weighted degree at most
$\delta n+o_\delta(n)$, for any fixed $\delta>0$. Choose
$\delta<\zeta/8$. Lemma~\ref{lem:general-cleanup} then uses at most
$\zeta n$ additional colors for large $n$; bounded $n$ is absorbed
into the additive constant. The rational constructions in
Appendix~\ref{sec:algorithms} implement the mixing, local distributions,
and bounded-multiplicity rounding in expected polynomial time.
All constants are uniform in graph order and minimum weight.
\end{proof}

\subsection{Proof of the general bound}
\begin{proof}[Proof of Theorem~\ref{thm:general-bp-baseline}]
Put $n=\bp(G)$; the empty graph is immediate. Fix a small
$b=1/T\le1/4$. On each parallel bundle, group edges lighter than $b$
into objects of weight in $[b,2b)$, leaving at most one object of weight
below $b$. Edges of weight at least $b$ stay as individual objects.
Expanding an object after coloring preserves feasibility.

The objects at each vertex fit in at most
\[
 n'=\left\lceil\frac{n}{1-2b}\right\rceil+1
\]
bins. To see this, keep the original edges of weight at least $b$ in
their original $n$ bins and insert the new objects greedily. Each new
object weighs less than $2b$. Whenever a new bin is opened, all earlier
bins have load greater than $1-2b$, and the total load is at most $n$.

Next fix $\sigma>0$. In a bundle containing $m$ objects of weight at
least $b$, sort them by decreasing weight and set $q=\lfloor\sigma m\rfloor$.
If $q=0$, leave them unchanged. Otherwise set aside the first $q$
objects, divide the rest into consecutive groups of at most $q$, and
round each weight up to the largest in its group. There are at most
$D=\lceil2/\sigma\rceil+2$ non-tiny types per bundle. The rounded
object at position $j>q$ is no heavier than the original object at
position $j-q$. This is an injection on each bundle, so the retained
graph $P$ still fits in $n'$ bins at each vertex. The set-aside objects
have total weight at most $\sigma n/b$ at a vertex.

Give retained objects the charges in
Lemma~\ref{lem:general-charge-finish}. Lemma~\ref{lem:general-bin-charge}
bounds the original incident charge by $h_\infty n$. Grouping tiny
edges adds at most $4bn$ to it: a new object of weight $w<2b$ has
charge at most $w/(1-2b)\le(1+4b)w$, while the original charge was
at least its total weight. The tiny remainder obeys the same bound.
The rank injection cannot increase total charge. Thus
\begin{equation}\label{eq:general-charge-budget}
 Q_P(v)\le(h_\infty+4b)n.
\end{equation}

Solve the fractional count-packing problem on each bundle of $P$.
Choose an optimum solution supported on at most $D+1$ patterns, one
per possible type constraint. A tiny type may have a very small
weight; its pattern set is finite, and there is at most one actual
tiny object in that bundle. From a pattern of mass $\lambda$, form
$\lfloor\lambda\rfloor$ packets using distinct available objects.
The demand equalities ensure that enough objects of each type exist.
The resulting ordinary multigraph is $H$; write $R$ for the objects
left over. Since a pattern contains at most $1/b$ non-tiny objects,
\[
 \text{multiplicity}(R)\le\lceil(D+1)/b\rceil+1.
\]

Every packet used has total charge at least one. Indeed, the dual
packing problem assigns a price $p_t$ to each type, with total price
at most one in every pattern. Prices can be chosen nonnegative:
replace negative prices by zero, using the constraint for the pattern
with those items deleted. The pattern containing
$\lfloor1/w_t\rfloor$ copies of type $t$ gives $p_t\le a(w_t)$.
For a tiny type, $a(w_t)\le w_t/(1-b)$ as well. Equality of primal
and dual optima makes the total price exactly one on every positive
pattern. Therefore
\begin{equation}\label{eq:general-packet-charge}
 d_H(v)+Q_R(v)\le Q_P(v)\le(h_\infty+4b)n.
\end{equation}

For a triangle, let $A,B,C$ be its three retained bundles. Each pair
fits into $n'$ local bins. Lemma~\ref{lem:general-two-bags} gives
$F(A)+F(B)+F(C)\le2n'$. The number of actual packets on a bundle
is at most its fractional cost, so the triangle has at most $2n'$
packet edges. For every odd vertex set $S$ of size $s\ge5$,
\[
 \frac{2|E_H(S)|}{s-1}
 \le\frac{s}{s-1}\Delta(H)
 \le\frac54(h_\infty+4b)n.
\]
Equation~\eqref{eq:general-fractional-index} now gives
\[
 \chi'_f(H)\le\max\{2n',\tfrac54(h_\infty+4b)n\}.
\]
Since $\frac54h_\infty>2$, Proposition~\ref{prop:general-ordinary-coloring}
and an even-palette
adjustment give a proper coloring with
\[
 K=\left(\frac54h_\infty+O(b)+\eta\right)n+O_{b,\eta}(1)
\]
colors for any fixed $\eta>0$. Add unused colors if needed to ensure
$K\ge\frac54(h_\infty+4b)n$. By
\eqref{eq:general-packet-charge}, the number missing at each vertex
exceeds its residual charge by at least $h_\infty n/4$.
Lemma~\ref{lem:general-charge-finish} colors $R$ with these colors
and $\zeta n+O_{b,\sigma,\zeta}(1)$ extra colors, for any fixed
$\zeta>0$. The set-aside objects cost at most
$4\sigma n/b+1$ further colors by Lemma~\ref{lem:general-cleanup}.

Choose $b$, then $\eta$, then $\sigma$ and $\zeta$ sufficiently small
in terms of the requested $\varepsilon$. The total is
$(\frac54h_\infty+\varepsilon)n+C_\varepsilon$.
All choices depend only on $\varepsilon$. Small $n$ are absorbed
into $C_\varepsilon$ by the same greedy bound. Finally expand all
packets and objects. This proves the theorem.
\end{proof}

\begin{algorithm}[The packet bound for general multigraphs]
\label{alg:general-bp-baseline}
Group tiny edges, set aside the prescribed largest objects on each
bundle, and round retained weights by rank. Use the polynomial-size
bundle programs described below to form integer packets. Properly
color their multigraph by the Sanders--Steurer algorithm and add
unused colors to obtain the required even palette. Apply the
missing-color and rounding constructions of
Lemma~\ref{lem:general-charge-finish}, then greedily color the remaining
objects and expand all packets. Check the output by summing its color
loads.
\end{algorithm}

The preceding proof gives more than its displayed constant.  It is useful
to keep the actual harmonic charge instead of replacing it immediately by
$h_\infty n$.

\begin{lemma}[A charge-sensitive general bound]
\label{lem:general-charge-sensitive}
For a weighted loopless multigraph $J$, put
\[
 A(J)=\max_v\sum_{e\in E_J(v)}a(w_e),
 \qquad a(w)=\frac1{\lfloor1/w\rfloor}.
\]
If $n\ge1$ and $\bp(J)\le n$, then
\[
 \chi'_w(J)\le
 \max\left\{2n,\frac54A(J)\right\}+o(n).
\]
Equivalently, for every $\varepsilon>0$ there is $C_\varepsilon$ such
that the right-hand side may be replaced by
$\max\{2n,5A(J)/4\}+\varepsilon n+C_\varepsilon$.
The error is uniform in the number of vertices, the multiplicities, and
the smallest weight.  For rational weights and fixed accuracy, the proof
is an expected polynomial-time algorithm.
\end{lemma}
\begin{proof}
Repeat the proof of Theorem~\ref{thm:general-bp-baseline}.  The only
change is to retain $A(J)$.  After the tiny-edge grouping and rank
rounding, equations~\eqref{eq:general-charge-budget} and
\eqref{eq:general-packet-charge} become
\[
 d_H(v)+Q_R(v)\le A(J)+4bn.
\]
The triangle bound is still $2n'$, and every odd set with at least five
vertices has density at most $\frac54(A(J)+4bn)$.  Hence
\[
 \chi'_f(H)\le
 \max\left\{2n',\frac54(A(J)+4bn)\right\}.
\]

Write $L=\max\{2n,\frac54A(J)\}$.  The elementary inequality
\begin{equation}\label{eq:charge-margin}
 L-A(J)\ge\frac25n
\end{equation}
follows by splitting at $A(J)=8n/5$: below that value use $L=2n$,
and above it use $L=5A(J)/4$.  Proposition~\ref{prop:general-ordinary-coloring}
therefore colors $H$ with $L+O(bn)+\eta n+O_{b,\eta}(1)$ colors.
Let $K$ be the least even integer at least both the number of colors just
used and $L$, adding unused colors if needed.  Then
$K\le L+O(bn)+\eta n+O_{b,\eta}(1)$.  Since
$d_H(v)+Q_R(v)\le A(J)+4bn$, the number of colors missing at $v$ then
exceeds $Q_R(v)$ by at least
\[
 L-A(J)-4bn\ge(2/5-4b)n.
\]
Choose $b<1/20$.  Thus the margin required by
Lemma~\ref{lem:general-charge-finish} inserts $R$ with $o(n)$ new
colors.  The set-aside objects also cost $o(n)$ colors, exactly as in
the preceding proof.  Letting the fixed accuracy parameters tend to
zero proves the formula.  The same operations give the stated algorithm.
\end{proof}

\begin{theorem}[A simpler algorithmic general multigraph bound]
\label{thm:general-bp-improved}
Every weighted multigraph without loops satisfies
\[
 \chi'_w(G)\le
 \left(h_\infty+\frac25\right)\bp(G)+o(\bp(G)),
 \qquad h_\infty+\frac25=2.091030206757\ldots.
\]
More precisely, for each $\varepsilon>0$ there is $C_\varepsilon$ such
that at most
$\left(h_\infty+\frac25+\varepsilon\right)\bp(G)+C_\varepsilon$
colors suffice.  The constant is independent of graph order, parallel
multiplicity, and the smallest weight.  For rational weights and
explicitly listed edges, each fixed accuracy has a randomized algorithm
with expected polynomial running time.
\end{theorem}
\begin{proof}
Put $n=\bp(G)$.  The empty graph is immediate, so assume $n\ge1$ and
define
\[
 \mathcal A=(1/2,3/5],\qquad
 \mathcal B=(1/3,2/5],\qquad
 t_* =\frac23\left(h_\infty-\frac85\right).
\]
The recurrence defining $h_\infty$ gives $8/5<h_\infty<2$, so
$0<t_*<1$; numerically $t_*=0.0606868045\ldots$.  We first choose a
$t_*$-fraction of
the edges in each range, with an additive error of two at every vertex.
For one range, give each edge the two options ``selected'' and
``unselected.''  Use selected-count constraints at its two endpoints
for the first option and unselected-count constraints for the second.
Lemma~\ref{lem:two} gives a selected subgraph whose degree differs from
$t_*$ times the original degree by at most two.  Apply it separately to
the two ranges, obtaining $F_{\mathcal A}$ and $F_{\mathcal B}$, and put
$F=F_{\mathcal A}\cup F_{\mathcal B}$.

Every one of the $n$ local bins contains at most one $\mathcal A$-edge
and at most two edges from $\mathcal A\cup\mathcal B$.  Therefore
\[
 d_{F_{\mathcal A}}(v)\le t_*n+2,
 \qquad d_F(v)\le2t_*n+4.
\]
Set $K=\lceil t_*n\rceil+2$.  Split each vertex into two copies.  Put
all $\mathcal A$-incidences at the first copy and distribute the
$\mathcal B$-incidences so that both copies have degree at most $K$.
Indeed, the two available capacities have sum
$2K-d_{F_{\mathcal A}}(v)\ge d_F(v)-d_{F_{\mathcal A}}(v)$.
The copied graph is a loopless multigraph of maximum degree at most $K$.
Lemma~\ref{lem:three-per-two} properly colors it with
\[
 3\left\lceil\frac K2\right\rceil
   =\frac32t_*n+O(1)
\]
colors.  After the two copies are joined, a color has at most two edges
at a vertex and at most one of them is in $\mathcal A$.  Two
$\mathcal B$-edges have total weight at most $4/5$, while one edge from
each range has total weight at most one.  Thus this is a feasible
coloring of $F$.

Let $R=G-F$.  A selected $\mathcal A$-edge removes harmonic charge one,
and a selected $\mathcal B$-edge removes charge $1/2$.  The lower degree
bounds from Lemma~\ref{lem:two} give
\[
 A(R)\le
 \max_v\sum_{e\in E_G(v)}q_{t_*}(w_e)+3,
\]
where $q_t$ is defined in Lemma~\ref{lem:discounted-charge-parameter}.
That lemma gives at most $8/5$ charge in each original local bin, so
\[
 A(R)\le\frac85n+3.
\]
Lemma~\ref{lem:general-charge-sensitive} now colors $R$ with
$2n+o(n)$ fresh colors.  Together with the coloring of $F$, this gives
\begin{align*}
 \chi'_w(G)
 &\le\left(2+\frac32t_*\right)n+o(n)\\
 &=\left(h_\infty+\frac25\right)n+o(n).
\end{align*}

For the algorithm, replace $t_*$ by a rational number just above it.
The six charge bounds in Lemma~\ref{lem:discounted-charge-parameter}
can only decrease, and the extra first-stage cost is absorbed into
$\varepsilon$.  Lemma~\ref{lem:two}, vertex splitting, constructive
Shannon coloring, and the charge-sensitive packet construction are all
polynomial for fixed accuracy.
\end{proof}

\begin{algorithm}[The simpler algorithmic general multigraph bound]
\label{alg:general-bp-improved}
Choose a rational $t>t_*$ close enough for the requested accuracy.
Use Lemma~\ref{lem:two} to select the two compatible parts, split every
vertex into two copies, and Shannon-color their union.  On the remaining
graph, run Algorithm~\ref{alg:general-bp-baseline}, retaining its actual
harmonic charge $A$.  For its grouping parameter $b$, choose the packet
palette from
$\max\{2n',5(A+4bn)/4\}$, with the ordinary-coloring slack, and pad it
to at least $\max\{2n,5A/4\}$.  Use disjoint palettes for the two stages
and check the final load of every vertex--color pair.
\end{algorithm}

\subsection{Why the algorithm is polynomial for fixed accuracy}
The number of types in a bundle is fixed by $\varepsilon$, but a tiny
type might fit into a bin an enormous number of times. We can represent
all those patterns using only two choices for each non-tiny pattern.

List all non-tiny count vectors $a$ of load $\ell(a)\le1$. There
are at most $(\lfloor1/b\rfloor+1)^D$, a constant for fixed accuracy.
If a tiny type of weight $t$ remains, put
\[
 k_a=\left\lfloor\frac{1-\ell(a)}{t}\right\rfloor.
\]
Keep only the two patterns $(a,0)$ and $(a,k_a)$, removing a duplicate
when $k_a=0$. When $k_a>0$, every omitted pattern $(a,k)$ lies between them:
\[
 (a,k)=\left(1-\frac{k}{k_a}\right)(a,0)
             +\frac{k}{k_a}(a,k_a)\qquad(0\le k\le k_a).
\]
Both endpoint patterns cost one, so this replacement preserves all
type demands and total packing cost. The smaller LP therefore has
exactly the same optimum as the full count-packing LP. If there is no
tiny type, only the non-tiny patterns are needed.

Although $k_a$ can be large, its binary encoding has polynomial
length. The LP has a constant number of rows and columns for fixed
accuracy and rational coefficients of polynomial encoding length.
It can be solved exactly in polynomial time, with a basic optimum
using at most $D+1$ patterns. Flooring their masses uses only available
objects: in particular, a pattern with at least one full unit of mass
cannot require more than the one available tiny object. Thus the
packet graph has at most as many edges as the original graph.

No exact computation of $\bp(G)$ is required. Use the local packing
algorithm in Appendix~\ref{sec:local-input-algorithms} to obtain a
computable common bound
\[
 n\le m\le(1+\xi)n+O_\xi(1),\qquad n=\bp(G).
\]
Run the construction with $m$ in place of $n$ and choose fixed
$\xi>0$ small enough to absorb this change in the final accuracy.
All its load, charge, and local-bin bounds remain valid.  The recurrence
for $r_j$ gives rational upper and lower approximations to $h_\infty$;
the tail after $r_j$ is at most $2/r_{j+1}$ because the terms at least
double.  Use these approximations to choose the rational selection
fraction $t>t_*$.  In the residual packet stage, compute its actual
rational charge bound $A$ and pad the palette to a rational upper bound
for
\[
 \max\left\{2\left(\left\lceil\frac{m}{1-2b}\right\rceil+1\right),
                  \frac54(A+4bm)\right\},
\]
plus the fixed accuracy slack required by the ordinary coloring and
completion steps.

If $m$ is below the fixed finishing threshold, use decreasing-weight
greedy coloring directly; it uses at most $4m+1=O_\varepsilon(1)$
colors. Otherwise use the construction above.

The ordinary coloring algorithm is deterministic and polynomial.
The mixing and rounding algorithms already proved in
Appendix~\ref{sec:algorithms} have expected polynomial running time.
Sorting, grouping, packet formation, and
greedy cleanup are polynomial. Consequently the entire algorithm has
expected polynomial running time in the number of listed edges and
their rational encoding length, for each fixed $\varepsilon$.
It always returns a feasible coloring. No bound polynomial jointly
in the input length and $1/\varepsilon$ is claimed.

The next two sections improve the general asymptotic coefficient
from $h_\infty+2/5$ to $1.999$. We keep the simpler construction above
because its fixed-accuracy expected polynomial running time has been
established separately.

\clearpage
\begin{theorem}[The unrestricted $1.999$ bound]
\label{thm:general-bp-subtwo}
Every weighted multigraph without loops satisfies
\[
 \chi'_w(G)\le\frac{1999}{1000}\bp(G)+o(\bp(G)).
\]
More precisely, for every $\varepsilon>0$ there is a constant
$C_\varepsilon$ such that
\[
 \chi'_w(G)\le(1.999+\varepsilon)\bp(G)+C_\varepsilon.
\]
The constant is independent of graph order, parallel multiplicity,
and the smallest positive weight.
\end{theorem}

\subsection{Preparing and packing the parallel bundles}
\begin{lemma}[Grouping, rounding, and forming packets]
\label{lem:general-packet-preprocess}
Fix $b=1/T\le1/4$ and $\sigma>0$. A weighted multigraph with
$\bp(G)=n$ can be replaced by a retained object graph $P$ and
set-aside objects as follows. Put
\[
 N=\left\lceil\frac{n}{1-2b}\right\rceil+1,
 \qquad D=\lceil2/\sigma\rceil+2.
\]
Each retained star fits in $N$ bins. Each original parallel bundle has
at most $D$ non-tiny types and one tiny object. The set-aside objects
have total weight at most $\sigma n/b$ at each vertex. The retained
objects can be packed into whole packets on each original pair; their
number is at most that pair's fractional count-packing cost. The
remaining objects have multiplicity at most $\lceil(D+1)/b\rceil+1$.
The charge bounds \eqref{eq:general-charge-budget} and
\eqref{eq:general-packet-charge} hold, with the charges from
Lemma~\ref{lem:general-charge-finish}. Expanding any packet or object
back into its original edges preserves feasibility.
\end{lemma}
\begin{proof}
Use the grouping, rank rounding, and basic fractional bundle packing
in the proof of Theorem~\ref{thm:general-bp-baseline}.
The rank injection gives the stated local bin bound; flooring at most
$D+1$ positive pattern masses leaves at most
$\lceil(D+1)/b\rceil+1$ objects on each original pair.
Equations~\eqref{eq:general-charge-budget} and
\eqref{eq:general-packet-charge} give the charge bounds.
Every group stays on one original pair, so expanding its objects
preserves all color loads.
\end{proof}

\section{The five-vertex argument}\label{sec:subtwo-local}
We work with the fractional count model defined in
Section~\ref{sec:general-bp}: a type may appear repeatedly in a pattern,
while the prescribed demands may be real. An item is \emph{heavy} if
its weight is greater than $1/2$, and \emph{light} otherwise.
The key observation is that two local patterns containing the same
heavy item can share one copy of that item in a feasible whole pattern.
We then assign light profits to the two endpoints so that these local
bounds add up to a bound on the entire demand.


\begin{lemma}[Light-item harmonic charges]
\label{lem:subtwo-light-charge}
For an item of weight $w\in(0,1]$, put
$\tau(w)=1/\lfloor1/w\rfloor$. A collection of items of total weight
at most one, each of weight at most $1/2$, has total charge at most
\[
 A=\frac{205}{144}.
\]
A collection of light items of total weight strictly below $1/2$ has
total charge strictly below
\[
 R=\frac{25}{36}.
\]
If $m_e$ is any upper bound on the price of a feasible packet on
bundle $e$, then an item $t$ on $e$ of weight $w_t$ and price $p_t$
satisfies
$p_t\le m_e\tau(w_t)$.
\end{lemma}
\begin{proof}
For $w\le1/j$, the elementary estimate
$\tau(w)\le(1+1/j)w$ follows by considering the integer
$\lfloor1/w\rfloor$. We use it for $j=3,4,6,12$.

In a light-only unit bin, there are at most two items heavier than
$1/3$. If there are none, the charge is at most $4/3$. If there is
one, it contributes $1/2$, and the remaining load is below $2/3$;
the charge is below $1/2+(4/3)(2/3)=25/18$. If there are two, they
contribute one and leave load below $1/3$. When all remaining items
weigh at most $1/4$, their charge is below $5/12$, giving $17/12$.
Otherwise exactly one of them weighs more than $1/4$; it contributes
$1/3$ and leaves load below $1/12$. The final charge is below
$(13/12)(1/12)=13/144$, giving $A$.

For total light load below $1/2$, if every item weighs at most $1/3$,
the charge is below $2/3$. Otherwise one item weighs more than
$1/3$, contributes $1/2$, and leaves load below $1/6$; the remaining
charge is below $(7/6)(1/6)=7/36$. This gives $R$. Finally,
$\lfloor1/w_t\rfloor$ copies of $t$ fit on its bundle, so their total
price is at most $m_e$.
\end{proof}

\begin{lemma}[Splitting profit between endpoints]
\label{lem:subtwo-endpoint-potential}
Let $V$ have five vertices. Each bundle contains finitely many item
types of positive weight and nonnegative price. A pattern may repeat
a type. Assume that the maximum price of a globally feasible pattern
is one and that the prescribed item-count vector lies in the
set of mean counts of a random feasible star pattern at every vertex. Write
$C=\sum_t p_tx_t$, and choose any nonnegative upper bound $m_e$ on
the price of a feasible packet using only bundle $e$.

Fix numbers $0\le\rho_e\le1$ for the bundles and $h_v\ge0$ for the
vertices. Suppose there are nonnegative numbers $z_{ve},c_v,P_v$
and $\alpha_{ve}$, for incident $v,e$, satisfying, for every
$e=uv$,
\begin{align}
 z_{ue}+z_{ve}&=2m_e,                                      \label{eq:subtwo-pot-sum}\\
 z_{vf}&\le c_v+h_vm_f\quad(f\ni v),                     \label{eq:subtwo-pot-light}\\
 P_v&\ge Ac_v+h_v,                                        \label{eq:subtwo-pot-base}\\
 \alpha_{ve}&\ge z_{vf}-\rho_em_f\quad(f\ni v,\ f\ne e), \label{eq:subtwo-pot-off}\\
 \alpha_{ve}&\ge z_{ve}-(2-\rho_e)m_e,                     \label{eq:subtwo-pot-one}\\
 \alpha_{ue}+\alpha_{ve}&\ge\rho_em_e,                    \label{eq:subtwo-pot-two}\\
 P_u+P_v&\ge(2-\rho_e)m_e+\rho_e
                     +R(\alpha_{ue}+\alpha_{ve}).          \label{eq:subtwo-pot-edge}
\end{align}
Then $2C\le\sum_{v\in V}P_v$.
\end{lemma}
\begin{proof}
Pad each local fractional packing with empty bins so that it is a
probability law on feasible integer star patterns. At the two ends
of a positive-price bundle $e$, give light-item profit respective
weights $z_{ue}/m_e$ and $z_{ve}/m_e$; give heavy-item profit weight
one at each end. Equation~\eqref{eq:subtwo-pot-sum} and the common
item-count means show that the sum of these five weighted star means
is $2C$. A zero-price bundle contributes zero and needs no division.

In a star state containing no item heavier than $1/2$, write $S_f$
for its ordinary profit on bundle $f$. By
\eqref{eq:subtwo-pot-light}, its weighted profit is at most
$c_v\sum_fS_f/m_f+h_v\sum_fS_f$. Lemma~\ref{lem:subtwo-light-charge}
bounds the first sum by $A$; terms with $m_f=0$ have $S_f=0$ and are
omitted. The whole star is itself a feasible
global pattern, so the second sum is at most one. Thus its weighted
profit is at most $Ac_v+h_v\le P_v$.

Now fix a heavy type $t$ on $e=uv$ and choose any feasible star at
each endpoint containing $t$. Keep one copy of $t$ and all items
from both stars on bundles other than $e$. This is a feasible global
pattern: the endpoint loads do not increase, and at any third vertex
the two incoming off-$e$ loads are each below $1/2$. If $p$ is the
price of $t$, and $O_u,O_v$ are the off-$e$ profits, then
\begin{equation}\label{eq:subtwo-pot-pair}
                         p+O_u+O_v\le1.
\end{equation}
At either endpoint $i\in\{u,v\}$, the heavy item and the remaining
items on $e$ form a packet of profit at most $m_e$. The light
remainder of its star has harmonic charge below $R$.

Write $(x)_+=\max\{0,x\}$. For the same-bundle light remainders, the two linear
conditions~\eqref{eq:subtwo-pot-one} and
\eqref{eq:subtwo-pot-two}, together with
$z_{ue}+z_{ve}=2m_e$, imply
\begin{equation}\label{eq:subtwo-pot-positive-part}
 (z_{ue}-\alpha_{ue})_+
 +(z_{ve}-\alpha_{ve})_+
                  \le(2-\rho_e)m_e.
\end{equation}
Indeed, if both terms are positive, add them and use
\eqref{eq:subtwo-pot-two}; if at most one is positive, use
\eqref{eq:subtwo-pot-one}. Let $E_i$ be the light profit remaining
on $e$ in endpoint $i$'s star, let $O_{if}$ be its profit on
$f\ne e$, and put $O_i=\sum_{f\ne e}O_{if}$. The endpoint packet
bound gives $E_i\le m_e-p$. The residual normalized profit
\[
 r_i={E_i\over m_e}+\sum_{f\ne e}{O_{if}\over m_f}
\]
is at most $R$ by Lemma~\ref{lem:subtwo-light-charge}, omitting
zero-price summands. For $m_e>0$, the paired weighted profit $Q$
therefore satisfies
\begin{align*}
 Q
 &\le 2p+\rho_e(O_u+O_v)+\sum_i\alpha_{ie}r_i
       +\sum_i{(z_{ie}-\alpha_{ie})_+\over m_e}E_i\\
 &\le 2p+\rho_e(1-p)+R(\alpha_{ue}+\alpha_{ve})
       +(2-\rho_e)(m_e-p)\\
 &=\rho_e+(2-\rho_e)m_e+R(\alpha_{ue}+\alpha_{ve}).
\end{align*}
The first line uses \eqref{eq:subtwo-pot-off} on off-$e$ light
items and separates the positive part of each same-bundle
coefficient. The second uses the feasible merge
$p+O_u+O_v\le1$ and~\eqref{eq:subtwo-pot-positive-part}.
If $m_e=0$, then $p=E_i=0$ and $z_{ue}=z_{ve}=0$; omit the
same-bundle fractions and the same bound follows directly.

For each heavy type, the two endpoint laws put the same mass $x_t$
on states containing that type; they need not have the same
conditional laws. Couple them arbitrarily and apply the preceding
bound. A star contains at most one heavy item. If $H_e$ is the total
mean count of heavy items on $e$, summing the no-heavy and paired
states gives
\[
 2C\le\sum_vP_v\Bigl(1-\sum_{e\ni v}H_e\Bigr)
       +\sum_{e=uv}H_e(P_u+P_v)=\sum_vP_v.
\]
\end{proof}

\begin{figure}[htbp]
\centering
\begin{tikzpicture}[x=1cm,y=1cm,font=\small]
 \node[vertex,label=left:$u$] (u) at (-2,1) {};
 \node[vertex,label=right:$v$] (v) at (2,1) {};
 \node[vertex,label=below:$a$] (a) at (-2,-1) {};
 \node[vertex,label=below:$b$] (b) at (0,-1) {};
 \node[vertex,label=below:$c$] (c) at (2,-1) {};
 \foreach \x in {a,b,c}{
   \draw[inkblue,thick] (u)--(\x);
   \draw[green!45!black,thick] (v)--(\x);
 }
 \draw[warmred,line width=2pt] (u)--node[above=5pt,black]
   {keep one copy of the shared item, $w>1/2$} (v);
 \node[inkblue,align=center] at (-3.5,.0)
   {light items\\from the star at $u$};
 \node[green!45!black,align=center] at (3.5,.0)
   {light items\\from the star at $v$};
 \node[align=center] at (0,-2.0)
   {At each lower vertex, either star contributes less than $1/2$.};
\end{tikzpicture}
\caption{Pair two local patterns containing the same heavy item.
Keep that item once and both collections of light items on other
bundles. The resulting whole pattern is feasible. Light items on
the shared bundle are omitted here and charged separately in the proof.}
\label{fig:subtwo-shared-heavy}
\end{figure}

\begin{theorem}[A five-vertex local bound]
\label{thm:subtwo-local-five}
Consider finitely many repeatable item types on at most five vertices.
Suppose their demand has a fractional unit-bin packing of mass at most
one at every vertex. Give the types nonnegative prices so that every
globally feasible integer-count pattern has price at most one. Then
its total demand price satisfies $C\le157/80$.
Equivalently, the demand has a fractional packing into globally
feasible integer-count patterns of total mass at most $157/80$.
\end{theorem}
\begin{proof}
Add isolated vertices if necessary. If every item price is zero there
is nothing to prove. Otherwise normalize the maximum whole-pattern
price to one; multiplying back at the end only weakens the bound. Let $\beta_e$ be the actual maximum profit of a feasible
packet on bundle $e$. A packet on $e$ and one on a disjoint bundle
$f$ fit in one whole pattern, so
\[
 0\le\beta_e\le1,\qquad
 \beta_e+\beta_f\le1\quad(e\cap f=\varnothing).
\]
Consequently all edges with $\beta_e>1/2$ meet pairwise; they all share one vertex or lie in one triangle.

Fix their largest price $a\in(1/2,1]$ and put $s=1-a$.
At this fixed $a$ consider two simple upper-price vectors:
$S_0$, the star with all four edges at $0$ priced $a$; and $T$,
the triangle with $01,02,12$ priced $a$. In both vectors every other edge has upper price $s$. Let $S_1$ be $S_0$ with
vertices $0,1$ swapped. These three \emph{labelled} vectors are
\begin{equation}
\begin{array}{c|cccccc}
 &01&02&03,04&12&13,14&23,24,34\\ \hline
 S_0&a&a&a&s&s&s\\
 S_1&a&s&s&a&a&s\\
 T  &a&a&s&a&s&s
\end{array}                                                       \label{eq:subtwo-model-prices}
\end{equation}
Appendix~\ref{sec:subtwo-potential-tables} gives endpoint bounds of total at most
$157/40$ for $S_0$ and $T$, for every $a$. The bounds for $S_1$
come from those for $S_0$ by swapping $0,1$. This relabelling
preserves their common parameters
\begin{equation}
\rho_{34}={2s\over3a},\quad \rho_e={2\over3}\ (e\ne34),
 \qquad
 h_0=h_1=h_2={2a-1\over3a},\quad h_3=h_4=0.             \label{eq:subtwo-common-parameters}
\end{equation}

The price vectors between these three are easy to describe. For
$s\le y\le x\le a$, let $M(x,y)$ have prices
\begin{equation}
\begin{array}{c|cccccc}
 &01&02&03,04&12&13,14&23,24,34\\ \hline
 M(x,y)&a&x&y&1-y&1-x&s
\end{array}                                                       \label{eq:subtwo-majorant}
\end{equation}
A coordinate check in~\eqref{eq:subtwo-model-prices} gives
\begin{equation}
\begin{gathered}
M(x,y)=\lambda_0S_0+\lambda_1S_1+\lambda_TT,\\
 \lambda_0={y-s\over2a-1},\quad
 \lambda_1={a-x\over2a-1},\quad
 \lambda_T={x-y\over2a-1}.
\end{gathered}
\label{eq:subtwo-model-mixture}
\end{equation}
The three coefficients are nonnegative and sum to one;
see Figure~\ref{fig:subtwo-two-models}.

Suppose first that the high-price edges lie in a star centred at
$0$. Order its four spoke prices as
$a\ge b\ge c\ge d$, padding any value below $1/2$ to $1/2$.
The actual prices are dominated by $M(b,c)$. Indeed its spoke
upper prices are $a,b,c,c$, with the fourth at least $d$.
On a core edge $ij$, its price is at most $1/2$ and at most
$1-\beta_{0k}$ for either disjoint spoke $0k$. These facts give,
in the order of~\eqref{eq:subtwo-majorant}, the core bounds
$1-c$ on $12$, $1-b$ on $13,14$, and $s$ on $23,24,34$.
Thus~\eqref{eq:subtwo-model-mixture} supplies a convex combination of the three upper-price vectors dominating all actual bundle prices.

If the high-price edges instead lie in a triangle, label their
prices $a\ge b\ge c>1/2$ on $01,02,12$. A nontriangle edge is
disjoint from one of these three edges, giving the upper prices
$1-c$ on $03,04$, $1-b$ on $13,14$, and $s$ on $23,24,34$.
This is exactly $M(b,1-c)$; the condition
$s\le1-c\le1/2<b\le a$ makes~\eqref{eq:subtwo-model-mixture} applicable.

The parameters in~\eqref{eq:subtwo-common-parameters} depend only on $a$, so every inequality of
Lemma~\ref{lem:subtwo-endpoint-potential} is linear in the upper
prices and the chosen bounds. Average all the bounds for
$S_0,S_1,T$ with the coefficients in~\eqref{eq:subtwo-model-mixture}. This averages numerical bounds only; the original local pattern laws
are unchanged. The resulting upper prices dominate the actual
$\beta_e$, and the total endpoint bound
is at most $157/40$. If no actual price exceeds $1/2$, use
$S_0$ at $a=1/2$; all its upper prices are $1/2$.
The endpoint lemma gives $C\le157/80$.

The equivalent global fractional packing statement follows from
finite linear-programming duality. Nonnegative prices suffice
because deleting items preserves feasibility.
\end{proof}

\begin{figure}[htbp]
\centering
\begin{tikzpicture}[x=1cm,y=1cm,font=\small,>=Stealth]
\foreach \j/\captiontext in {0/{$S_0$},1/{$S_1$},2/{$T$}}{
 \begin{scope}[xshift=2.65cm*\j]
  \foreach \v/\ang in {0/90,1/18,2/-54,3/-126,4/162}{
   \coordinate (p\j\v) at (\ang:.72);
  }
  \foreach \u/\v in {0/1,0/2,0/3,0/4,1/2,1/3,1/4,2/3,2/4,3/4}{
   \draw[gray!30] (p\j\u)--(p\j\v);
  }
  \node at (0,-1.16) {\captiontext};
 \end{scope}
}
\foreach \u/\v in {0/1,0/2,0/3,0/4}{\draw[warmred,line width=1.5pt] (p0\u)--(p0\v);}
\foreach \u/\v in {0/1,1/2,1/3,1/4}{\draw[warmred,line width=1.5pt] (p1\u)--(p1\v);}
\foreach \u/\v in {0/1,0/2,1/2}{\draw[warmred,line width=1.5pt] (p2\u)--(p2\v);}
\foreach \j in {0,1,2}{
 \foreach \v/\ang in {0/90,1/18,2/-54,3/-126,4/162}{
  \node[vertex,inner sep=1.6pt] at (p\j\v) {};
  \node at ($(p\j\v)+(\ang:.22)$) {\scriptsize$\v$};
 }
}
\begin{scope}[xshift=7.4cm,yshift=-.75cm]
 \fill[softblue] (0,0)--(2.6,0)--(2.6,1.8)--cycle;
 \draw[inkblue,thick] (0,0)--(2.6,0)--(2.6,1.8)--cycle;
 \draw[->] (-.35,0)--(3.05,0) node[right] {$x$};
 \draw[->] (0,-.3)--(0,2.1) node[above] {$y$};
 \draw[gray,dashed] (0,1.8)--(2.6,1.8);
 \node[left] at (0,1.8) {$a$};
 \node[below left] at (0,0) {$s$};
 \node[below] at (2.6,0) {$a$};
 \node[below right=2pt] at (0,0) {$S_1$};
 \node[above right] at (2.6,0) {$T$};
 \node[above] at (2.6,1.8) {$S_0$};
 \fill[inkblue] (1.9,.67) circle (1.6pt);
 \node[left=3pt] at (1.9,.67) {$M(x,y)$};
\end{scope}
\end{tikzpicture}
\caption{Two kinds of price vectors suffice. In the three labelled
drawings, red edges have upper price $a$ and all other edges have
upper price $s=1-a$. The right-hand triangle represents
$s\le y\le x\le a$: every $M(x,y)$ is a weighted average of the
three corners. The bounds for $S_0$ and $S_1$ are the same after
swapping vertices $0$ and $1$.}
\label{fig:subtwo-two-models}
\end{figure}

\section{Dense sets and the proof}\label{sec:general-subtwo}
Ordinary packet coloring can lose colors on small odd sets. We remove
that loss by treating dense triangles and dense sets of five vertices
separately. Contract these sets, color the resulting packet graph, and
then put their internal edges back. A packet leaving a contracted set
remembers its original endpoint. This information tells us which
internal edges can use its color.

Throughout this section put
\[
 M=\frac{1691031}{1000000},\qquad
 \kappa=\frac{1999}{1000},\qquad
 D_3=3M-2\kappa=1.075093,\qquad
 D_5=5M-4\kappa=0.459155.
\]
These are rational constants. In particular,
$h_\infty<M<\kappa$ and $7M/6<\kappa$.
Costs in the next two lemmas are divided by a common local bin bound.
Accordingly, fractional packing programs allow nonnegative real type
demands. Scaling every demand scales their optimum by the same factor;
the patterns themselves still have integer type counts.

\begin{lemma}[Finding the dense sets]\label{lem:subtwo-blocks}
Assign a nonnegative cost $t_{uv}$ to each pair of distinct vertices.
Suppose each cost is at most one, every vertex has incident cost at
most $M$, and the three costs of every triangle sum to at most two.
There are disjoint sets, each of size three or five, with these
properties:
\begin{enumerate}[leftmargin=*,itemsep=2pt]
\item A selected three-set has internal cost greater than $\kappa$;
      a selected five-set has internal cost greater than $2\kappa$.
\item After these sets are contracted and their internal pairs removed,
      every quotient vertex has incident cost at most $M$.
\item Every odd quotient set $S$ of size at least three has internal
      cost at most $\kappa(|S|-1)/2$.
\end{enumerate}
Quotient costs are sums of original costs between the two sets. A
contracted three-set has boundary cost less than $D_3$, and a
contracted five-set has boundary cost less than $D_5$.
\end{lemma}
\begin{proof}
Call a triangle dense if its cost exceeds $\kappa$.
Two dense triangles cannot share just one vertex: each contributes
more than $\kappa-1$ there, while $2(\kappa-1)>M$.
If they share an edge, their combined cost is at most the sum of
the incident costs at its endpoints, at most $2M<2\kappa$.
Thus all dense triangles are disjoint. Contract them. The boundary
cost of one such triangle is less than $3M-2\kappa=D_3$.

A dense triangle in this quotient has exactly one contracted vertex
and two ordinary vertices. Indeed, an all-ordinary triangle was
already considered, while a triangle containing at least two blocks
has cost at most $(2D_3+M)/2<\kappa$.
These new dense triangles are also disjoint. Sharing a quotient
edge is impossible by the $2M<2\kappa$ degree argument. If two
share their old block and no other vertex, their seven original
vertices contain more than $3\kappa$ cost, contrary to
$7M/2<3\kappa$. If they use different old blocks and share one
ordinary vertex, their nine original vertices contain more than
$4\kappa$, contrary to $9M/2<4\kappa$.
Contract each new triangle to a five-set. Its internal cost exceeds
$2\kappa$, and its boundary cost is less than $D_5$.

Among vertices still uncontracted, consider five-sets of cost greater
than $2\kappa$. They are pairwise disjoint. For two such sets with
intersection size one, two, or three, the cost of their union plus
that of their intersection is at most, respectively,
\[
 9M/2,\qquad 4M+1,\qquad 7M/2+\kappa.
\]
Each is less than $4\kappa$. In the third case their common triangle
has cost at most $\kappa$, since it was not contracted at the first
step. If the intersection has four vertices, summing the incident
costs at those four vertices counts all the costs of the two
five-sets, and possibly more. This contradicts
$4M<4\kappa$. Contract these disjoint five-sets as well.

Every remaining ordinary vertex has degree at most $M$, and the
block bounds $D_3,D_5$ are smaller than $M$.
An all-ordinary triangle is not dense. A triangle with one old
three-block and two ordinary vertices was treated at the second
step. A triangle with a five-block and two ordinary vertices has
cost at most $D_5+1<\kappa$. Triangles with at least two blocks
have cost at most $(2D_3+M)/2<\kappa$.

For a quotient five-set containing $a$ three-blocks and $b$ five-blocks,
expand those blocks and subtract their internal costs. Its cost is
less than
\[
 \frac{5M}{2}+(M-\kappa)(a+2b).
\]
If it contains a block, this is at most $7M/2-\kappa<2\kappa$.
All-ordinary dense five-sets were contracted. For any odd quotient
set of size $s\ge7$, the same expansion and subtraction bound its
cost by $sM/2$, and
\[
 \frac{sM}{s-1}\le\frac{7M}{6}<\kappa.
\]
This proves every required odd-set inequality.
\end{proof}

\begin{lemma}[Putting back a triangle]\label{lem:subtwo-triangle}
Let a three-vertex weighted graph have local fractional bin costs
at most one. Write $t_1,t_2,t_3\le1$ for its separate bundle costs,
with bundle $i$ opposite vertex $i$, and suppose
$\kappa<T=t_1+t_2+t_3\le2$.
For each $i$ reserve a disjoint class of $q_i$ colors that may be
used internally only by bundle $i$. If
\[
 0\le q_i\le M-T+t_i,
\]
the internal demand fits fractionally, together with these reserved
classes, in at most $1+M/2<\kappa$ colors.
\end{lemma}
\begin{proof}
Put $s=T-M>0$. In the $q_i$ reserved colors, pack a $q_i/t_i$
fraction of bundle $i$. The assumption gives
$q_i/t_i\le1-s$, since $t_i\le1$.
Thus a common $s$ fraction remains in every bundle.
The entire triangle has a fractional coloring with $3/2$ colors:
take one half of each of its three local star packings. Every edge
then receives total mass one. Color the common fraction in $3s/2$
colors and pack the remaining demand of each bundle separately.
Including the reserved classes, the total is at most
\[
 T-sT+\frac32s
   =T-(T-M)(T-3/2).
\]
This expression increases on $\kappa\le T\le2$, and at $T=2$
it is $1+M/2$. Empty patterns may fill any unused colors.
\end{proof}

\begin{figure}[htbp]
\centering
\begin{tikzpicture}[x=1cm,y=1cm,font=\small,>=Stealth]
 \foreach \i/\ang in {0/90,1/18,2/-54,3/-126,4/162}{
   \node[vertex] (v\i) at (\ang:1.15) {};
   \node at (\ang:1.48) {$\i$};
 }
 \foreach \i/\j in {0/1,1/2,2/3,3/4,4/0,0/2,0/3,1/3,1/4,2/4}
   \draw[gray!40] (v\i)--(v\j);
 \draw[very thick,inkblue] (v1)--(v2);
 \draw[very thick,inkblue] (v3)--(v4);
 \draw[very thick,inkblue,->] (v0)--(0,2.15);
 \node[align=center] at (0,-2.0) {In a boundary color labelled $0$,\\
 internal edges avoid vertex $0$.};
 \draw[->,thick] (2.05,0)--(3.35,0);
 \node[above] at (2.7,.05) {contract};
 \node[circle,draw=inkblue,fill=softblue,minimum size=1.15cm] (B)
       at (4.65,0) {$B$};
 \draw[very thick,inkblue,->] (B)--(4.65,2.15);
 \node[right] at (4.7,1.4) {label $0$};
 \draw[gray!60,thick] (B)--(6.15,-.7);
 \draw[gray!60,thick] (B)--(3.55,-1.05);
 \node[align=center] at (4.65,-2.0) {Each outgoing packet keeps\\
 its original endpoint label.};
\end{tikzpicture}
\caption{The information retained by a contracted block. The blue
internal matching is one allowed use of a color reserved at vertex
$0$; general allowed patterns may contain several edges at a vertex,
provided their total weight is at most one.}
\label{fig:subtwo-labels}
\end{figure}

\subsection{Putting back a five-vertex set}

\begin{lemma}[Dense five-vertex palette extension]
\label{lem:subtwo-five}
Put $\kappa=1999/1000$ and $M=1691031/1000000$.
Consider repeatable item types on five vertices whose demand at each
vertex has a fractional unit-bin packing of mass at most one.
Let $t_e$ be the separate fractional packing cost of bundle $e$, let
$T=\sum_e t_e$, and let $q_v\ge0$ satisfy
\begin{equation}\label{eq:subtwo-five-hypotheses}
 T\ge2\kappa,\qquad d_v+q_v\le M,
 \qquad d_v:=\sum_{e\ni v}t_e.
\end{equation}
Give the internal types nonnegative prices, normalized so that the
largest price of a feasible whole pattern is one. Let $C$ be the
price of the prescribed demand, and let $r_v$ be the largest price
of a pattern avoiding $v$. Then
\begin{equation}\label{eq:subtwo-five-dual}
 C+\sum_vq_v(1-r_v)<\kappa.
\end{equation}
Consequently, a palette of mass $\kappa$ can meet the internal demand
when a disjoint class of mass $q_v$ is reserved for patterns avoiding
each vertex $v$.
\end{lemma}
\begin{proof}
Theorem~\ref{thm:subtwo-local-five} gives $C\le c:=157/80$.
Write $m_e$ for the largest price of a bundle packet and put
$Z=\sum_vq_v(1-r_v)$. Separate bundle packing gives
$C\le\sum_e m_et_e$.
We will combine these two bounds by a matching on six vertices.

For cleaner exact arithmetic, relax $M$ to
$M_0=16911/10000$. Put
\[
 \lambda=\max\{2,T/2\},\qquad Q_0=5M_0-2T.
\]
Summing \eqref{eq:subtwo-five-hypotheses} gives
$Q:=\sum_vq_v\le Q_0$. If $Q_0=0$, then $Z=0$ and the conclusion
follows from $c<\kappa$. Otherwise choose
\[
 k=1+\frac{\lambda-M_0}{Q_0}>1.
\]
Add a vertex $\infty$, give an internal edge $e$ amount $t_e$,
and give edge $\infty v$ amount $kq_v$. These amounts have a
fractional proper edge coloring of mass at most $\lambda$.
Here are the degree and odd-set checks required by Edmonds'
matching theorem~\cite{Edmonds65}.

An original vertex has degree
$d_v+kq_v\le M_0+(k-1)Q\le\lambda$.
Also
\[
 kQ\le Q_0+\lambda-M_0\le\lambda-1,
\]
since $Q_0\le5M_0-4\kappa=4595/10000<M_0-1$.
This bounds the new vertex's degree. Every original triangle has
cost at most two by Lemma~\ref{lem:general-two-bags}; every triangle
containing $\infty$ has cost at most $kQ+1\le\lambda$, since
$t_e\le1$. The original five-set has cost $T\le2\lambda$.
For a five-set containing $\infty$ and omitting $v$, write
$R=Q-q_v$. The other four original degree bounds imply
$2T-d_v+R\le4M_0$. The set's internal cost is therefore at most
\[
 T-d_v+kR\le4M_0-T+(k-1)R
 \le3M_0-T+\lambda\le2\lambda.
\]
The last inequality uses $T+\lambda\ge2\kappa+2>3M_0$.
These are all odd sets of size at least three on six vertices.

Assign price $m_e$ to internal edge $e$ and price $1-r_v$ to
$\infty v$. Every matching has price at most one. Without an edge
at $\infty$, its bundle packets form one feasible pattern. With
$\infty v$, its other packets avoid $v$ and have total price at
most $r_v$. Averaging over the fractional matching coloring gives
\[
 \sum_e m_et_e+kZ\le\lambda.
\]
It follows that
\begin{equation}\label{eq:subtwo-one-variable-transfer}
 C+Z\le(1-1/k)C+\lambda/k
       \le c+\frac{\lambda-c}{k}.
\end{equation}

For $2\kappa\le T\le4$, we have $\lambda=2$ and
$k\ge7684/4595$. Thus the last expression is at most
$1220173/614720<\kappa$.
For $T\ge4$, denote that expression by $F$. Substitution gives
\[
 F=\frac{-2T^2+(5M_0+c)T-2cM_0}{8M_0-3T}.
\]
Its denominator is positive because $T\le5M_0/2$.
One completed square proves the required strict inequality:
\[
 (\kappa-F)(8M_0-3T)
 =2\left(T-\frac{3283}{800}\right)^2
       +\frac{4423}{40000000}>0.
\]

Finally, the reserved palette can supply price
$\kappa-\sum_vq_v+\sum_vq_vr_v$ for every such price vector.
Inequality~\eqref{eq:subtwo-five-dual} says this exceeds the demand
price. Finite linear programming duality gives the desired fractional
packing. Nonnegative prices suffice because deleting items from a
feasible pattern preserves feasibility.
\end{proof}

\begin{proof}[Proof of Theorem~\ref{thm:general-bp-subtwo}]
Put $n=\bp(G)$. Apply the grouping and rank rounding in
Lemma~\ref{lem:general-packet-preprocess}. Retained objects
fit locally in
\[
 N=\left\lceil\frac{n}{1-2b}\right\rceil+1
\]
bins, and each original parallel bundle has a bounded number of
non-tiny types and at most one tiny remainder object. The parameter
$b$ will be sufficiently small and at most $10^{-7}$.
The set-aside objects will use a fresh palette at the end.

On each original pair put $t_{uv}=F(P_{uv})/N$. These costs are at
most one. Projecting each local bin onto its neighbors and applying
the harmonic charge bound shows that their sum at a vertex is at
most $h_\infty<M$. Lemma~\ref{lem:general-two-bags} gives triangle
cost at most two. Apply Lemma~\ref{lem:subtwo-blocks}.

For a block $B$ and $v\in B$, let $q_v$ be the charge of its retained
external objects, divided by $N$. Use charge $a(w)$ except that
an external tiny object has charge $w/(1-b)$. In a local bin,
keep each internal neighbor projection together and leave each
external object separate. These weights still sum to at most one.
The tiny modification raises their harmonic charge by at most
$b/(1-b)$, so
\begin{equation}\label{eq:subtwo-local-charge}
 \sum_{u\in B\setminus\{v\}}t_{uv}+q_v\le M.
\end{equation}
The choice of $b$ permits this exact inequality. For example, the
harmonic recurrence and its geometric tail give
\[
 h_\infty\le\frac{2759290}{1631721},\qquad
 \frac{2759290}{1631721}+\frac{10^{-7}}{1-10^{-7}}<M.
\]
The first bound is the sum of the first five terms plus twice the
sixth term. Thus the total boundary charge of a three-block is
less than $D_3N$, and that of a five-block is less than $D_5N$.

Form integer packets separately on each original pair, using a
basic optimum fractional bundle packing as before. Only then
contract their endpoints. In particular, a quotient packet never
mixes different original endpoint pairs. Delete internal packets
from the quotient; all internal objects will be colored when their
block is restored. The ordinary quotient packet graph $H$ has
fractional chromatic index at most $\kappa N$, by
Lemma~\ref{lem:subtwo-blocks} and
\eqref{eq:general-fractional-index}. Properly color it with
$\kappa N+o(N)$ colors using
Proposition~\ref{prop:general-ordinary-coloring}.

Let $R$ be the remaining external objects. Their multiplicity on a quotient pair is at most
$25(\lceil(D+1)/b\rceil+1)$: two blocks contain at most $25$
original pairs, and the preprocessing lemma bounds each pair separately. As in \eqref{eq:general-packet-charge}, every used
packet has charge at least one. Therefore, at an ordinary quotient
vertex the number of packet colors plus residual charge is at
most $MN$. At a three-block or five-block it is at most $D_3N$
or $D_5N$. Since all these constants are smaller than $\kappa$,
there is a fixed margin of missing colors.

Apply Lemma~\ref{lem:general-charge-finish} with one extra record
at every block: each used color has one original endpoint label,
and at most $q_vN+\delta N+o_\delta(N)$ colors have label $v$,
where $\delta>0$ is an accuracy parameter chosen sufficiently small
at the end. We justify this record explicitly.
Packet colors already have a unique label. The path switches in
Lemma~\ref{lem:weighted-moment-mix} preserve each packet and its original
endpoint. The coloring remains proper, so packets of a given label
still occupy distinct colors at that block. Their label count is
therefore unchanged.
After these switches, at each quotient vertex and color choose a
residual local set by first choosing one original endpoint label and
then one size class. If $x_{e,c}$ are the
fractions in that lemma, the total probability of choosing label
$v$ at a fixed color is at most
\[
 \sum_{\substack{e\in R\text{ with}\\\text{label }v}}q_ex_{e,c}
       +J\max_e x_{e,c},
\]
where $J$ is the fixed number of size classes. Summing over colors
bounds the expected label count by its residual charge plus $O(1)$.
For a fixed color, summing the displayed bound over at most five labels
adds at most $5J\max_e x_{e,c}=O(1/N)$ to the total charge contribution.
Thus the total choice probability is below one by the fixed margin.

Compress the profile distributions separately for each label by
Lemma~\ref{lem:profiles}; its repair profiles keep that label.
The smoothing step adds at most an arbitrarily small multiple of
$N$ to the expected label count. In the proof of
Lemma~\ref{lem:bp-mult-round}, retain only vertex lists whose label
counts are at most their means plus $N^{3/4}$, as well as satisfying
the existing progress test. Independent choices over colors and
Hoeffding's inequality show that these at most five extra tests
discard probability $\exp(-\Omega(\sqrt N))$. The proof's estimates
change only by the same $1+o(1)$ conditioning factor. There is no
union bound over all vertices: each vertex chooses from its own
retained lists. Deleting an edge later can only decrease a label
count. Choosing the endpoint label first prevents two residual labels
from sharing a color at a block. Using only colors missing in the
switched packet coloring prevents a residual label from sharing a
packet color there.
Packet charge and residual charge together pay for the asserted
$q_vN+\delta N+o_\delta(N)$ bound.

We now restore the blocks. Label classes at one block are disjoint;
pad them with unused colors if necessary. Their total size is below
$\kappa N$, up to $O(\delta N)+o_\delta(N)$.
A three-block satisfies Lemma~\ref{lem:subtwo-triangle}, by
\eqref{eq:subtwo-local-charge}. A five-block satisfies the
palette lemma, Lemma~\ref{lem:subtwo-five}, with its full
internal cost $T>2\kappa$. That lemma packs its internal demand
in $\kappa N$ fractional colors respecting the reserved classes.
To handle the small excess in actual label counts, keep at most
$\lfloor q_vN\rfloor$ colors of label $v$ in the reserved palette
program. The omitted label-colors are left unused by the internal
coloring. Add one fresh color for each omitted color, so that the
usable palette again has the required mass. At most
$5\delta N+o_\delta(N)+O(1)$ fresh colors suffice. The retained
classes satisfy the original local charge inequalities, so the
palette lemma is applied within its stated hypotheses. The same
fresh colors can be reused in different blocks, which are disjoint.

For the finite construction, leave the at most ten internal tiny
objects of a block for cleanup. Keep its original full costs and
density in the palette lemma. Every price on its non-tiny types
extends by zero to the omitted types; deleting zero-price items
changes neither the maximum pattern price nor any maximum price
avoiding a specified vertex. Thus the same local inequality proves
feasibility of the non-tiny pattern program. This step does not
apply the density test again after deleting tiny objects.
There are boundedly many non-tiny types and patterns. Take a basic
feasible solution and round its pattern masses down. The demand
equalities ensure distinct actual objects can fill the slots.
Only boundedly many objects per original vertex remain.

The blocks are disjoint, so their internal colorings coexist.
In a color labelled $v$, the internal pattern avoids $v$; all
external objects of that color meet the block only at $v$.
Thus all vertex-color loads stay at most one.
Use fresh palettes for the residual insertion leftovers, the
bounded internal leftovers, and the rank-grouping set-asides.
Their total additional cost is an arbitrarily small multiple of
$n$, plus an accuracy-dependent constant, by
Lemma~\ref{lem:general-cleanup}.
Choose $b$ and then the remaining fixed accuracy parameters small
enough. Since $N=n/(1-2b)+O(1)$, the result is
$(\kappa+\varepsilon)n+C_\varepsilon$.
The same greedy bound handles bounded $n$.
\end{proof}

\begin{algorithm}[Coloring below twice the local bin bound]
\label{alg:general-bp-subtwo}
For the finite construction, find optimal local packings by enumerating
partitions of each finite incident edge set, and put $n=\bp(G)$. Group and round objects, compute the original bundle
costs, and enumerate triangles and five-sets to find the disjoint
blocks in Lemma~\ref{lem:subtwo-blocks}. Form packets on original
pairs and color their quotient. Insert residual external objects
using profiles with one endpoint label per block and color.
Solve the bounded pattern program within each block, respecting
its boundary label classes, and round pattern masses down.
Color all remaining objects in fresh palettes, expand the packets
and objects, and check every vertex-color load.
\end{algorithm}

\section{Algorithms and asymptotic guarantees}\label{sec:algorithm-scope}
For rational weights, the exact general total-weight bound has a
deterministic polynomial-time algorithm.  It returns a feasible coloring
with at most $3\lceil\Delta_w(G)\rceil$ colors and has no accuracy
parameter; see Algorithm~\ref{alg:total}. For every fixed accuracy
$\varepsilon>0$, the simple-graph and bipartite bounds, and the
preparatory general bound of Theorem~\ref{thm:general-bp-improved},
have randomized algorithms with expected polynomial running time.
Theorem~\ref{thm:general-bp-subtwo} has the finite construction in
Algorithm~\ref{alg:general-bp-subtwo}; a polynomial running-time bound
for its additional label-preserving step is not claimed here. Every algorithm returns an explicit,
checked coloring.  For a leading constant $c$ and its stated parameter
$n$, their guarantee is $(c+\varepsilon)n+O_\varepsilon(1)$. Edges are
listed individually in the input; multiplicities are not supplied only
as binary counts.

In the polynomial-time constructions, local bin packings are computed approximately. Their extra factor is
absorbed into $\varepsilon$. These constructions assume no exact bin-packing oracle.
Profile choices use computable second moments and resampling.
The disjoint-option step uses the fixed-slack matching algorithm proved
in Appendix~\ref{sec:matching-algorithm}. The sorting, grouping, proper
edge-coloring, and small-edge rounding steps are direct polynomial-time
procedures. For the general $h_\infty+\frac25$ bound,
Section~\ref{sec:general-bp} explains the constant-size bundle programs
and cites the polynomial-time ordinary coloring algorithm.
Appendix~\ref{sec:algorithms} gives the remaining implementation.

Below fixed thresholds, decreasing-weight greedy cleanup supplies a
feasible coloring with only an accuracy-dependent additive cost.
The thresholds are independent of graph order, multiplicity, and minimum
edge weight. As usual, the fixed-accuracy guarantees imply the $o(n)$
statements: for accuracy $1/k$, choose increasing thresholds $N_k$, and
use relative error at most $1/k$ on $N_k\le n<N_{k+1}$.
We claim polynomial running time for fixed accuracy, not a polynomial
bound jointly in input length and $1/\varepsilon$.


\clearpage
\appendix
\section{Formulas for the uniform star and triangle}
\label{sec:subtwo-potential-tables}

We give the bounds used in Theorem~\ref{thm:subtwo-local-five}.
Each formula holds throughout $1/2\le a\le1$.
Put
\[
 A=\frac{205}{144},\qquad R=\frac{25}{36},\qquad
 s=1-a,\qquad q=\frac{2s}{3a},\qquad h=\frac{2a-1}{3a}.
\]
For both price vectors use $\rho_{34}=q$, $\rho_e=2/3$ for $e\ne34$,
and $h_0=h_1=h_2=h$, $h_3=h_4=0$.
The identities
\[
 qa=\frac{2s}{3},\qquad ha=\frac{a-s}{3},\qquad
 q+2h=\frac23
\]
explain this balancing choice: as the high price increases, some of
the charge is moved from a paired-star bound to the bound of one on
the price of a whole star.

The endpoint splits have the same form. Choose a number $\phi_v$
at each vertex and set
\begin{equation}\label{eq:subtwo-offset-split}
 z_{ue}=\min\{2m_e,\max\{0,m_e+\phi_u-\phi_v\}\}
                           \qquad(e=uv).
\end{equation}
The analogous formula at $v$ gives $z_{ue}+z_{ve}=2m_e$.
Thus $\phi_u>\phi_v$ moves some light-item profit toward $u$,
and the minimum and maximum keep both endpoint shares nonnegative.
For the uniform star $S_0$, set
$\phi_0=3(2a-1)/4$ and $\phi_1=\cdots=\phi_4=0$.
For the uniform triangle $T$, set
$\phi_0=\phi_1=\phi_2=(2a-1)/3$ and $\phi_3=\phi_4=0$.
The star $S_1$ uses the bounds for $S_0$ with vertices $0,1$ swapped.

We verify the inequalities of
Lemma~\ref{lem:subtwo-endpoint-potential} using these splits.
For brevity, write
\[
 B_v=Ac_v+h_v,\qquad
 f_e=(2-\rho_e)m_e+\rho_e+R(\alpha_{ue}+\alpha_{ve}).
\]
After the light and residual bounds $c,\alpha$ have been chosen,
we only need $P_v\ge B_v$ and $P_u+P_v\ge f_{uv}$.
The following checks also cover $a=1$, where some upper prices vanish.


\subsection{Four high spokes}\label{app:subtwo-four-high}
For the uniform star $S_0$, set $m_{0i}=a$ and
$m_{ij}=s$ for $1\le i<j\le4$. Put $z=5a/2-3/4$.
The split in~\eqref{eq:subtwo-offset-split} gives the root share $z$
on every spoke and equal shares $s$ on every core edge.
The no-heavy baselines and residual multipliers can be taken as
\[
\begin{aligned}
 c_0&=z-ha,&
 c_1=c_2&=\max\{2a-z-ha,\ s-hs\},&
 c_3=c_4&=2a-z,\\
 \alpha_{0,0i}&=z-2a/3,&
 \alpha_{i,0i}&=s/3,&
 \alpha_{i,ij}=\alpha_{j,ij}&=s/3\quad(ij\ne34),\\
 \alpha_{3,34}&=\alpha_{4,34}
   =u:=\max\{(1-q)s,\ 2a-z-qa\}.
\end{aligned}
\]
The relations
\[
 z-{5a-1\over3}={5(2a-1)\over12}\ge0,\qquad q\le{2\over3}
\]
verify the off-edge, same-bundle, and pair-sum rows of the
endpoint-potential rule. The spoke and ordinary-core pair right sides
are, respectively,
\[
 F_s={4a\over3}+{2\over3}
          +R\left(z-{2a\over3}+{s\over3}\right),
 \qquad
 F_c={4s\over3}+{2\over3}+{2Rs\over3}.
\]
Define
\[
 P_H={F_c\over2}={133-97a\over108},\qquad
 F_1=(2-q)s+q+2R(1-q)s,
\]
\[
 P_L=\max\left\{P_H,\ {F_1\over2},\ A(2a-z)\right\},
 \qquad
 P_0={157\over40}-2P_H-2P_L,\qquad
 P_1=P_2=P_H,\quad P_3=P_4=P_L.                         \tag{S}
\]
The core edges other than $34$ are covered because their pair right
side is $2P_H$. For edge $34$, the second possible branch
$u=2a-z-qa$ has right side $F_2$ with
\[
                    P_H-{F_2\over2}
                         ={23(2a-1)\over144}\ge0;
\]
the first branch is covered by $P_L\ge F_1/2$.
Since $P_L\ge P_H$, all spoke pairs follow from
$P_0+P_H\ge F_s$.

For the leaf no-heavy rows, $P_L\ge Ac_3=Ac_4$ by construction.
The two possible baselines at leaves $1,2$ give
\[
 P_H-A(2a-z-ha)-h
 ={1318a^2-1689a+576\over1728a}>0,\qquad
 P_H-A(s-hs)-h
 ={61(1-a)(3a-1)\over432a}\ge0.
\]
The first numerator has discriminant $-183951$.
It remains to check the root no-heavy row and the spoke pair row.
For a candidate $Y$ in the maximum defining $P_L$, write
$P_0(Y)=157/40-2P_H-2Y$. Direct substitution gives
\[
\begin{aligned}
 P_0(P_H)-Ac_0-h
   &={8490a^2-9283a+2880\over8640a},&
 P_0(P_H)+P_H-F_s
   &={690a-317\over2160},\\
 P_0(F_1/2)-Ac_0-h
   &={36010a^2-39043a+10880\over8640a},&
 P_0(F_1/2)+P_H-F_s
   &={7570a^2-7757a+2000\over2160a},\\
 P_0(A(2a-z))-Ac_0-h
   &={5270a^2-6453a+2880\over8640a},&
 P_0(A(2a-z))+P_H-F_s
   &={781-230a\over4320}.
\end{aligned}                                                \tag{S-slack}
\]
All displayed slacks are nonnegative on $1/2\le a\le1$.
The linear numerators are positive there; the quadratic numerators
have positive leading coefficients and, in the left column from
top to bottom and then the middle right entry, discriminants
$-11630711,-42799351,-19069191,-388951$.
The maximum in~(S) only decreases $P_0$, so checking all three
candidates proves the two remaining rows. Hence
$2C\le\sum_vP_v=157/40$, where $C$ is the demand price in
Theorem~\ref{thm:subtwo-local-five}.



\subsection{Three high edges forming a triangle}\label{app:subtwo-high-triangle}
Give the edges $01,02,12$ upper price $a$ and the other seven
edges upper price $s$.
Let $x=\min\{2s,(2-a)/3\}$ and $y=2s-x$.  The split in~\eqref{eq:subtwo-offset-split} gives shares $(x,y)$
on a spoke, with $x$ at the triangle vertex, and equal shares on
triangle edges and on $34$.  Choose the residual multipliers to be
$(a/3,a/3)$ on a triangle edge, $(a/3,s/3)$ on a spoke, and
$(w,w)$ on $34$, where
\[
                 w=\max\{0,y-qs,qs/2\}.
\]
Direct substitution verifies the same-bundle, off-edge, and pair-sum
conditions of Lemma~\ref{lem:subtwo-endpoint-potential}; the $qs/2$
term is needed for the outer pair-sum condition.  The no-heavy
star bounds are $c_H=(a+1)/3$ at each triangle vertex and $c_O=s$
at each outer vertex.  Accordingly, write
\[
 B_H=A\frac{a+1}{3}+h,\qquad B_O=As,\qquad
 f_H=\frac{97a+36}{54},\qquad
 f_S=\frac23+\frac{4s}{3}+\frac R3.
\]
Here $f_H$ and $f_S$ are the heavy-pair right sides on a triangle
edge and a spoke.  For the outer edge define the two expressions
\[
 f_O^+=-\frac{319a^2-344a+50}{54a},\qquad
 f_O^{\rm pair}=(2-q)s+q+Rqs.
\]
When $a\leq4/5$, the actual outer right side is
$\max\{f_O^+,f_O^{\rm pair}\}$.  When $a\geq4/5$, it is
$f_O^{\rm pair}$.  The following identities show that its positive
part above the no-heavy baseline is always covered by $f_O^+$:
\[
 f_O^{\rm pair}-f_O^+
       =\frac{25(2a-1)(4a-3)}{54a},\qquad
 f_O^{\rm pair}-2B_O
       =\frac{(a-1)(139a-100)}{216a}.
\]
The first difference is nonpositive for $a\leq3/4$, and the
second is nonpositive for $a\geq3/4$; hence
$f_O\leq\max\{f_O^+,2B_O\}$ throughout $[1/2,1]$.

Now give each triangle vertex potential
$P_H=\max\{B_H,f_H/2\}$ and each outer vertex potential
$P_O=\max\{B_O,f_O^+/2\}$.  Triangle and outer edges are covered
by construction.  Every spoke is covered because
\[
 \frac{f_H+f_O^+}{2}-f_S
       =-\frac{(2a-1)(39a-50)}{108a}\geq0.
\]
Thus the endpoint potential lemma yields $2C\leq3P_H+2P_O$.
Expanding the two maxima gives four possible scalar bounds:
\[
\begin{array}{c|c|c}
 (P_H,P_O)&(3P_H+2P_O)/2&\text{rational upper bound}\\ \hline
 (B_H,B_O)&(903-205a-144/a)/288&187/96\\
 (B_H,f_O^+/2)&(4231-1937a-832/a)/864&1693/864\\
 (f_H/2,B_O)&(277-11a)/144&181/96\\
 (f_H/2,f_O^+/2)&199/54-347a/216-25/(54a)&847/432
\end{array}
\]
For the first, second, and fourth rows, AM--GM proves the stated
bounds using, respectively,
$29520>171^2$, $1611584>1269^2$, and
$347>(149/8)^2$.  The third row decreases with $a$.
All four rational bounds are at most $847/432$, and
\[
                   C\leq\frac{847}{432}
                         <\frac{157}{80}.
\]


\subsection{Formal verification of the stated bound}
\label{app:subtwo-formal-verification}
The exact local inequalities and an unrestricted proof of the final
asymptotic theorem have also been checked in Lean. The public theorem is
\begin{center}
\small\path{Sosa27.subtwoBPAsymptotic_unrestricted}.
\end{center}
It states that, for every $\varepsilon>0$, there is a threshold $N$
such that every finite loopless weighted multigraph with $\bp(G)\ge N$
satisfies
$\chi'_w(G)\le(1999/1000+\varepsilon)\bp(G)$.
The check uses only Lean's standard axioms for classical mathematics;
the final theorem has no added mathematical assumptions or unproved
placeholders. The formal proof implements the rounding step through
explicit distributions and variance bounds, so it is an independent
formal derivation of the same result, not a line-by-line transcription
of the presentation above. The scalar inequalities in this appendix
are given explicitly and do not require a numerical solver.

\section{Elementary tools}\label{sec:elementary}
We use K\H{o}nig's theorem: a bipartite multigraph of maximum degree $d$
has a proper edge coloring with $d$ colors. The following first lemma
also handles many edges of very small weight, even when the ordinary
degree is large.

\Needspace{13\baselineskip}
\begin{lemma}[Coloring edges of bounded weight]\label{lem:cleanup}
If $G$ is bipartite, every edge has weight at most $b<1$, and
$W_G(v)\le R$ for every vertex, then
\[
 \chi'_w(G)\le\left\lceil\frac{R}{1-b}\right\rceil.
\]
Without the bound $w_e\le b$, at most $2\lceil2R\rceil$ colors suffice.
\end{lemma}
\begin{proof}
There are no edges if $R=0$. Otherwise put $q=\lceil R/(1-b)\rceil$.
At each vertex list its edges in decreasing order of weight and cut the
list into consecutive blocks of $q$ edges; the last block may be shorter.
Replace the vertex by one new vertex for each block. An original edge now
joins the two new vertices representing its blocks at its two ends.
The new graph is bipartite and has maximum degree at most $q$. Give it
a proper $q$-edge-coloring and transfer the colors back.

Fix an original vertex $v$ and write its sorted weights as
$a_1\ge a_2\ge\cdots$. A color appears at most once in each block.
Its load is therefore at most
\[
 a_1+a_{q+1}+a_{2q+1}+\cdots,
\]
where the sum stops at the last existing term. The first term is at most
$b$. For every later term, all $q$ weights in the preceding block are
at least as large, so
\[
 q\,a_{tq+1}\le\sum_{j=(t-1)q+1}^{tq}a_j\qquad(t\ge1).
\]
Adding these inequalities shows that the load is at most
$b+W_G(v)/q\le1$. Thus the transferred coloring is feasible.

For the last assertion, properly color the edges heavier than $1/2$
with at most $\lceil2R\rceil$ colors. Apply the first assertion, with
$b=1/2$, to the other edges, using another $\lceil2R\rceil$ colors.
\end{proof}

\Needspace{10\baselineskip}





The next lemma is a rounding statement of the kind used in approximation
algorithms. Each job must eventually choose one option. A fractional
solution gives an option a value between zero and one, with the values
for each job summing to one. A capacity constraint limits the total
amount used by the chosen options. In our application, a job is an edge,
and its two constraints belong to its two endpoints.

\begin{lemma}[Choosing one option per job]\label{lem:two}
Suppose each job has a probability distribution over finitely many options.
Each option uses a nonnegative amount from at most two capacity constraints,
and each amount is at most $\tau$. If $b_R$ is the expected amount used
in constraint $R$, then there is a choice of one option per job using
at most $b_R+2\tau$ in every constraint.
\end{lemma}
\begin{proof}
Let $x_j$ be the current value of option $j$. Values zero and one will
remain fixed. Initially the values are the given probabilities.
Keep the equation that the options for each unfinished job sum to one.
For a capacity constraint, also keep its equation fixing the current load
until the first time
\[
 \sum_{j\text{ using that constraint},\ 0<x_j<1}(1-x_j)\le2.
\]
At that time drop the equation permanently. The sum measures how much the
remaining option values could increase if all of them were rounded to one.

Suppose $N>0$ option values are still strictly between zero and one, and
$r$ jobs are unfinished. Each such job has at least two of these values,
so $r\le N/2$. Their values sum to $r$, hence $\sum_j(1-x_j)=N-r$.
If $s$ capacity equations are still kept, counting each option in at most
two constraints gives $s<N-r$. Thus fewer than $N$ linear equations are
being kept for $N$ variables. Linear algebra gives a nonzero direction
in which all these equations stay unchanged. Move in that direction until
one variable first reaches zero or one. Repeat; each step fixes another
variable, so the process ends.

When a capacity equation is dropped, its load is still $b_R$. Its possible
increase afterward is at most $\tau\sum_j(1-x_j)\le2\tau$. A constraint
kept to the end has load exactly $b_R$. Every job has exactly one option
with final value one, which proves the result.
\end{proof}

\Needspace{20\baselineskip}
\subsection{Two theorems for making many choices}\label{sec:inputs}
We use two known theorems. The first chooses many local objects while
controlling all specified sums. The second chooses disjoint sets. We only
need their stated forms; their proofs can be found in the cited papers.

\begin{proposition}[Controlling many sums, Harris--Srinivasan]\label{prop:HS}
Let $Z_i$ be independent random choices, each with a finite set of possible
values. For each index $t$ let
\[
 F_t=\sum_i f_{t,i}(Z_i),\qquad 0\le f_{t,i}(z)\le1.
\]
Define
\[
 R=\max_t\E F_t,\qquad
 D=\max_{i,z}\sum_t f_{t,i}(z),\qquad C=\max\{2,R,D\}.
\]
There is one choice of all the $Z_i$ for which every sum satisfies
\[
 F_t\le\E F_t+
 O\bigl(\sqrt{R\log(C+2)}+\log(C+2)\bigr).
\]
The constant hidden in $O(\cdot)$ is absolute.
\end{proposition}

The quantity $R$ bounds the expected size of a sum. The quantity $D$
bounds the total contribution that one possible choice can make to all
the sums together. The error bound depends on these two quantities,
even if the number of choices and sums is very large. This is what lets
us obtain bounds independent of the number of vertices.

For completeness, this form follows from Theorem~4.7 of~\cite{HS} by
using upper bounds $c_t=\max\{1,\E F_t\}$, the parameter
$\widehat D=C+1$, and its error parameter $1/\widehat D$.
If $c_t\le\log\widehat D$, that theorem gives $O(\log(C+2))$.
Otherwise it gives
$c_t+O(\sqrt{c_t\log(C+2)}+1)$. Since $c_t\le R+1$, both cases
give the displayed bound.

\Needspace{18\baselineskip}
\begin{proposition}[Choosing disjoint sets, Delcourt--Postle]\label{prop:DP}
Let $A$ and $B$ be disjoint sets of elements. Consider a finite list of
sets, each containing exactly one element of $A$ and at most two elements
of $B$. The same set may occur several times in the list. Write $d(x)$
for the number of listed sets containing $x$, and $\Delta_2$ for the
largest number of listed sets containing any given pair of distinct
elements. All counts include repetitions.

Fix $\zeta>0$. There is a constant $a>0$ such that, for all sufficiently
large $D$, the conditions
\[
 d(x)\ge(1+D^{-a})D\quad(x\in A),\qquad
 d(y)\le D\quad(y\in B),\qquad
 \Delta_2\le D^{1-\zeta}
\]
guarantee a collection of pairwise disjoint listed sets covering every
element of $A$.
\end{proposition}

In our use, an element of $A$ represents an edge still to be colored.
An element of $B$ represents a place available at one of its endpoints.
A listed set is one possible choice for that edge. Disjointness means
that no edge is chosen twice and no place is used twice. The inequalities
say that each edge has more choices than any one place has competing
choices, and no pair of elements occurs together too often. This is
Theorem~2.6 of~\cite{DP}, restricted to sets of size at most three and
with no extra forbidden collections.

\section{Describing the choices at one vertex}\label{sec:local}
We first consider only one vertex, or equivalently one finite set of items
with weights $w_e\le1$. A subset is feasible if its total weight is at
most one. We say a vector $x=(x_e)$ is an \emph{average of feasible sets}
if there are nonnegative numbers $p_S$ such that
\[
 \sum_Sp_S=1,\qquad x_e=\sum_{S\ni e}p_S,
\]
where only feasible sets $S$ are used. In probability terms, choose $S$
with probability $p_S$; then item $e$ is chosen with probability $x_e$.
Write $\mathcal K$ for all such vectors. At a graph vertex we write $\mathcal K_v$.

Fix a small positive number $\beta$. Call an item \emph{large} if its
weight is at least $\beta$, and \emph{small} otherwise. A feasible set
has at most $K=\lfloor1/\beta\rfloor$ large items. We describe a bin by
slots of sizes $b_1,\ldots,b_K$, allowing zero-size slots, and by a
space of size $R$ reserved for small items. The description is valid if
\[
 \sum_i b_i+R\le1.
\]
A large item may occupy a slot only if its weight is no greater than
the slot size. Several small items may share the reserved space as long
as their total weight is at most $R$. See Figure~\ref{fig:bin-description}.





We next replace a possibly long list of profiles by a list whose
length depends only on $\beta$ and a chosen error $\delta$. It will be
enough to assign fractions of items to these descriptions. The later
rounding step will make each edge a single, whole choice.

\begin{lemma}[A bounded list of profiles]\label{lem:profiles}
Fix $0<\delta<1$ and $0<\beta\le r_0=\delta/4$. If $x\in\mathcal K$,
there are at most $M=M(\beta,\delta)$ valid profiles, with
probabilities $\lambda_P$, and nonnegative assigned fractions $a(e;P,i)$
with the following properties. Put $y_e=(1-\delta)x_e$. Then
\begin{align}
 \sum_{P,i}a(e;P,i)&=y_e,\label{eq:marginal}\\
 \sum_e a(e;P,i)&\le\lambda_P
       &&\text{for each large-item slot }i,\label{eq:slot}\\
 \sum_{e:w_e<\beta}w_ea(e;P,R)&\le\lambda_PR
       &&\text{for the space reserved for small items}.\label{eq:reservoir}
\end{align}
Large items are assigned only to slots that can hold their weight. Small
items are assigned only to positive reserved spaces, each of which has
size at least $r_0$. Every individual assigned fraction is at most $y_e$.
\end{lemma}

The first equation preserves a $1-\delta$ fraction of each item's
original probability. The next two equations bound the amount assigned
to each place. Notice the different units: a large-item slot holds one
item, while the space for small items is measured by total weight.

\begin{proof}
Choose a random feasible set whose inclusion probabilities are $x_e$.
Give each large item in the chosen set a slot equal to its weight, and
put $R=1-$ the sum of those slot sizes. The small items in the set fit
in this remaining space.

For each distinct positive slot size $b$, let $m(b)$ be the expected
number of slots of that size. This is a finite list of nonnegative
numbers with $\sum_bm(b)\le K$. Put $q=\delta/4$. We keep a short
list of slot sizes as follows. Keep the largest size. Scan toward
smaller sizes, and keep the first size for which the sum of the $m(b)$
at smaller sizes than the previous kept size reaches $q$. Repeat, and keep the
smallest size at the end. If there are no slots, the list is empty.

Each completed step uses at least $q$ of the total sum $\sum_bm(b)$.
Thus at most $L=\lfloor K/q\rfloor+2$ sizes are kept. The sum of the
$m(b)$ strictly between two consecutive kept sizes is less than $q$.
Round every slot size down to the largest kept size no greater than it.
For any threshold $t$, this lowers the expected number of slots of size
at least $t$ by less than $q$: only sizes between the two kept values
around $t$ can cross below $t$.

Also round each reserved space down to $r_0\lfloor R/r_0\rfloor$.
Every resulting profile is valid, and every positive reserved
space has size at least $r_0$. Sort the slots and pad with zeros.
There are now only boundedly many descriptions, since there are at most
$L+1$ choices for each of $K$ slot sizes and at most
$\lfloor1/r_0\rfloor+1$ choices for the reserved space.

Give this rounded distribution total probability $s=1-\delta$.
Add probability $sq$ on a profile consisting of one slot of size
one, and probability $sr_0$ on a profile consisting of reserved
space of size one. These probabilities sum to at most
\[
 s(1+q+r_0)=(1-\delta)(1+\delta/2)\le1.
\]
Put any probability left over on the empty description. In particular,
one may use
\[
 M\le(L+1)^K\bigl(\lfloor1/r_0\rfloor+1\bigr)+3.
\]

The extra size-one slot repairs every lost threshold count. The total
slot capacity, counted as fractions of slots, for sizes at least $t$
is at least
\[
 s\sum_{b\ge t}m(b)
 =\sum_{e:w_e\ge\max\{t,\beta\}}sx_e
 \qquad(0<t\le1).
\]
Assign each large item its required fraction $sx_e$, processing items
from largest weight to smallest. Split fractions across compatible
slots as needed. This process cannot fail. If it stopped at an item
of weight $t$, all fractions already placed came from items of weight
at least $t$, and all slots of size at least $t$ would be full. That
would contradict the displayed inequality. This proves
\eqref{eq:marginal} and \eqref{eq:slot} for large items.

The total expected reserved space is
$Q=\sum_P\lambda_PR(P)$. Rounding a space down loses less than $r_0$,
and the added full-space description repairs that loss. Hence
\[
 Q\ge s(\E R-r_0)+sr_0
 \ge\sum_{e:w_e<\beta}w_e y_e.
\]
Assign a small item's fraction among the reserved spaces in proportion
to their expected capacities:
\[
 a(e;P,R)=y_e\frac{\lambda_PR(P)}{Q}\quad\text{if }Q>0.
\]
These fractions sum to $y_e$. Their total weight in description $P$
is at most $\lambda_PR(P)$, giving \eqref{eq:reservoir}. If $Q=0$,
all small-item demands are zero. Finally, since all fractions are
nonnegative and sum to $y_e$, each one is at most $y_e$.
\end{proof}

The assigned fractions need not describe one actual set of distinct
items in each profile. We only need their sums and capacity
bounds. When we return to graph edges, the rounding step will ensure
that no edge is chosen more than once.


\section{Choosing compatible profiles and whole edges}\label{sec:round-proof}
\begin{proof}[Proof of Theorem~\ref{thm:round}]
All constants hidden in $O(\cdot)$ below may depend on the fixed parameters,
but not on $n$ or on the number of vertices.


\subsection*{Dividing the vertices into groups}
Choose a fixed integer $r\ge\max\{2,4/\eps\}$. Assign each vertex
independently and uniformly to one of $r$ groups. For a vertex $v$ and
group $j$, let
\[
 L_{v,j}=\sum_{u:uv\in E(G)}w_{uv}\one_{\{u\text{ is in group }j\}}.
\]
This is the weight from $v$ to group $j$, and its expectation is
$W_G(v)/r$. Apply Proposition~\ref{prop:HS} to all these sums. One
choice at a source vertex $u$ contributes at most one to each sum,
because $G$ is simple, and its contributions over all sums add up to
$W_G(u)\le n$. Thus the proposition has $R\le n/r$ and $D\le n$.
It gives a division into groups $V_1,\ldots,V_r$ for which
\[
 L_{v,j}\le \frac{W_G(v)}r+O(\sqrt{n\log n})
 \quad\text{for every }v,j.
\]
This bound allows arbitrarily many edges of very small weight.

Let $H$ contain just the edges joining different groups. Set aside the
other edges for now. Their total weight at every vertex is at most
$n/r+O(\sqrt{n\log n})$. Restrict the proposed fractions to $H$.
They remain averages of feasible sets: intersect each original feasible
set with $E(H)$. If
\[
 T_v^G=\sum_{e\in E_G(v)}w_e\sum_cx_{e,c},\qquad
 T_v^H=\sum_{e\in E_H(v)}w_e\sum_cx_{e,c},
\]
then $\sum_cx_{e,c}\le1$ gives
\begin{equation}\label{eq:deleted-weight}
 T_v^G-T_v^H\le \frac nr+O(\sqrt{n\log n}).
\end{equation}
Until the final error calculation, work only with $H$: every edge,
incident-edge set, and sum below belongs to $H$.

\subsection*{Choosing profiles and measuring the lost weight}
Choose small positive constants $\delta,\theta,\sigma$, put $r_0=\delta/4$,
and choose a sufficiently small positive constant $\beta\le r_0$.
The precise requirements on these constants will be given at the end.
Call an edge large if $w_e\ge\beta$, and small otherwise.

Apply Lemma~\ref{lem:profiles} to every vertex $v$ and color $c$.
Write $\lambda_{v,c,P}$ for the probability of profile $P$, and
$a_{v,c}(e;P,i)$ for the fraction assigned to edge $e$ in its slot $i$.
For a small edge, $i$ instead denotes the space for small edges.
In either case these fractions sum, over all descriptions and places, to
\[
 y_{e,c}=(1-\delta)x_{e,c}.
\]
Use the same number $M$ of profile labels everywhere, adding empty
descriptions of probability zero if necessary. The bound on $M$ depends
only on $\beta,\delta$.

It is useful to give every label a positive probability. Set
\[
 \lambda'_{v,c,P}=(1-\theta)\lambda_{v,c,P}+\theta/M,
 \qquad \rho=(1-\theta)(1-\sigma).
\]
Initially, choose all descriptions $P_{v,c}$ independently according to
$\lambda'$. For an edge $e=uv$, a color $c$, and a place $i$ at $u$ and
a place $j$ at $v$, define
\[
 X(e,c,i,j)=
 \frac{\rho\,a_{u,c}(e;P_{u,c},i)a_{v,c}(e;P_{v,c},j)}
 {y_{e,c}\lambda'_{u,c,P_{u,c}}\lambda'_{v,c,P_{v,c}}}.
\]
If $y_{e,c}=0$, the value is zero. A positive value is a fractional option
to give $e$ color $c$ and use the stated places at its two ends. Large
edges use ordinary slots at both ends; small edges use the space for small
edges at both ends. Put
\[
 S_e=\sum_{c,i,j}X(e,c,i,j),\qquad
 \mu_e=\rho\sum_c y_{e,c}\le1.
\]
The factor $\lambda'$ in each denominator cancels the probability of
choosing that description. Summing the assigned fractions at each end
therefore gives $\E S_e=\mu_e$. More explicitly, for each color the two endpoint sums give $y_{e,c}^2$; division by $y_{e,c}$ leaves $y_{e,c}$.

For one color, the sum of the $X$ values is $O(1/n)$. Indeed,
$\sum_i a_{v,c}(e;P,i)\le y_{e,c}$,
$\lambda'_{v,c,P}\ge\theta/M$, and $y_{e,c}\le C_0/n$.
There are $O(n)$ colors, so $S_e=O(1)$ for every possible set of choices.
The terms from distinct colors are independent. Adding their variances
gives $\operatorname{Var}(S_e)=O(1/n)$.

At a vertex, measure the total error by adding the absolute errors of its
incident edges, multiplied by their weights:
\[
 Z_v=\sum_{e\ni v}w_e|S_e-\mu_e|.
\]
The Cauchy--Schwarz inequality gives
$\E|S_e-\mu_e|\le\sqrt{\operatorname{Var}(S_e)}$.
Since $\sum_{e\ni v}w_e\le n$, we obtain
\begin{equation}\label{eq:L1}
 \operatorname{Var}(S_e)=O(1/n),\qquad \E Z_v=O(\sqrt n).
\end{equation}

We must also respect the slots and the space for small edges. For an
ordinary slot, let $d(e)=1$ and let its capacity be $h(P,i)=1$; it can
hold at most one edge. For the space for small edges, let $d(e)=w_e$ and
$h(P,i)=R(P)$; its capacity measures total weight. Its load is the sum
of $d(e)X(e,c,i,j)$ over all options using it.
Fix a profile $P$ at $(v,c)$ and one of its places $i$. Averaging over
the choices at the other endpoints gives load at most
\begin{equation}\label{eq:meanload}
 \frac{\rho}{\lambda'_{v,c,P}}
       \sum_{e\ni v}d(e)a_{v,c}(e;P,i)
 \le (1-\sigma)h(P,i).
\end{equation}
Here we used the capacity bounds of Lemma~\ref{lem:profiles} and
$\lambda'\ge(1-\theta)\lambda$. We omit spaces of capacity zero.
Every remaining capacity is at least $r_0$.

\subsection*{Choosing the groups in order}
The groups $V_1,\ldots,V_r$ have no edges within them in $H$.
We choose the lists $P_v=(P_{v,c}:c\in\cC)$ one group at a time,
as shown in Figure~\ref{fig:group-passes}. All references to the original distribution of a profile list mean
the independent description choices with probabilities $\lambda'$.


\begin{figure}[htbp]
\centering
\begin{tikzpicture}[x=1cm,y=1cm,font=\small]
\foreach \xx/\name in {0/1,3.5/2,7/3}{
 \draw[fill=softblue,draw=inkblue,rounded corners] (\xx-.95,-.4) rectangle (\xx+.95,1.6);
 \node at (\xx,1.25) {$V_{\name}$};
 \node[vertex] (a\name) at (\xx-.45,.55) {};
 \node[vertex] (b\name) at (\xx+.45,.05) {};
 \draw[black!45,thick,dashed] (a\name)--(b\name);
}
\draw[inkblue,thick] (a1)--(a2);
\draw[inkblue,thick] (b1)--(b2);
\draw[inkblue,thick] (a2)--(b3);
\draw[inkblue,thick] (b2)--(a3);
\draw[inkblue,thick] (a1) to[bend left=24] (a3);
\node at (9,.6) {$\cdots\quad V_r$};
\node[align=center] at (3.5,2.6) {Keep edges between groups; set aside edges within groups.};
\node[panel,fill=softgreen,minimum width=1.9cm] (s1) at (0,-1.45) {choose $V_1$};
\node[panel,fill=softgreen,minimum width=1.9cm] (s2) at (3.5,-1.45) {choose $V_2$};
\node[panel,fill=softgreen,minimum width=1.9cm] (s3) at (7,-1.45) {choose $V_3$};
\draw[-{Latex},thick] (s1)--(s2);
\draw[-{Latex},thick] (s2)--(s3);
\node at (9,-1.45) {$\cdots\quad V_r$};
\end{tikzpicture}
\caption{Divide the vertices into a fixed number of groups. The dashed
edges have small total weight at each vertex and will be colored later.
Choose descriptions one group at a time, keeping all earlier choices.
There are no edges within a group in the remaining graph, so choices
made in one step can be treated as independent.}
\label{fig:group-passes}
\end{figure}

We use second moments to compute an upper bound on the remaining
absolute error. Choose a power of two $\kappa$ with $(2n^2)^{-1}<\kappa\le n^{-2}$,
and put
\[
 \phi(z)=\kappa\left\lceil\frac{\sqrt z}{\kappa}\right\rceil
 \qquad(z\ge0).
\]
Thus $\sqrt z\le\phi(z)<\sqrt z+\kappa$. After groups
$V_1,\ldots,V_k$ have been chosen, let $\mathcal F_k$ record their lists
and define
\begin{equation}\label{eq:computable-progress}
 U_e^{(k)}=\E[(S_e-\mu_e)^2\mid\mathcal F_k],\qquad
 Q_v^{(k)}=\sum_{e\ni v}w_e\phi(U_e^{(k)}).
\end{equation}
Every unchosen list in this expectation has its original distribution;
this is a way to measure progress, not a rule for the later choices.
The moments are computed by adding the first and second moments of the
independent color contributions. Section~\ref{sec:computable-progress}
gives the formulas and the rational implementation of $\phi$.
Initially $Q_v^{(0)}=O(\sqrt n)$, by~\eqref{eq:L1} and
$\sum_{e\ni v}w_e\kappa\le n\kappa\le1/n$.

Under the original distributions for the lists in the next group,
the tower property and the concavity of the square root give
\[
 \E[\phi(U_e^{(k)})\mid\mathcal F_{k-1}]
 \le\sqrt{U_e^{(k-1)}}+\kappa
 \le\phi(U_e^{(k-1)})+\kappa.
\]
Consequently
\[
 \E[Q_v^{(k)}\mid\mathcal F_{k-1}]
 \le Q_v^{(k-1)}+1/n.
\]
These inequalities replace the need to compute conditional expectations
of absolute values.

Put $t=\lceil n^{3/4}\rceil$. At step $k$, for each $u\in V_k$, retain only lists
for which $Q_u^{(k)}\le t$. Since no two vertices of $V_k$ are adjacent,
this condition depends on the list at $u$ and on earlier choices, but
not on any other list in $V_k$. We choose independently from the retained
lists, with their original probabilities rescaled to sum to one.

We show inductively that, just before step $k$, every unchosen vertex
satisfies
\[
 Q_v^{(k-1)}=O_r(\sqrt{n\log n}).
\]
Thus Markov's inequality, $\Prb(Y>t)\le\E Y/t$ for $Y\ge0$,
bounds the probability removed at each vertex of $V_k$ by the computable
number
\[
 p_k=\max_{u\in V_k}\frac{Q_u^{(k-1)}+1/n}{t}
     =O_r(n^{-1/4}\sqrt{\log n})=o(1).
\]
For an empty group take $p_k=0$. The constant can be chosen to cover all
$r$ steps. In particular, every retained distribution has positive total
probability for sufficiently large $n$. To sample it, repeatedly draw an
original list until its computed value of $Q_u^{(k)}$ is at most $t$.
The expected number of draws is at most $1/(1-p_k)$.

Apply Proposition~\ref{prop:HS} to the independent retained lists in
$V_k$, imposing the following bounds.
\begin{itemize}
\item For every vertex $v$, color $c$, possible profile $P$ at
$(v,c)$, and place $i$ in $P$, bound the load contributed by neighbours
of $v$ in $V_k$, on the assumption that $v$ uses $P$.
\item For every $v\notin V_k$, bound the part of $Q_v^{(k)}$ contributed
by neighbours in $V_k$.
\end{itemize}
The first kind of row is zero when $v\in V_k$. Every nonzero row is a
sum of nonnegative functions, one function of each neighbouring list
being chosen. For an error row, if $v$ was chosen earlier, its list is
fixed; if it will be chosen later, average over its original list.
Thus the contribution of a neighbour $u\in V_k$ is respectively
\[
 w_{uv}\phi\bigl((S_{uv}(P_u,P_v)-\mu_{uv})^2\bigr)
 \quad\text{or}\quad
 w_{uv}\phi\bigl(\E_{P_v}[(S_{uv}(P_u,P_v)-\mu_{uv})^2]\bigr).
\]
Each value is computable from first and second moments of the color
contributions. The contributions of all other neighbours are unchanged
during this step.

We give the size estimates needed for the proposition. In the following
table, a ``source'' is one vertex whose list we are choosing. A row is one
of the quantities to be bounded. The last column adds the contributions
of one fixed source choice over all rows. The scaling makes every
individual contribution at most one. Here $B_0$ is a sufficiently large
fixed constant.
\begin{equation}\label{eq:APtable}
\begin{array}{c|c|c|c}
 \text{quantity bounded}&\text{one source's contribution}
 &\text{scaling}&\text{sum over all rows}\\ \hline
 \text{load}&O(1/n)&\text{multiply by }n/B_0&O(n^2)\\
 \text{error bound}&O(w_{uv})&\text{divide by }B_0&O(n)
\end{array}
\end{equation}
All expected row sums after scaling are $O(n)$.
These estimates hold at every step.

Here is why the table does not involve the number of vertices. A source
vertex has at most one edge to a given vertex, because $H$ is simple.
Its contribution to one ordinary-slot load is therefore $O(1/n)$; its
contribution to a small-edge load is $O(w_e/n)$. For a fixed source list
and color, add the contributions over all possible profiles and places
at all neighbouring endpoints. There are only boundedly many description
labels and places at an endpoint. Summing the factors for the other
endpoint in the formula for $X$ cancels $y_{e,c}$, up to a fixed factor,
because those assigned fractions sum to $y_{e,c}$ and $\lambda'\ge\theta/M$.
The remaining sum at the source is bounded by its total assigned
ordinary-slot fraction plus its total assigned small-edge weight, divided
by its profile probability. Lemma~\ref{lem:profiles} bounds these by
$O(1/\beta)$ and $O(1)$, respectively, since
$\lambda/\lambda'\le1/(1-\theta)$. Thus the sum is $O(1)$ per color,
or $O(n)$ over all colors before scaling. This gives $O(n^2)$
in the first row of the table. For the second row,
$|S_e-\mu_e|=O(1)$ implies $\phi(U_e^{(k)})=O(1)$ for every
choice of fixed lists. Thus one edge contributes $O(w_e)$, and the sum
of the weights at a source is at most $n$.


To check the expected row sums, first use the original list distributions.
A load row has mean at most the full load in~\eqref{eq:meanload}, hence
at most one. Restricting a list to an event of probability at least
$1-p_k$ increases the expectation of each nonnegative contribution by
at most $1/(1-p_k)$. An error row is $O(n)$ for every possible set of
lists, since $\phi(U_e^{(k)})=O(1)$ and the incident weights sum to at
most $n$.
After scaling as in the table, Proposition~\ref{prop:HS} therefore has
$R=O(n)$ and $D=O(n^2)$ at every step.

For an error row, let $C$ be the total contribution from neighbours
outside $V_k$. This is a fixed nonnegative quantity at this step.
The original expected contribution from neighbours in $V_k$ is at most
$Q_v^{(k-1)}-C+1/n$, by the square-root inequality above.
Use this computable upper bound, divided by $1-p_k$, for the retained
lists. Proposition~\ref{prop:HS} therefore gives
\[
 Q_v^{(k)}
 \le C+\frac{Q_v^{(k-1)}-C+1/n}{1-p_k}
       +O(\sqrt{n\log n})
 \le\frac{Q_v^{(k-1)}}{1-p_k}+O(\sqrt{n\log n}).
\]
Only upper bounds on row means are needed; the sampling procedure does
not evaluate expectations under the retained distributions.
Section~\ref{sec:computable-progress} explains why the
Harris--Srinivasan algorithm can use these implicitly described lists.
There are only $r$ steps. For a vertex not yet chosen, induction from
$Q_v^{(0)}=O(\sqrt n)$ proves the claimed
$O_r(\sqrt{n\log n})$ bound, justifying the removed-probability bound
used above. More explicitly, at each induction step we may first use
$1/(1-p_k)\le2$. The constants then grow by at most a fixed factor
per step. Since there are only $r$ steps, choosing $n$ sufficiently
large makes every removed probability less than $1/2$, as required.
At the step when $v$ itself is chosen, the retained-list
condition gives $Q_v^{(k)}\le t$. At each later step the displayed
inequality applies again. Consequently, after the last step,
\[
 Z_v\le Q_v^{(r)}=O_r(n^{3/4})\qquad\text{for every }v.
\]

To check the load bounds, fix a vertex $v$, a color $c$, a possible
description $P$, and one of its places $i$. Under the original lists,
the expected contribution from neighbours in $V_k$ is exactly
\[
 m_k=\frac{\rho}{\lambda'_{v,c,P}}
       \sum_{\substack{e=uv\ni v\\u\in V_k}}
                d(e)a_{v,c}(e;P,i).
\]
In particular, $\sum_{k=1}^r m_k\le(1-\sigma)h(P,i)$ by
\eqref{eq:meanload}. Proposition~\ref{prop:HS}, followed by undoing
the load scaling, bounds the actual contribution from group $V_k$ by
\[
 \frac{m_k}{1-p_k}
 +O\left(\sqrt{\frac{\log n}{n}}+\frac{\log n}{n}\right).
\]
We imposed this bound for every possible $P$, so it holds for the
description eventually chosen at $v$, whether that choice was made
before, during, or after step $k$. Put $p=\max_{1\le k\le r}p_k=o(1)$.
Adding the $r$ bounds gives total load at most
\[
 \frac{(1-\sigma)h(P,i)}{1-p}
 +O_r\left(\sqrt{\frac{\log n}{n}}+\frac{\log n}{n}\right)
 \le(1-\sigma/2)h(P,i)
\]
for sufficiently large $n$, since every positive capacity is at least
$r_0$. This proves both required bounds. We can now turn the fractional options into a partial coloring.


An edge may currently have total assigned fraction $S_e>1$.
Divide all its option fractions by $\max\{1,S_e\}$, and call the new
fractions $\widehat X(e,c,i,j)$. Their sum and the unused fraction are
\[
 s_e=\min\{1,S_e\},\qquad r_e=1-s_e.
\]
This only decreases loads. Since $\mu_e\le1$, we have
$s_e\ge\mu_e-|S_e-\mu_e|$. Hence
\begin{equation}\label{eq:capture}
 \sum_{e\ni v}w_es_e
 \ge\sum_{e\ni v}w_e\mu_e-Z_v
 \ge\rho(1-\delta)\sum_{e\ni v}w_e\sum_cx_{e,c}
       -O(n^{3/4}).
\end{equation}

\subsection*{Making whole choices for the large edges}
For large edges, each option uses one ordinary slot at each end. We use
Proposition~\ref{prop:DP} to choose options without using a slot twice.
We also provide places for edges left uncolored, so that the number of
such edges can be controlled at every vertex.

Divide $[\beta,1]$ into finitely many weight intervals, assigning an edge
on a shared boundary to just one interval. Let $\ell_h,u_h$ be the lower
and upper endpoints of interval $h$, with
$u_h\le(1+\xi)\ell_h$, where $\xi>0$ is a small fixed constant.
Let $E(v,h)$ be the edges at $v$ in that interval, and put
\[
 R_{v,h}=\sum_{e\in E(v,h)}r_e,
 \qquad
 B_{v,h}=\left\lceil\frac{R_{v,h}}{1-\gamma}\right\rceil
                   +\lceil n^{2/3}\rceil,
\]
where $\gamma>0$ is another small fixed constant.
Create $B_{v,h}$ extra slots at $(v,h)$, reserved for uncolored edges.
Each such slot may receive at most one edge. These are counting devices;
they are separate from the slots belonging to colors.

Construct a family of three-element sets, allowing repeated copies of a
set. Introduce one new element for each large edge; call it the job
element for that edge. Let $A$ be all job elements. The remaining elements,
forming a set $B$, are the ordinary slots in the selected profiles
and the extra slots just created. Each three-element set will consist of
one job element and one slot at each endpoint of that edge.

Put $N=\lceil n\rceil^6$. For each color option $(e,c,i,j)$, add
\[
 \lfloor N\widehat X(e,c,i,j)\rfloor
\]
copies of the set containing the job element of $e$ and its two ordinary
slots. For an edge $e=uv$ in interval $h$, divide the unused fraction
$r_e$ equally among all pairs of extra slots, one at $(u,h)$ and one at
$(v,h)$. For each such pair add
\[
 \left\lfloor\frac{Nr_e}{B_{u,h}B_{v,h}}\right\rfloor
\]
copies of the set containing that pair and the job element of $e$.
Selecting a set of the first kind will color the edge; selecting a set
of the second kind will leave it uncolored.


\begin{figure}[htbp]
\centering
\begin{tikzpicture}[x=1cm,y=1cm,font=\small]
\node[panel,minimum width=2.8cm] (uc) at (-3.5,1) {slot at $u$\\color $c$};
\node[panel,minimum width=2.8cm] (vc) at (3.5,1) {slot at $v$\\color $c$};
\node[panel,fill=black!4,draw=black!50,minimum width=2.8cm] (ud) at (-3.5,-1) {extra slot at $u$\\leave uncolored};
\node[panel,fill=black!4,draw=black!50,minimum width=2.8cm] (vd) at (3.5,-1) {extra slot at $v$\\leave uncolored};
\node[draw,rounded corners,fill=white,minimum width=1.3cm,minimum height=.7cm] (e) at (0,0) {edge $e$};
\draw[inkblue,very thick] (uc)--(e)--(vc);
\draw[black!55,thick,dashed] (ud)--(e)--(vd);
\node[align=center] at (0,1.9) {two possible choices for the same edge};
\end{tikzpicture}
\caption{Each option is a set containing an edge label and one slot at each
endpoint. The upper option gives color $c$; the lower option leaves the
edge uncolored and uses an extra slot at each end. The two options share
the edge label, so at most one can be selected. Disjoint selected sets
also prevent a slot from being used by two edges.}
\label{fig:matching-options}
\end{figure}

% One large edge e has a central job element. Draw one available triple
% connecting it to an ordinary slot at u and an ordinary slot at v for a
% color c; draw another available triple connecting it to one extra slot
% at each endpoint. Use distinct line styles for 'give color c' and
% 'leave uncolored'. Each selected triple uses all three elements, so
% selecting one excludes every other triple touching any of its elements.

For an element $z$, let $d(z)$ be the number of sets containing it,
counting copies. Let $\Delta_2$ be the largest number of sets containing
any given pair of distinct elements, again counting copies. We claim
\begin{equation}\label{eq:degrees}
 \begin{aligned}
 d(a)&=N-O(n^2) &&(a\in A),\\
 d(b)&\le(1-\eta)N &&(b\in B),\\
 \Delta_2&=O(Nn^{-2/3}),
 \end{aligned}
 \qquad \eta=\min\{\sigma/2,\gamma\}>0.
\end{equation}
To check the first line, a job has $O(n)$ color options, since there
are $O(n)$ colors and boundedly many slots per color. Also,
$R_{v,h}\le n/\beta$, so $B_{v,h}=O(n)$ and the job has $O(n^2)$
pairs of extra slots. Before taking integer parts, the numbers of copies
for that job sum to $N(s_e+r_e)=N$. Each integer part loses less than
one, proving the first line.

For an ordinary slot, the number of sets using it is at most its
fractional load multiplied by $N$, which is at most $(1-\sigma/2)N$.
For an extra slot at $(v,h)$, the corresponding bound is
$NR_{v,h}/B_{v,h}\le(1-\gamma)N$. This proves the second line.

For the last line, consider the possible pairs. A job element and an
ordinary slot lie together in $O(N/n)$ sets: the slot fixes the color,
and the total fraction for one edge and color is $O(1/n)$.
A job element and an extra slot lie together in at most
$N/B_{v,h}=O(Nn^{-2/3})$ sets. Two ordinary slots at opposite endpoints
can occur together only for the unique original edge joining those
endpoints, and again the bound is $O(N/n)$. Two extra slots at opposite
endpoints occur together in at most $N/(B_{u,h}B_{v,h})$ sets, for the
same reason. All other pairs of slot elements occur together in no set,
as do pairs of job elements. This proves \eqref{eq:degrees}.

Take $D=\lfloor(1-\eta/2)N\rfloor$ and any fixed $0<\zeta<1/9$.
For large $n$, the hypotheses of Proposition~\ref{prop:DP} hold.
Indeed, every job element belongs to more than $D$ sets by a fixed
positive fraction of $N$, whereas every slot element belongs to fewer
than $D$ sets. The extra factor $D^{-a}$ required by that proposition
tends to zero. Finally, since $N$ has order $n^6$,
\[
 Nn^{-2/3}=o(D^{1-\zeta}).
\]
The proposition supplies pairwise disjoint sets covering all job
elements. Thus it chooses exactly one option for each large edge and
uses each ordinary or extra slot at most once.

At most $B_{v,h}$ edges of $E(v,h)$ are left uncolored, because each
uses a different extra slot. Their total weight is therefore at most
\begin{equation}\label{eq:residual}
 u_hB_{v,h}
 \le\frac{1+\xi}{1-\gamma}
                \sum_{e\in E(v,h)}w_er_e+O(n^{2/3}).
\end{equation}
The number of intervals depends only on the fixed constants, so summing
this bound over all intervals still gives an $O(n^{2/3})$ error.

\subsection*{Making whole choices for the small edges}
For a small edge $e$, use its color options of fractions $\widehat X$
and one additional option of fraction $r_e$ that leaves it uncolored.
These fractions sum to one. For the space for small edges of size $R>0$,
divide every load by $R$, so the capacity becomes one.
A color option then adds $w_e/R\le\beta/r_0$ to each of two capacity
constraints, one at each endpoint. The current fractional load on each
such space is at most $1-\sigma/2$.

At each vertex add one more constraint that records the total weight of
small edges left uncolored. The uncolored option for $e=uv$ adds $w_e$
to this sum at $u$ and at $v$. Thus every option contributes to at most
two constraints, and every contribution is at most
$\tau=\beta/r_0$; here $r_0<1$.
Choose $\beta$ small enough that
\[
 2\beta/r_0\le\sigma/2.
\]
Lemma~\ref{lem:two} now chooses one whole option for every small edge.
It increases each load by at most $2\beta/r_0$.
Consequently all the spaces for small edges obey their capacities, and
the small uncolored weight at $v$ is at most
\[
 \sum_{\substack{e\ni v\\w_e<\beta}}w_er_e+2\beta/r_0.
\]

For each vertex and color, a large edge occupies a distinct ordinary
slot and has weight at most that slot's size. The total weight of the
small edges fits the space set aside for them. A profile has
total size at most one, so the resulting partial coloring is feasible.
Every color option used had positive assigned fraction, and therefore
had $x_{e,c}>0$.

It remains to add the errors. Write
$T_v=T_v^H=\sum_{e\ni v}w_e\sum_cx_{e,c}$ and
$A=(1+\xi)/(1-\gamma)$. We have $T_v\le W_H(v)\le n$.
The bounds for large and small uncolored edges give total uncolored
weight at most
\[
 A\sum_{e\ni v}w_er_e+O(n^{2/3}).
\]
Subtract this from $W_H(v)$, use $r_e=1-s_e$, and then apply
\eqref{eq:capture}. The colored weight is at least
\[
 \begin{split}
 W_H(v)-A\sum_{e\ni v}w_er_e-O(n^{2/3})
 &\ge \sum_{e\ni v}w_es_e-(A-1)n-O(n^{2/3})\\
 &\ge T_v-\bigl[1-\rho(1-\delta)+(A-1)\bigr]n
                 -O(n^{3/4}).
 \end{split}
\]
Choose $\delta,\theta,\sigma,\xi,\gamma$ positive and small enough that
the quantity in square brackets is less than $\eps/2$.
Then set $r_0=\delta/4$ and choose $\beta$ as required above.
By~\eqref{eq:deleted-weight}, the colored weight in the original graph
is therefore at least
\[
 T_v^G-\left(\frac1r+\frac\eps2\right)n-O(n^{3/4})
 \ge T_v^G-\eps n-O(n^{3/4}).
\]
Here $1/r\le\eps/4$, and $\sqrt{n\log n}=O(n^{3/4})$.
The integer $r$ and all other choices depend only on $C_0,C_1,\eps$.
Taking $n$ sufficiently large proves the theorem, with no dependence
on the number of vertices.
\end{proof}

\section{Making missing colors agree}\label{sec:mix-proof}
We next study a proper edge coloring of a multigraph $H$ with
colors $1,\ldots,K$. A color is \emph{missing at $v$} if no edge at $v$
has that color. Write $F_v$ for the set of missing colors and $f_v=|F_v|$.
We will change the coloring while keeping every number $f_v$ fixed.

Here is the change we use. Fix two colors $a,b$ and keep only the edges
with these colors. Since the coloring is proper, every vertex has degree
at most two in the resulting graph. Its components are paths, even cycles,
and isolated vertices. On any whole path component, we may exchange $a$ and $b$ on every
edge. The coloring stays proper. Only the missing colors at the two ends
of the path change; see Figure~\ref{fig:path-switch}. This operation is usually called a \emph{Kempe switch}.

\begin{figure}[htbp]
\centering
\begin{tikzpicture}[x=1cm,y=1cm,font=\small]
\foreach \idx/\shift/\title/\missing in {0/0/Before/b,1/6.3/After/a}{
 \begin{scope}[xshift=\shift cm]
 \node at (2.1,1.1) {\title};
 \foreach \xx/\name in {0/u,1.4/x,2.8/y,4.2/v}{
  \node[vertex,label=below:$\name$] (\name\idx) at (\xx,0) {};
 }
 \ifnum\idx=0
  \draw[inkblue,very thick] (u\idx)--node[above] {$a$} (x\idx);
  \draw[warmred,very thick,dashed] (x\idx)--node[above] {$b$} (y\idx);
  \draw[inkblue,very thick] (y\idx)--node[above] {$a$} (v\idx);
 \else
  \draw[warmred,very thick,dashed] (u\idx)--node[above] {$b$} (x\idx);
  \draw[inkblue,very thick] (x\idx)--node[above] {$a$} (y\idx);
  \draw[warmred,very thick,dashed] (y\idx)--node[above] {$b$} (v\idx);
 \fi
 \node at (2.1,-1) {Both ends miss $\missing$.};
 \end{scope}
}
\draw[-{Latex},thick] (4.75,0)--(5.7,0);
\end{tikzpicture}
\caption{Exchanging two colors on a whole path component preserves a
proper edge coloring. Only the missing colors at the two ends change. In the path shown,
both ends miss the same color before the switch and the same color
afterward. The argument below also permits even-length paths.}
\label{fig:path-switch}
\end{figure}


The weights in the next lemma belong to an additional graph $Q$ whose
edges still need colors. We will give an edge fractions of the colors
missing at its endpoints. Some edges may receive a total fraction below
one, but the total weight lost at each vertex will be small. The proof
uses averages and squared errors to obtain this conclusion without any
lower bound on the edge weights.

\Needspace{22\baselineskip}
\begin{lemma}[Controlling available colors and lost weight]\label{lem:weighted-moment-mix}
Fix positive constants $\alpha,C,L$. Suppose a multigraph without loops $H$
has a proper coloring with an even number $K\ge4$ of colors. Let $F_v$
be the colors missing at $v$, and write $f_v=|F_v|$. Suppose
\[
 f_v\ge\alpha n\quad\text{for every }v,\qquad K\le Cn.
\]
Let $Q$ be another multigraph without loops on the same vertex set, with
nonnegative weights $q_e$, such that
\[
 W_Q(v):=\sum_{e\ni v}q_e\le n,\qquad
 \sum_{e=uv}q_e\le L\quad\text{for every pair }u,v.
\]
For sufficiently large $n$, path switches in $H$ give missing-color sets
and fractions $x_{e,c}$ with the following properties:
\begin{align}
 x_{e,c}>0&\ \Longrightarrow\ c\in F_u\cap F_v
                      &&(e=uv),\label{eq:moment-support}\\
 0\le x_{e,c}&\le \frac{C}{\alpha^2n},\qquad
 \sum_cx_{e,c}\le1,\label{eq:moment-small}\\
 \sum_{e\ni v}q_ex_{e,c}
 &\le\frac{W_Q(v)}{f_v}
       +O\bigl(n^{-1/2}(\log n)^{3/2}\bigr),
                       \label{eq:moment-load}\\
 \sum_{e\ni v}q_e\Bigl(1-\sum_cx_{e,c}\Bigr)
 &\le O\bigl(n^{3/4}(\log n)^{3/4}\bigr).
                       \label{eq:moment-loss}
\end{align}
The constants and the required lower bound on $n$ depend only on
$\alpha,C,L$, not on the number of vertices. No positive lower bound on
$q_e$ is required.
\end{lemma}

\begin{proof}
Write $a_{u,c}=\one_{\{c\in F_u\}}$. We choose a short list of pairings
of the colors and switch paths for each pairing. In addition
to averages of $a_{u,c}$, we keep track of averages of $a_{u,c}a_{u,d}$
for distinct colors $c,d$.

For a list $h$, write $\|h\|^2=\sum_i h_i^2$; $h^T$ denotes the
same list written as a row rather than a column.

\medskip\noindent\textbf{A list that mixes single colors and pairs.}
For a pairing $M$ of the $K$ colors, let $A_M$ average within each pair.
Let $B_M$ act on lists indexed by unordered pairs $\{c,d\}$ of distinct
colors. To apply $B_M$, independently exchange or keep the two colors
of every pair in $M$, apply the resulting permutation to the indices of
the list, and average over all these choices. Both matrices are symmetric,
have nonnegative entries, and have every row sum equal to one. Applying
either averaging operation twice has the same effect as applying it once.
Consequently
\[
 \|A_Mh\|^2=h^TA_Mh,\qquad
 \|B_Mh\|^2=h^TB_Mh.
\]

Choose $M$ uniformly at random. Direct expansion gives
\[
 \E\|A_Mh\|^2=\frac{K-2}{2(K-1)}\|h\|^2
 \le\tfrac12\|h\|^2
 \quad\text{when }\sum_i h_i=0.
\]
We prove the same statement for $B_M$. An entry of $\E B_M$ depends only
on whether its two unordered-pair indices are equal, meet in one color,
or are disjoint. The respective entries are
\[
 p=\frac{K+2}{4(K-1)},\qquad
 q=\frac1{4(K-1)},\qquad
 r=\frac1{2(K-1)(K-3)}.
\]
For example, a pair $\{i,j\}$ is itself paired in $M$ with probability
$1/(K-1)$ and is then unchanged. Otherwise it is unchanged with
probability $1/4$, giving $p$. To change just $j$ to a specified third
color, those two colors must be paired, one swap must be made, and one
must be omitted, giving $q$. To reach a disjoint pair, there are two
possible pairings of the four colors, and both swaps must be made,
giving $r$.

Every symmetric array $(h_{ij})_{i\ne j}$ is the sum of three mutually
perpendicular arrays: a constant array; an array of the form $b_i+b_j$
with $\sum_i b_i=0$; and an array with
$\sum_{j\ne i}h_{ij}=0$ for every $i$. To see this directly, first
subtract the average entry. The remaining row sums add to zero.
Subtract $b_i+b_j$, where $b_i$ is the remaining $i$th row sum divided
by $K-2$. The final array has every row sum zero. Perpendicularity follows
by summing over rows.

The matrix $\E B_M$ multiplies these three arrays, respectively, by
\[
 1,\qquad
 p+(K-4)q-(K-3)r=\frac{K-2}{2(K-1)},\qquad
 p-2q+r=\frac{K-2}{4(K-3)}.
\]
The last two numbers are at most $1/2$ for even $K\ge4$. Therefore
\[
 \E\|B_Mh\|^2=h^T(\E B_M)h\le\tfrac12\|h\|^2
 \quad\text{when }\sum_{i<j}h_{ij}=0.
\]

Take $s=14\lceil\log_2K\rceil$ independent uniform pairings. Apply these
two contraction bounds successively to all unit coordinate lists after
subtracting their averages. The expected sum of the squared errors for
the two products is at most
\[
 \bigl(K+\tbinom K2\bigr)2^{-s}<K^{-10}.
\]
Thus some fixed list of pairings has the following property for both
products: it sends any list to within $K^{-4}$ times its original norm
of the constant list with the same average. Indeed, if $P$ is either product and $J$ replaces every entry by the
average, the Cauchy--Schwarz inequality gives
\[
 \|(P-J)x\|^2\le\|x\|^2\sum_j\|(P-J)e_j\|^2,
\]
where the $e_j$ are the unit coordinate lists. Fix such a list of pairings.

\medskip\noindent\textbf{The sums we control.}
During the switches define
\begin{align*}
 z_{v,c}&=\sum_{e=uv}\frac{q_e}{f_u}a_{u,c},\qquad
 d_{v,c}=\sum_{e=uv}\frac{q_e}{f_u^2}a_{u,c},\\
 Z_{v,\{c,d\}}&=
       \sum_{e=uv}\frac{q_e}{f_u^2}a_{u,c}a_{u,d}\quad(c\ne d).
\end{align*}
The numbers $f_u$ do not change during path switches. For a fixed
color pairing $M$, flip an independent fair coin on every nontrivial
two-color path and switch exactly when its coin says so. At one vertex,
the missing indicators for a paired pair of colors are either equal,
and unchanged, or unequal, and fairly exchanged. The coins for different
color pairs are independent. Hence, conditional on the coloring before
the stage, the means of the three displayed lists after the stage are
\[
 A_Mz_v,\qquad A_Md_v,\qquad B_MZ_v.
\]
This statement concerns one vertex $u$ at a time. It does not require
independence between different vertices.

We now choose the switches while keeping every coordinate close to this
mean. Process the $K/2$ color pairs one at a time. At one step the
independent choices are the path coins for the current pair of colors.
Switching within one pair does not change the edges or path components
of another pair, since their colors are disjoint. Thus the independent
coins for the entire stage are fixed before this process begins.
For every coordinate of $z,d$ or $Z$, track its conditional mean: average
only over the color pairs that have not yet been processed. A coordinate
$z_{v,c}$ or $d_{v,c}$ can change only when the pair containing $c$ is processed. A
coordinate $Z_{v,\{c,d\}}$ can change only when the pair containing $c$
or the pair containing $d$ is processed. Thus it changes at most twice.
If $c,d$ are paired together, their product at every vertex is unchanged.

Whenever such a conditional mean changes, its new value is a constant
plus a sum of nonnegative functions of the independent path coins at
that step. For a product involving a color from an unprocessed pair,
replace that missing indicator by its mean, a number in $[0,1]$. If its
pair was processed earlier, use its fixed indicator instead. A path coin
changes indicators at at most two vertices. Its contribution is therefore
at most
\[
 \frac{2L}{\alpha n}\quad\text{to a coordinate of }z,
 \qquad
 \frac{2L}{\alpha^2n^2}\quad\text{to a coordinate of }d\text{ or }Z.
\]
To control lower as well as upper errors, also use the sum obtained by
replacing each summand by its maximum possible value minus that summand.
All of these new summands are nonnegative.

Apply Proposition~\ref{prop:HS}, multiplying the $z$ sums by a sufficiently
small fixed constant times $n$, and the $d,Z$ sums by a sufficiently small
fixed constant times $n^2$. Every summand then lies in $[0,1]$. The
expected size of each scaled sum is $O(n)$, since before scaling
\[
 z_{v,c}\le 1/\alpha,\qquad
 d_{v,c},Z_{v,\{c,d\}}\le1/(\alpha^2n).
\]
The same estimates apply to the sums used for lower errors.

The total contribution of any one value of a path coin to all scaled
sums is $O(n^2)$. Here is the count. Each of the coin's at most two ends
$u$ contributes to target vertices $v$ with total weight at most
$W_Q(u)\le n$. For the $z,d$ sums there are only the two current colors.
For the $Z$ sums one color must be current, and the other has at most
$K$ possibilities. The factors $1/f_u$ and $1/f_u^2$ cancel the respective
scalings up to fixed constants. Thus the $z,d$ contribution is $O(n)$ and
the $Z$ contribution is $O(Kn)=O(n^2)$. The sums used for lower errors
only double these bounds.

Proposition~\ref{prop:HS} now bounds all scaled errors by
$O(\sqrt{n\log n}+\log n)$. Each coordinate is affected at most twice
as the color pairs are processed. We can therefore complete one whole
stage with coordinate errors
\[
 O\bigl(n^{-1/2}\sqrt{\log n}\bigr)\quad\text{for }z,
 \qquad
 O\bigl(n^{-3/2}\sqrt{\log n}\bigr)\quad\text{for }d,Z.
\]
All constants are independent of the number of vertices. Notice that
no ordinary-degree bound was used.

Apply this procedure for all $s=O(\log n)$ stages. Averaging does not
increase the largest coordinate error. Put
\[
 E_n=\sqrt n(\log n)^{3/2},\qquad
 D_v=\sum_{e=uv}\frac{q_e}{f_u}.
\]
The fixed sums of the lists are
\[
 \sum_cz_{v,c}=W_Q(v),\qquad
 \sum_cd_{v,c}=D_v,\qquad
 \sum_{c<d}Z_{v,\{c,d\}}=\frac{W_Q(v)-D_v}{2}.
\]
The mixing property of the chosen pairings, together with the accumulated
coordinate errors, consequently gives
\begin{align}
 z_{v,c}&=\frac{W_Q(v)}K+O(E_n/n),\label{eq:moment-first}\\
 d_{v,c}&=\frac{D_v}K+O(E_n/n^2),\label{eq:moment-diagonal}\\
 Z_{v,\{c,d\}}
 &=\frac{W_Q(v)-D_v}{K(K-1)}+O(E_n/n^2).
 \label{eq:moment-second}
\end{align}
The errors from imperfect averaging are smaller than the displayed errors:
the initial norms are $O(\sqrt n)$, $O(n^{-1/2})$, and $O(1)$,
respectively.

\medskip\noindent\textbf{Turning these sums into fractions.}
For an edge $e=uv$, write
\[
 \lambda_e=\frac{K|F_u\cap F_v|}{f_uf_v}.
\]
We do not require every $\lambda_e$ to be close to one. Instead we bound
the weighted sum of their squared errors at each vertex.
By \eqref{eq:moment-first},
\[
 \sum_{e\ni v}q_e\lambda_e
 =\frac K{f_v}\sum_{c\in F_v}z_{v,c}
 =W_Q(v)+O(E_n).
\]
For the second moment, expand the square of $|F_u\cap F_v|$ and use
\eqref{eq:moment-diagonal} and \eqref{eq:moment-second}. We get
\begin{align*}
 \sum_{e\ni v}q_e\lambda_e^2
 &=\frac{K^2}{f_v^2}\left(
       \sum_{c\in F_v}d_{v,c}
       +2\sum_{\{c,d\}\subseteq F_v}Z_{v,\{c,d\}}\right)\\
 &=\frac{K D_v}{f_v}
   +\frac{K(f_v-1)}{f_v(K-1)}\bigl(W_Q(v)-D_v\bigr)+O(E_n)\\
 &\le W_Q(v)+O(E_n).
\end{align*}
For the last inequality, the first term is $O(1)$ because
$D_v\le1/\alpha$ and $K/f_v\le C/\alpha$; the coefficient of
$W_Q(v)-D_v$ is at most one because $f_v\le K$.
Combining the first and
second moments gives
\[
 \sum_{e\ni v}q_e(\lambda_e-1)^2=O(E_n).
\]
The Cauchy--Schwarz inequality therefore gives
\begin{equation}
 \sum_{e\ni v}q_e|\lambda_e-1|
 \le\sqrt{W_Q(v)\sum_{e\ni v}q_e(\lambda_e-1)^2}
 =O\bigl(n^{3/4}(\log n)^{3/4}\bigr).
 \label{eq:moment-total-error}
\end{equation}

Finally set
\[
 x_{e,c}=\frac{K}{f_uf_v}\,
       \frac{\one_{\{c\in F_u\cap F_v\}}}{\max\{1,\lambda_e\}}
       \qquad(e=uv).
\]
These fractions have the required support and individual upper bound,
and their sum on an edge is $\min\{1,\lambda_e\}$. Equation
\eqref{eq:moment-total-error} proves \eqref{eq:moment-loss}. For a color
missing at $v$, division by $\max\{1,\lambda_e\}$ only decreases the load,
so \eqref{eq:moment-first} gives
\[
 \sum_{e\ni v}q_ex_{e,c}
 \le\frac K{f_v}z_{v,c}
 \le\frac{W_Q(v)}{f_v}+O(E_n/n).
\]
For a color present at $v$ the load is zero. This proves
\eqref{eq:moment-load} and completes the proof.
\end{proof}

\subsection{Two elementary coloring facts}
The first fact turns small fractional loads into actual feasible sets.

\begin{lemma}[Feasible choices from a half-load bound]\label{lem:half-load}
Suppose $0<w_i\le1/2$ for a finite set of items $I$. If
\[
 0\le x_i\le1\quad(i\in I),\qquad \sum_iw_ix_i\le1/2,
\]
then there is a random subset $S\subseteq I$ such that
$\sum_{i\in S}w_i\le1$ always and $\Prb(i\in S)=x_i$ for every $i$.
Equivalently, $x$ is a weighted average of the zero--one vectors of sets
of total weight at most one.
\end{lemma}
\begin{proof}
We turn the entries of $x$ into zeros and ones by random choices that
preserve the expected value of each entry. Call an entry \emph{fractional}
if it is strictly between zero and one.

Suppose two entries $x_i,x_j$ are fractional. Replace them temporarily by
\[
 x_i(t)=x_i+t w_j,\qquad x_j(t)=x_j-t w_i.
\]
Their combined weight stays fixed. The values of $t$ for which both
entries lie in $[0,1]$ form an interval $[-a,b]$, with $a,b>0$. At each
end of this interval, at least one of the two entries is zero or one.
Choose $t=b$ with probability $a/(a+b)$ and $t=-a$ with probability
$b/(a+b)$. Then $\E t=0$, so the expected value of each entry after this
choice is its value before the choice. All other entries stay fixed.

Repeat while at least two fractional entries remain. Each step reduces their
number, preserves total weight, and preserves all entry values in
expectation. Thus the process stops after at most $|I|$ steps, with total
weight still at most $1/2$ and with at most one fractional entry. If that
entry has value $q$, round it to one with probability $q$ and to zero
otherwise. This last choice also preserves its expected value. It adds
at most $\max_i w_i\le1/2$ to the total weight, so the final zero--one
vector always describes a set of weight at most one. Repeatedly taking
conditional expectations shows that the probability its $i$th entry is
one is the original value $x_i$, as required.
\end{proof}


The exact general-multigraph proof uses the following constructive form
of Shannon's ordinary edge-coloring theorem.  The later bin-packing proof
also uses it to color its large groups.

\begin{lemma}[A proper coloring with few colors]\label{lem:three-per-two}
Every multigraph without loops and of maximum degree $d$ has a proper
edge coloring with at most $3\lceil d/2\rceil$ colors.
\end{lemma}
\begin{proof}
There is nothing to do if $d=0$. Put $r=\lceil d/2\rceil$.
Take two disjoint copies of the graph. For every original vertex $v$,
join its two copies by $2r-d(v)$ new parallel edges. Every vertex of
the enlarged graph now has degree $2r$.

Orient an Euler tour in each connected component. Each vertex has
exactly $r$ incoming and $r$ outgoing edges. Make two new vertices
$v_{\rm out},v_{\rm in}$ for each vertex $v$, and replace each directed
edge $u\to v$ by the edge $u_{\rm out}v_{\rm in}$. This graph is
bipartite and has maximum degree $r$. K\H{o}nig's theorem colors its
edges properly with $r$ colors.

Transfer this coloring to the enlarged graph. At a vertex, each color
now appears on at most one incoming edge and at most one outgoing edge.
The edges of that color therefore form paths and cycles, which can be
properly colored with three colors. Use a separate set of three colors
for each of the $r$ colors, and then restrict to the original graph.
\end{proof}

\subsection{Completing the coloring of the groups}
The next argument works whether or not the graph is bipartite. The
coloring of $H$ is the only input that will differ between the two cases.

\begin{lemma}[Completing the coloring of parallel-edge groups]\label{lem:finish-packets}
Fix $\tau>0$ and $C>0$. Suppose the grouped graph above has a proper
coloring of $H$ with an even number $K$ of colors, where
\[
 (2+\tau)n\le K\le Cn.
\]
Then the original graph has a feasible coloring with $K+o(n)$ colors.
The error depends only on $\tau,C$, not on the number of vertices.
\end{lemma}
\begin{proof}
For every vertex,
\[
 d_H(v)\le2W_H(v)\le2(n-W_R(v)).
\]
Thus its number of missing colors satisfies
\[
 f_v=K-d_H(v)\ge\tau n+2W_R(v).
\]
Apply Lemma~\ref{lem:weighted-moment-mix} to the simple graph $R$, giving
an edge of weight $w_e$ charge $q_e=2w_e$. Its weighted degree in these
charges is at most $2n$, and the charge between any pair is at most one.
The scale parameter in that lemma can therefore be $2n$, with
$\alpha=\tau/2$ and color-bound constant $C$.
It supplies fractions supported on common missing colors, each
$O_{\tau,C}(1/n)$, whose sum on each edge is at most one. At every vertex,
\[
 \sum_{e\ni v}w_e\left(1-\sum_cx_{e,c}\right)=o(n).
\]
For each color, their local weighted load is at most
\[
 \sum_{e\ni v}w_ex_{e,c}
 \le\frac{W_R(v)}{f_v}+o(1)
 \le\frac12-\frac{\tau}{2C}+o(1)<\frac12
\]
for sufficiently large $n$. Here we used $f_v\le K\le Cn$.
Lemma~\ref{lem:half-load} makes each local vector an average of feasible
edge sets.

Fix $\delta>0$. Theorem~\ref{thm:round}, applied to the simple graph
$R$, gives a feasible partial coloring using the same $K$ colors and
leaving weight at most $\delta n+o(n)$ at every vertex. Each color used
on an edge of $R$ is missing at both endpoints in $H$, so the colorings
fit together. Lemma~\ref{lem:general-cleanup} colors the uncolored
edges with at most $4\delta n+o(n)$ new colors. Expand the groups back
into their original edges; this preserves every color load. Taking
$\delta$ arbitrarily small proves the lemma.
\end{proof}

\section{The bin-packing upper bound in detail}\label{sec:bp-details}
Proposition~\ref{prop:lower-four-thirds} proves the lower bound $4/3$.
This appendix proves the upper bound
$\frac{81}{50}\bp(G)+o(\bp(G))$ for all bipartite multigraphs.

\subsection{An upper bound for all bipartite multigraphs}
\label{sec:bp-unrestricted}



The proof starts with two parts. First, assign each edge the charge
$a(w)=1/\lfloor1/w\rfloor$. We show that the largest total charge at
a vertex, plus a smaller-order error, pays for a coloring. Second,
color a balanced tenth of the edges in $(1/2,3/5]$ and $(1/3,2/5]$
together. This costs about $\bp(G)/10$ colors and reduces the charge
of the remaining graph to at most $(371/240)\bp(G)+O(1)$.
The sum is $79/48$.  The final subsection replaces this one selection
by the three choices in Table~\ref{tab:three-compatible-modes}; their
cost plus the remaining charge is $81/50$.

We begin with the charge bound. Its proof uses
Lemma~\ref{lem:weighted-moment-mix} and the following two tools.
The \emph{multiplicity} of a graph is the largest number of edges
joining any one pair of vertices.

\begin{lemma}[Rounding with bounded multiplicity]\label{lem:bp-mult-round}
Fix $L\in\mathbb N$ and $C_0,C_1,\delta>0$. Let $Q$ be a weighted multigraph without loops of multiplicity at most $L$, with $W_Q(v)\le n$,
and let $\mathcal C$ contain at most $C_1n$ colors. Suppose
\[
 0\le x_{e,c}\le C_0/n,\qquad \sum_cx_{e,c}\le1,
 \qquad (x_{e,c}:e\ni v)\in\mathcal K_v
 \quad\text{for every }v,c.
\]
Then, for sufficiently large $n$, a feasible partial coloring uses
color $c$ on $e$ only when $x_{e,c}>0$, and leaves colored weight at
least
\[
 \sum_{e\ni v}w_e\sum_cx_{e,c}
 -\delta n-O_{L,C_0,C_1,\delta}(n^{3/4})
\]
at every vertex. The threshold and error constant are independent of
the number of vertices.
\end{lemma}
\begin{proof}
Use the proof of Theorem~\ref{thm:round}, keeping every parallel edge
as a separate edge. We give all changes where simplicity was used.

In the division into vertex groups, one source vertex contributes at
most $L$ to a row, while its sum over rows is still at most $n$.
Divide the rows by $L$ before applying Proposition~\ref{prop:HS} and
then multiply back. The error is still
$O_L(\sqrt{n\log n})$.

The list at a source $u$ now contributes the sum over its at most $L$
edges to a given target $v$. Its ordinary-slot contribution is
$O(L/n)$; its contribution to an error row is
$O(\sum_{e: e=uv}w_e)$. Increase the fixed scaling constant $B_0$ in
\eqref{eq:APtable} by a factor of $L$. The total contributions over
all target rows are still $O_L(n^2)$ for the scaled load rows and
$O_L(n)$ for the error rows, because sums over incident edges still
use the bound $W_Q(u)\le n$. Thus the Harris--Srinivasan bounds,
the induction on the vertex groups, and the $O(n^{3/4})$ total-error
bound hold with constants depending on $L$.

The variance calculation is for one edge at a time, using independent
colors, and is unchanged. The deletion of edges inside vertex groups
also remains valid: there are no loops, and all edges joining the same
two vertices are kept or deleted together.

In the construction of three-element sets, job--slot pairs have the
same bounds as before. A pair of ordinary slots can now belong to
options for at most $L$ original edges, giving $O(LN/n)$ copies.
A pair of extra slots can likewise belong to at most $L$ original
edges, giving $O(LNn^{-4/3})$ copies. Consequently the largest pair
count is still $O_L(Nn^{-2/3})$. With $N=\lceil n\rceil^6$, this is
$o(D^{1-\zeta})$ for every fixed $0<\zeta<1/9$, exactly as required
in the original proof. The job degrees, the slot degrees, the rounding
of small edges, and the final weight estimates do not use simplicity.
This proves the stated extension.
\end{proof}

\begin{lemma}[Choosing from finitely many size classes]
\label{lem:bp-classchoice}
Partition a finite set of items into classes $I_2,\ldots,I_T$, with
every item in $I_k$ having weight at most $1/k$. Let $x_i\ge0$, and
put
\[
 p_k=\max\left\{\frac1k\sum_{i\in I_k}x_i,
                       \max_{i\in I_k}x_i\right\},
\]
where the maximum over an empty class is zero. If $\sum_{k=2}^Tp_k\le1$,
there is a random feasible set with inclusion probability $x_i$ for
every item $i$.
\end{lemma}
\begin{proof}
Choose class $k$ with probability $p_k$, and choose the empty set with
the remaining probability. When $p_k>0$, the conditional numbers
$z_i=x_i/p_k$ satisfy $0\le z_i\le1$ and $\sum_{i\in I_k}z_i\le k$.
They are the inclusion probabilities of a random subset of at most
$k$ items. Here is an elementary way to see this last fact. While two
entries are fractional, increase one and decrease the other by the
same amount, randomly choosing one of the two endpoints of the
allowable interval so that all expectations stay fixed. Each step
makes an entry zero or one and preserves the sum. At most one fractional
entry remains at the end; round it randomly with its own value as the
probability of rounding up. The resulting sum is at most
$\lceil\sum z_i\rceil\le k$. Taking conditional expectations shows
that the inclusion probabilities are the original $z_i$.

All chosen items have weight at most $1/k$, so the chosen set is
feasible. Multiplication by the probability $p_k$ gives unconditional
inclusion probability $x_i$.
\end{proof}

\Needspace{17\baselineskip}
\begin{lemma}[Coloring from the total charge]
\label{lem:harmonic-load}
For an edge weight $0<w\le1$, put
\[
 a(w)=\frac1{\lfloor1/w\rfloor},\qquad
 A(G)=\max_v\sum_{e\ni v}a(w_e).
\]
Every bipartite weighted multigraph with $W_G(v)\le n$ at every vertex
satisfies
\[
 \chi'_w(G)\le A(G)+o(n).
\]
More precisely, for every $\eta>0$ there is $N(\eta)$ such that the
number of colors is at most $A(G)+\eta n$ whenever $n\ge N(\eta)$.
The threshold is independent of the number of vertices, the number of
parallel edges, and the smallest weight.
\end{lemma}
\begin{proof}
The empty graph is immediate. Notice that $w\le a(w)\le2w$, so
$A(G)\le2n$.
Fix $0<\varepsilon<1$. Choose an integer $T$ large enough that
\[
 b=1/T\le\min\{1/4,\varepsilon/8\}.
\]
Both $b$ and $T$ stay fixed while $n$ tends to infinity.

For each endpoint pair, group its original edges of weight below $b$
into packets of weight in $[b,2b)$, leaving at most one tiny packet of
weight below $b$. Keep every original edge of weight at least $b$
separately. Give each non-tiny object of weight $w$ charge $a(w)$;
give each tiny packet charge $w/(1-b)$.
Grouping preserves the weighted degree and can later be undone without
changing any color load. Figure~\ref{fig:tiny-packets} shows this step.

\begin{figure}[htbp]
\centering
\begin{tikzpicture}[x=1cm,y=1cm,font=\small]
\node[vertex,label=left:$u$] (u) at (0,0) {};
\node[vertex,label=right:$v$] (v) at (3.8,0) {};
\foreach \yy/\ww/\cc in {1.6/0.40b/inkblue,.8/0.35b/inkblue,0/0.30b/inkblue,-.8/0.20b/warmred,-1.6/0.10b/warmred}{
 \draw[thick,\cc] (u)..controls (1.1,\yy) and (2.7,\yy)..node[pos=.5,fill=white,inner sep=1pt] {$\ww$}(v);
}
\node at (1.9,1.8) {All weights below $b$};
\draw[-{Latex},thick] (4.55,0)--(5.5,0);
\node[vertex,label=left:$u$] (uu) at (6.2,0) {};
\node[vertex,label=right:$v$] (vv) at (10,0) {};
\draw[very thick,inkblue] (uu)..controls (7.3,1.1) and (8.9,1.1)..node[pos=.5,above=3pt] {$1.05b\in[b,2b)$}(vv);
\draw[very thick,warmred,dashed] (uu)..controls (7.3,-1.1) and (8.9,-1.1)..node[pos=.5,below=3pt] {$0.30b<b$}(vv);
\node at (8.1,1.8) {A packet and a tiny remainder};
\end{tikzpicture}
\caption{Collect small parallel edges until their total first reaches
$b$. Every completed packet has weight below $2b$, and at most one
tiny remainder is left between each pair. The original edge weights
may be arbitrarily small; $b$ is a fixed grouping threshold.}
\label{fig:tiny-packets}
\end{figure}

At each vertex, the increase in total charge is at most $4bn$.
Indeed, a completed small-edge packet of weight $w<2b$ has
\[
 a(w)<\frac{w}{1-w}\le(1+4b)w,
\]
where the last inequality uses $b\le1/4$.
A tiny packet also has charge at most $(1+4b)w$.
The original edges forming these packets had total charge at least
their total weight. Thus the increase is at most $4b$ times their
original total weight, which is at most $4bn$.
The charges of the other original edges stay unchanged.

Within each endpoint pair and each non-tiny class
$k=\lfloor1/w\rfloor\in\{1,\ldots,T\}$, combine exactly $k$ objects
at a time. Regard every full group as one edge of a multigraph $H$.
Its weight is at most one, and its charge is exactly one.
Let $R$ consist of the ungrouped non-tiny objects and the tiny packets.
It has multiplicity at most
\[
 L=1+T(T-1)/2.
\]
There is no class-$1$ remainder, so every edge of $R$ has weight at most
$1/2$. Retain its charge
\[
 q_e=\begin{cases}
 a(w_e),&w_e\ge b,\\
 w_e/(1-b),&w_e<b,
 \end{cases}
 \qquad Q_v=\sum_{e\in E_R(v)}q_e.
\]
In particular, $w_e\le q_e\le1/2$, and
\begin{equation}\label{eq:hload-charge}
 d_H(v)+Q_v\le A(G)+4bn\le A(G)+\varepsilon n/2.
\end{equation}

Let $K$ be the least even integer at least $A(G)+\varepsilon n$.
For large $n$, $K\le4n$.
K\H{o}nig's theorem gives a proper $K$-coloring of $H$.
Its missing-color counts satisfy
\[
 f_v=K-d_H(v)\ge Q_v+\varepsilon n/2.
\]
Apply Lemma~\ref{lem:weighted-moment-mix} to $R$ with charges $q_e$,
using size parameter $3n$. Indeed, $Q_v\le A(G)+4bn\le3n$,
$f_v\ge(\varepsilon/6)(3n)$, $K\le2(3n)$, and the total charge
between any pair is at most $L/2$.
The lemma gives fractions supported on the common missing colors,
with $x_{e,c}=O_\varepsilon(1/n)$, $\sum_cx_{e,c}\le1$, and
\[
 \sum_{e\in E_R(v)}q_ex_{e,c}\le\frac{Q_v}{f_v}+o(1),
 \qquad
 \sum_{e\in E_R(v)}q_e\left(1-\sum_cx_{e,c}\right)=o(n).
\]
All constants may depend on $\varepsilon$ and its chosen $T$, but not
on the graph or its smallest weight.
Also $Q_v/f_v\le1-\varepsilon/8$, because $f_v\le K\le4n$.

We spell out why these fractions are averages of feasible sets at
one vertex and one color. For $2\le k\le T$, let $R_k(v)$ be the
non-tiny edges at $v$ with $\lfloor1/w_e\rfloor=k$, and put
\[
 p_k=\max\left\{\frac1k\sum_{e\in R_k(v)}x_{e,c},
                             \max_{e\in R_k(v)}x_{e,c}\right\}.
\]
For tiny edges let
\[
 p_0=\max\left\{\sum_{e\in E_R(v):\,w_e<b}\frac{w_e}{1-b}x_{e,c},
                             \max_{e\in E_R(v):\,w_e<b}x_{e,c}\right\}.
\]
Maxima over empty sets are zero. The sum of these probabilities is at
most the local charge load plus $T\max_e x_{e,c}$, hence is at most one
for sufficiently large $n$.
Choose case $k$ with probability $p_k$, case zero with probability
$p_0$, and the empty set with the remaining probability.
In case $k$, Lemma~\ref{lem:bp-classchoice} supplies a set of at most $k$
edges with conditional inclusion probabilities $x_{e,c}/p_k$.
In case zero, these conditional probabilities have weighted sum at most
$1-b$ and every item has weight below $b$. While two entries are fractional, move them in opposite weighted
directions so that the total weight stays fixed. Randomly choose an
endpoint of the allowable interval, with probabilities that preserve
both expectations; at least one entry becomes zero or one. Repeat
until at most one fractional entry remains, then round it up with
probability equal to its value. Only this last step can increase the
total weight, and its increase is less than $b$. Thus the selected
set has weight at most one, with the required inclusion probabilities. This proves the local condition
required by Lemma~\ref{lem:bp-mult-round}.

Because $q_e\ge w_e$, the total original weight not proposed at each
vertex is $o(n)$. Apply Lemma~\ref{lem:bp-mult-round} with any fixed
error parameter $\delta>0$. It colors all but weight at most
$\delta n+o(n)$ at every vertex. Its colors are missing from $H$ at both
ends, so the two colorings fit together.
All remaining edges have weight at most $1/2$.
Lemma~\ref{lem:cleanup} colors them with at most
$2\delta n+o(n)+1$ additional colors.
Expanding the groups and packets gives
\[
 \chi'_w(G)\le A(G)+(\varepsilon+2\delta)n+o(n).
\]
Choose $\varepsilon,\delta$ sufficiently small compared with the given
$\eta$, then choose $T$, and finally take $n$ sufficiently large.
Every choice depends only on $\eta$. This proves the stated uniform bound.
\end{proof}


\subsection{A short five-thirds proof}\label{sec:short-five-thirds}

\begin{theorem}[A short proof of the five-thirds bound]
\label{thm:bp-five-thirds-short}
Every bipartite weighted multigraph satisfies
\[
 \chi'_w(G)\le\frac53\bp(G)+o(\bp(G)).
\]
The error is independent of the number of vertices, the multiplicities,
and the smallest edge weight. In fact,
$\chi'_w(G)\le(5/3)\bp(G)$ whenever $\bp(G)$ is sufficiently large,
with a threshold independent of these graph parameters.
\end{theorem}

The proof uses the harmonic-charge coloring lemma
(Lemma~\ref{lem:harmonic-load}): if $W_J(v)\le n$ at every vertex, then
a bipartite multigraph $J$ can be colored with
\[
 \max_v\sum_{e\ni v}a(w_e)+o(n)
 \quad\text{colors},\qquad a(w)=\frac1{\lfloor1/w\rfloor}.
\]
That lemma is proved in Appendix~\ref{sec:bp-details}, including its
uniform error bound. The short argument below is the additional step
that turns this charge bound into a $5/3$ bound in terms of $\bp(G)$.

\begin{proof}
Put $n=\bp(G)$ and consider the two ranges
\[
 A=(1/2,5/8],\qquad B=(1/3,3/8].
\]
An $A$-edge and a $B$-edge fit together, as do two $B$-edges.
At each vertex, $n$ local bins contain at most $n$ edges from $A$
and at most $2n$ edges from $A\cup B$.

\textit{Set aside a balanced sixth.}
Choose one sixth of each range, to within one edge at every vertex.
Here is a direct way to do this. In the graph of one range, divide the
incidences at each vertex into groups of six and one possible smaller
group. Split the vertex into one copy per group. The resulting graph
is bipartite and has maximum degree six. Properly color it with six
colors by K\H{o}nig's theorem, and keep any one color. Every full group
contributes one selected edge; the smaller group contributes zero or
one. Apply this construction separately to $A$ and $B$.

Let $F$ be the selected union and put $K=\lceil n/6\rceil+1$.
At every vertex, $F$ has at most $K$ incident $A$-edges and at most
$2K$ incident edges in total. Split each vertex into two copies,
put all its $A$-edges at the first copy, and distribute its $B$-edges
so that each copy has degree at most $K$. Properly $K$-color the graph
of copies. After joining the copies back together, each color has at
most two edges at a vertex, with at most one from $A$. Every such pair
fits, so $F$ needs at most $n/6+O(1)$ colors.

\begin{figure}[htbp]
\centering
\begin{tikzpicture}[x=1cm,y=1cm,font=\small]
\fill[softblue] (0,.8) rectangle (5,1.4);
\fill[softgold] (5,.8) rectangle (8,1.4);
\draw[thick] (0,.8) rectangle (8,1.4);
\draw (5,.8)--(5,1.4);
\node at (2.5,1.1) {$A$: at most $5/8$};
\node at (6.5,1.1) {$B$: at most $3/8$};
\fill[softgold] (0,-.35) rectangle (6,.25);
\fill[black!5] (6,-.35) rectangle (8,.25);
\draw[thick] (0,-.35) rectangle (8,.25);
\draw (3,-.35)--(3,.25);
\draw (6,-.35)--(6,.25);
\node at (1.5,-.05) {$B$: at most $3/8$};
\node at (4.5,-.05) {$B$: at most $3/8$};
\node at (7,-.05) {unused};
\node at (4,1.85) {The possible pairs in a selected color};
\end{tikzpicture}
\caption{Splitting each vertex into two copies allows these pairs.
Both have total weight at most one; singleton and empty colors are
also allowed.}
\label{fig:short-five-thirds}
\end{figure}

\textit{Charge the remaining edges.}
Define
\[
 q(w)=
 \begin{cases}
  \frac56 a(w),&w\in A\cup B,\\
  a(w),&\text{otherwise}.
 \end{cases}
\]
The key observation is just
\begin{equation}\label{eq:short-five-thirds-light}
 0<w\le\frac12\quad\Longrightarrow\quad q(w)\le\frac43w.
\end{equation}
For $w\le1/3$, put $k=\lfloor1/w\rfloor\ge3$; then
$a(w)=1/k\le(4/3)w$. For $1/3<w\le3/8$, the discounted charge is
$5/12\le(4/3)w$. For $3/8<w\le1/2$, it is $1/2\le(4/3)w$.

Every original unit bin now has charge at most $3/2$. If it has no
edge heavier than $1/2$, equation~\eqref{eq:short-five-thirds-light}
gives charge at most $4/3$. Otherwise let $h$ be its unique heavier
edge. The two possibilities give
\[
 \begin{array}{ll}
 h\le5/8:&\displaystyle
    q(h)+\sum_{w\ne h}q(w)
    \le\frac56+\frac43(1-h)\le\frac32,\\[5pt]
 h>5/8:&\displaystyle
    q(h)+\sum_{w\ne h}q(w)
    \le1+\frac43(1-h)\le\frac32.
 \end{array}
\]
Here the sum means the other items in this bin, with multiplicities.

Let $R=G-F$. Selecting a sixth of $A$ removes a sixth of its charge,
to within one at each vertex. Selecting a sixth of $B$ does the same,
to within $1/2$. Consequently,
\[
 \sum_{e\in E(R),\ e\ni v}a(w_e)
 \le\sum_{e\ni v}q(w_e)+\frac32
 \le\frac32n+\frac32.
\]
The harmonic-charge lemma colors $R$ with $3n/2+o(n)$ colors, since
$W_R(v)\le n$. Use fresh colors for $R$ and combine the two colorings:
\[
 \chi'_w(G)\le\frac n6+\frac{3n}{2}+o(n)
             =\frac53n+o(n).
\]
All constant losses above are independent of the graph, so the uniform
error bound follows from that of the harmonic-charge lemma.
\end{proof}

\begin{proof}[Absorbing the error for sufficiently large $\bp(G)$]
The same charge calculation has a small uniform margin. For $w\le1/2$,
the inequality $q(w)\le(5/4)w$ holds except on two intervals:
\[
 \begin{array}{c|c}
 \text{interval}&q(w)-(5/4)w\text{ is at most}\\\hline
 (1/4,1/3]&1/48\\
 (3/8,1/2]&1/32.
 \end{array}
\]
These follow directly from the charges $1/3$ and $1/2$ on those
intervals. Outside them, use $a(w)\le(5/4)w$ for $w\le1/4$ and
$q(w)=5/12\le(5/4)w$ on $B$.

A bin containing a heavy edge has remaining weight less than $1/2$,
so it contains at most one exceptional small item. If $h\le5/8$,
its charge is at most
\[
 \frac56+\frac54\cdot\frac12+\frac1{32}=\frac{143}{96}.
\]
If $h>5/8$, the remaining weight is less than $3/8$, so the second
exception cannot occur; the charge is at most
\[
 1+\frac54\cdot\frac38+\frac1{48}=\frac{143}{96}.
\]
Bins without a heavy edge already have charge at most $4/3$.
Thus the same coloring uses at most
\[
 \left(\frac16+\frac{143}{96}\right)n+o(n)
 =\left(\frac53-\frac1{96}\right)n+o(n).
\]
The uniform $o(n)$ term is at most $n/96$ once $n$ is sufficiently
large, proving the final assertion.
\end{proof}


\subsection{Using two compatible weight ranges}
\label{sec:compatible-ranges}

We can improve the bound by first coloring a small part of the graph in
which two edges of different sizes can share a color. The remaining graph
then has a smaller total charge. We need two elementary ways to split a
bipartite multigraph.

\begin{lemma}[Choosing a balanced fraction of the edges]\label{lem:balanced-fraction}
Let $J$ be a bipartite multigraph and let $k$ be a positive integer. There
is a subgraph $F$ such that every vertex $v$ satisfies
\[
 \left\lfloor\frac{d_J(v)}k\right\rfloor
 \le d_F(v)\le
 \left\lceil\frac{d_J(v)}k\right\rceil.
\]
Here degrees count parallel edges separately.
\end{lemma}
\begin{proof}
At each vertex, divide its incident edges into groups of $k$, with at most
one smaller group left over. Replace the vertex by one copy for each
nonempty group, and attach each edge to the copy containing its incidence.
Make these choices separately at the two ends of every edge. The resulting
graph is still bipartite, and its maximum degree is at most $k$.

By K\H{o}nig's edge-coloring theorem, properly color the new graph with
$k$ colors. Every copy of degree $k$ sees each color exactly once. A
smaller copy sees each color at most once. Take the edges of any one
color and return them to the original graph. Every original vertex then
has one chosen edge from each full group and either zero or one from its
smaller group. These are exactly the required degree bounds.
\end{proof}

Define the two weight ranges
\[
 \mathcal A=(1/2,3/5],\qquad
 \mathcal B=(1/3,2/5].
\]
An edge whose weight lies in $\mathcal A$ can share a color with any edge
whose weight lies in $\mathcal B$, since $3/5+2/5=1$.

\begin{lemma}[An exact coloring for two weight ranges]\label{lem:compatible-ranges}
Suppose every edge of a bipartite multigraph $J$ has weight in
$\mathcal A\cup\mathcal B$. Let $d_{\mathcal A}(v)$ and
$d_{\mathcal B}(v)$ count its incident edges in the two ranges. Then
\[
 \chi'_w(J)=\bp(J)=
 \max_v\max\left\{
   d_{\mathcal A}(v),\,
   \left\lceil\frac{d_{\mathcal A}(v)+d_{\mathcal B}(v)}2\right\rceil
 \right\}.
\]
\end{lemma}
\begin{proof}
Call the displayed maximum $K$. A feasible color contains at most one
$\mathcal A$-edge at any vertex, because each such edge weighs more than
$1/2$. It contains at most two edges altogether, because every edge
weighs more than $1/3$. Thus $K$ is a lower bound on $\bp(J)$.

For the upper bound, replace each vertex $v$ by two copies $v_1,v_2$.
Attach all its $\mathcal A$-edges to $v_1$. Distribute its
$\mathcal B$-edges between the two copies so that both degrees are at
most $K$. This is possible because
\[
 d_{\mathcal A}(v)\le K,\qquad
 d_{\mathcal A}(v)+d_{\mathcal B}(v)\le2K:
\]
put as many $\mathcal B$-edges as possible at $v_1$ without exceeding
$K$, and put the rest at $v_2$. Again, make the choices independently at
the two ends of an edge.

The graph of vertex copies is bipartite and has maximum degree at most
$K$, so properly color it with $K$ colors. When we join the two copies
of each vertex back together, one color contains at most two incident
edges, of which at most one lies in $\mathcal A$. Two
$\mathcal B$-edges have total weight at most $4/5$, and an
$\mathcal A$-edge together with a $\mathcal B$-edge has total weight
at most one. Thus the coloring is feasible. This proves
$\chi'_w(J)\le K$ and hence both equalities.
\end{proof}

Figure~\ref{fig:compatible-copies} illustrates the construction.

\begin{figure}[htbp]
\centering
\begin{tikzpicture}[x=1cm,y=1cm,font=\small]
\node[vertex,label=right:$v$] (v) at (0,0) {};
\foreach \yy/\cc [count=\i] in {1.5/inkblue,.9/inkblue,.3/warmred,-.3/warmred,-.9/warmred,-1.5/warmred}{
 \node[vertex] (l\i) at (-1.9,\yy) {};
 \draw[thick,\cc] (l\i)--(v);
}
\node[align=center] at (-.9,2.35) {Two $\mathcal A$-edges\\and four $\mathcal B$-edges};
\draw[-{Latex},thick] (.8,0)--(2,0);
\node[vertex,label=left:$v_1$] (v1) at (3.1,1) {};
\node[vertex,label=left:$v_2$] (v2) at (3.1,-1) {};
\foreach \yy/\cc [count=\i] in {1.65/inkblue,1/inkblue,.35/warmred}{
 \node[vertex] (a\i) at (5.1,\yy) {};
 \draw[thick,\cc] (v1)--(a\i);
}
\foreach \yy [count=\i] in {-.35,-1,-1.65}{
 \node[vertex] (b\i) at (5.1,\yy) {};
 \draw[thick,warmred] (v2)--(b\i);
}
\node[align=center] at (4.1,2.35) {Two copies,\\each of degree three};
\node[inkblue] at (-.8,-2.05) {$\mathcal A=(1/2,3/5]$};
\node[warmred] at (3.6,-2.05) {$\mathcal B=(1/3,2/5]$};
\end{tikzpicture}
\caption{Put all $\mathcal A$-edges at the first copy, then distribute the
$\mathcal B$-edges so that neither copy has degree above $K$; here $K=3$.
A proper coloring of the graph of copies gives at most two edges of one
color at the original vertex, with at most one in $\mathcal A$.
The blue and red lines show weight ranges, not the chosen colors.}
\label{fig:compatible-copies}
\end{figure}

\subsection{The reduced charge of one bin}
\label{sec:discounted-bin}

Recall the harmonic charge
\[
 a(w)=\frac1{\lfloor1/w\rfloor}.
\]
We reduce this charge by one tenth on the two compatible ranges. Define
\[
 q(w)=
 \begin{cases}
   \dfrac9{10}a(w),&w\in\mathcal A\cup\mathcal B,\\[3pt]
   a(w),&\text{otherwise}.
 \end{cases}
\]

\begin{lemma}[The reduced charge of one bin]\label{lem:discounted-charge}
Every finite collection of positive item weights with sum at most one has total
$q$-charge at most $371/240$.
\end{lemma}
\begin{proof}
We repeatedly use the following elementary inequality: for any positive
integer $j$,
\begin{equation}
 w\le\frac1j\quad\Longrightarrow\quad
 a(w)\le\frac{j+1}{j}\,w.
 \label{eq:harmonic-ratio-small}
\end{equation}
Indeed, if $\ell=\lfloor1/w\rfloor$, then $\ell\ge j$ and
$w>1/(\ell+1)$. Hence
$1/\ell\le(\ell+1)w/\ell\le(j+1)w/j$.
Also $q(w)\le a(w)$ for every weight.

There is at most one item heavier than $1/2$. If there is one, call its
weight $h$. The following six cases cover all possibilities.
\begin{enumerate}
\item
\textit{No item is heavier than $1/2$.}
Use \eqref{eq:harmonic-ratio-small} with $j=2$. The total charge is at
most $3/2$.

\item
\textit{$h\in\mathcal A$, and every other item has weight at most $1/3$.}
The item $h$ has charge $9/10$. The remaining items have total weight
less than $1/2$. If all of them weigh at most $1/4$, their total charge
is at most $(5/4)(1/2)=5/8$.
Otherwise exactly one of them weighs more than $1/4$. That item has
harmonic charge $1/3$, and the items left after it have total weight less
than $1/4$. Their charge is at most $5/16$. Thus in either event the
remaining charge is at most
\[
 \frac13+\frac5{16}=\frac{31}{48}.
\]
The total charge is therefore at most
\[
 \frac9{10}+\frac{31}{48}=\frac{371}{240}.
\]

\item
\textit{$h\in\mathcal A$, and another item has weight
$r\in(1/3,2/5]$.}
The two items have charges $9/10$ and $9/20$. All other items together
weigh less than $1/6$, so \eqref{eq:harmonic-ratio-small} with $j=6$
bounds their charge by $7/36$. The total is at most
\[
 \frac9{10}+\frac9{20}+\frac7{36}
 =\frac{139}{90}<\frac{371}{240}.
\]

\item
\textit{$h\in\mathcal A$, and another item has weight $r>2/5$.}
Necessarily $r<1/2$, so its charge is $1/2$. The remaining weight is less
than $1/10$, with charge at most $11/100$. Thus the total is at most
\[
 \frac9{10}+\frac12+\frac{11}{100}
 =\frac{151}{100}<\frac{371}{240}.
\]

\item
\textit{$h>3/5$, and every other item has weight at most $1/3$.}
The item $h$ has charge one. The remaining weight is less than $2/5$,
with charge at most $(4/3)(2/5)=8/15$. Thus the total is at most
\[
 1+\frac8{15}=\frac{23}{15}<\frac{371}{240}.
\]

\item
\textit{$h>3/5$, and another item has weight $r>1/3$.}
We must have $r<2/5$, so this item lies in $\mathcal B$ and has charge
$9/20$. The remaining weight is less than $1/15$, with charge at most
$16/225$. The total is at most
\[
 1+\frac9{20}+\frac{16}{225}
 =\frac{1369}{900}<\frac{371}{240}.
\]
\end{enumerate}
This proves the bound in every case.
\end{proof}

The same six cases will also be used for the general multigraph bound.
Keeping the selected fraction as a variable gives the following form.

\begin{lemma}[Discounting the two compatible ranges]
\label{lem:discounted-charge-parameter}
For $0\le t\le1$, define
\[
 q_t(w)=
 \begin{cases}
  (1-t)a(w),&w\in(1/2,3/5]\cup(1/3,2/5],\\
  a(w),&\text{otherwise}.
 \end{cases}
\]
Every collection of positive weights with sum at most one has total
$q_t$-charge at most
\[
 M(t)=\max\left\{
 \frac32,\ \frac{79}{48}-t,\ h_\infty-\frac32t,\
 \frac{161}{100}-t,\ \frac{23}{15},\
 \frac{707}{450}-\frac12t
 \right\}.
\]
For
\[
 t_*=\frac23\left(h_\infty-\frac85\right)
\]
we have $M(t_*)=8/5$.
\end{lemma}
\begin{proof}
Use the six cases in the proof of Lemma~\ref{lem:discounted-charge}.
If there is no item above $1/2$, the charge is at most $3/2$.
If the heavy item lies in $(1/2,3/5]$ and every other item is at most
$1/3$, the remaining charge is at most $31/48$, giving
$79/48-t$.  If that heavy item is accompanied by an item in
$(1/3,2/5]$, their charges are $(1-t)$ and $(1-t)/2$.  The other
items have total weight less than $1/6$, so
Lemma~\ref{lem:harmonic-tail-six} gives
\[
 (1-t)+\frac{1-t}{2}+h_\infty-\frac32
 =h_\infty-\frac32t.
\]
If instead the second item is above $2/5$, the old $11/100$ tail bound
gives $161/100-t$.

If the heavy item is above $3/5$ and all other items are at most $1/3$,
the charge is at most $23/15$.  In the last case the second item lies
in $(1/3,2/5]$ and the old $16/225$ tail bound gives
\[
 1+\frac{1-t}{2}+\frac{16}{225}
 =\frac{707}{450}-\frac12t.
\]
These are all six displayed terms.

At $t=t_*$, the third term equals $8/5$.  The first, fifth, and sixth
terms are smaller immediately.  Also
$h_\infty\ge1+1/2+1/6+1/42=71/42$, so $t_*\ge19/315>11/240$.
This makes both $79/48-t_*$ and $161/100-t_*$ smaller than $8/5$.
Hence $M(t_*)=8/5$.
\end{proof}

\subsection{Combining the two colorings}
\label{sec:mixed-bound-proof}

\begin{proof}[Proof of Theorem~\ref{thm:bp-two-ranges}]
Let $n=\bp(G)$. Write $d_{\mathcal A}(v)$ and
$d_{\mathcal B}(v)$ for the numbers of incident edges in the two weight
ranges. Each of the $n$ local bins at $v$ contains at most one
$\mathcal A$-edge and at most two edges in
$\mathcal A\cup\mathcal B$. Therefore
\[
 d_{\mathcal A}(v)\le n,\qquad
 d_{\mathcal A}(v)+d_{\mathcal B}(v)\le2n.
\]

Apply Lemma~\ref{lem:balanced-fraction} with $k=10$ separately to
the subgraph of $\mathcal A$-edges and the subgraph of
$\mathcal B$-edges. Let $F_{\mathcal A},F_{\mathcal B}$ be the selected
subgraphs, and let $F$ be their union. At every vertex,
\[
 d_{F_{\mathcal A}}(v)\le n/10+1,\qquad
 d_{F_{\mathcal A}}(v)+d_{F_{\mathcal B}}(v)\le n/5+2.
\]
Lemma~\ref{lem:compatible-ranges} colors $F$ with at most
$\lceil n/10\rceil+1=n/10+O(1)$ colors.

Let $R=G-F$. Every selected $\mathcal A$-edge removes harmonic charge
one, and every selected $\mathcal B$-edge removes harmonic charge
$1/2$. The lower degree bounds in
Lemma~\ref{lem:balanced-fraction} show that the removed charge at
$v$ is at least
\[
 \frac1{10}\left(d_{\mathcal A}(v)
                   +\frac12d_{\mathcal B}(v)\right)-\frac32.
\]
Consequently,
\begin{align*}
 \sum_{\substack{e\ni v\\e\in E(R)}}a(w_e)
 &\le\sum_{e\ni v}a(w_e)
       -\frac1{10}\left(d_{\mathcal A}(v)
                         +\frac12d_{\mathcal B}(v)\right)+\frac32\\
 &=\sum_{e\ni v}q(w_e)+\frac32\\
 &\le\frac{371}{240}n+\frac32.
\end{align*}
The last inequality applies Lemma~\ref{lem:discounted-charge} to each of
the $n$ original local bins at $v$.

The charge bound in Lemma~\ref{lem:harmonic-load} now colors $R$
using at most $(371/240)n+o(n)$ colors. Its treatment of arbitrarily
small weights applies here without change. Use disjoint sets of colors
for $F$ and $R$. We obtain
\[
 \chi'_w(G)
 \le\left(\frac1{10}+\frac{371}{240}\right)n+o(n)
 =\frac{79}{48}n+o(n).
\]
Since $5/3-79/48=1/48$, this also gives
$\chi'_w(G)\le(5/3)\bp(G)$ for all sufficiently large $\bp(G)$.
\end{proof}

\subsection{Three compatible choices}

We now use three small selections at once.  The next lemma is the same
vertex-copy idea as Lemma~\ref{lem:compatible-ranges}, with more than two
copies allowed.

\begin{lemma}[Coloring compatible heavy and small edges]
\label{lem:compatible-r-copies}
Let $r\ge2$.  Suppose the edges of a bipartite multigraph $J$ are split
into sets $H$ and $S$.  Every $H$-edge has weight in $(1/2,h]$, and every
$S$-edge has weight at most $b$, where
\[
 rb\le1,
 \qquad h+(r-1)b\le1.
\]
Then $J$ has a feasible coloring with
\[
 L=\max_v\max\left\{d_H(v),
       \left\lceil\frac{d_H(v)+d_S(v)}r\right\rceil\right\}
\]
colors.
\end{lemma}
\begin{proof}
Replace each vertex by $r$ copies.  Put all its $H$-incidences at the
first copy, and distribute its $S$-incidences so that every copy has
degree at most $L$.  The two inequalities in the definition of $L$ make
this possible.  The copied graph is bipartite with maximum degree at
most $L$, so K\H{o}nig's theorem properly colors it with $L$ colors.

After the copies are joined, a color has at most $r$ incident edges and
at most one of them is in $H$.  Its weight is at most $rb$ if it contains
only $S$-edges, and at most $h+(r-1)b$ otherwise.  Hence the coloring is
feasible.
\end{proof}

\begin{lemma}[Combining finitely many compatible choices]
\label{lem:compatible-modes}
For each $j=1,\ldots,s$, choose integers $2\le r_j\le k_j$, numbers
$h_j,b_j,f_j>0$, and the two ranges
\[
 H_j=(1/2,h_j],\qquad S_j=(1/(k_j+1),b_j].
\]
Assume
\[
 b_j\le\frac1{k_j},\qquad
 h_j+(r_j-1)b_j\le1.
\]
An edge of weight $w$ is proposed for choice $j$ with probability
\[
 p_j(w)=
 \begin{cases}
  f_j,&w\in H_j,\\
  (r_j/k_j)f_j,&w\in S_j,\\
  0,&\text{otherwise}.
 \end{cases}
\]
Suppose $\sum_jp_j(w)\le1$ for every $w$, and put
\[
 q(w)=a(w)\left(1-\sum_jp_j(w)\right).
\]
If every unit bin has $q$-charge at most $Q$, then every bipartite
weighted multigraph satisfies
\[
 \chi'_w(G)\le\left(Q+\sum_jf_j\right)\bp(G)+o(\bp(G)).
\]
For fixed rational choices and rational, explicitly listed input, the
construction is randomized with expected polynomial running time for
each fixed final accuracy.
\end{lemma}
\begin{proof}
Put $n=\bp(G)$.  Give an edge the eligible labels $j$ with probabilities
$p_j(w)$ and give it the remainder label with the probability left over.
For a positive label $j$, use separate endpoint constraints for the
number of chosen $H_j$-edges and the number of chosen $S_j$-edges.  For
the remainder label, use its harmonic charge in the two endpoint
constraints.  Every option uses only its two endpoint constraints, so
Lemma~\ref{lem:two} chooses one label per edge with an additive error of
at most two in every constraint.

One local bin contains at most one $H_j$-edge.  If it contains no such
edge, it contains at most $k_j$ edges of $S_j$.  If it contains an
$H_j$-edge, it contains at most $\lfloor k_j/2\rfloor$ edges of $S_j$.
Thus in either case
\[
 |H_j|+\frac{r_j}{k_j}|S_j|\le r_j.
\]
Consequently, the part with label $j$ has at most $f_jn+2$ heavy edges
and at most $r_jf_jn+4$ edges in total at any vertex.
Lemma~\ref{lem:compatible-r-copies} colors it with $f_jn+O(1)$ colors.

The remainder has incident harmonic charge at most $Qn+2$ at every
vertex.  Lemma~\ref{lem:harmonic-load} colors it with $Qn+o(n)$ fresh
colors.  Summing the fixed number of palettes proves the claim.  The
proofs of Lemmas~\ref{lem:two}, \ref{lem:compatible-r-copies}, and
\ref{lem:harmonic-load} give the algorithm.
\end{proof}

Use the following three choices.

\begin{table}[htbp]
\centering
\begin{tabular}{c|c c|c|c|c}
$j$ & $k_j$ & $r_j$ & $H_j$ & $S_j$ & $f_j$\\ \hline
1&2&2&$(1/2,5/8]$&$(1/3,3/8]$&$1/20$\\
2&2&2&$(1/2,7/12]$&$(1/3,5/12]$&$1/10$\\
3&3&2&$(1/2,17/24]$&$(1/4,7/24]$&$1/12$
\end{tabular}
\caption{The three compatible choices.  Their coloring costs add to
$1/20+1/10+1/12=7/30$.}
\label{tab:three-compatible-modes}
\end{table}

The hypotheses concerning $h_j,b_j,k_j,r_j$ are immediate from the
table.  The largest total proposal probability is $7/30$, on the heavy
ranges, so every edge has a nonnegative remainder probability.

\begin{lemma}[The three-choice charge calculation]
\label{lem:three-mode-charge}
For the choices in Table~\ref{tab:three-compatible-modes}, every unit
bin has $q$-charge at most $104/75$.
\end{lemma}
\begin{proof}
Here is a short exact check.  For weights at most $1/13$,
\[
 a(w)\le\frac{14}{13}w.
\]
Above $1/13$, divide the weight line at the reciprocal points and at
the endpoints in Table~\ref{tab:three-compatible-modes}.  On each open-left,
closed-right cell $(\ell,u]$, the value of $q$ is constant.  Give that
cell the extra value
\[
 p(\ell)=q(w)-\frac{14}{13}\ell
 \qquad(w\in(\ell,u]).
\]
Cells with nonpositive extra value may be discarded.  A cell may also be
discarded if another cell has a no larger lower endpoint and a no smaller
extra value.  The cells left are the following.

\begin{center}
\small
\begin{tabular}{c|c@{\qquad}c|c}
cell & $p(\ell)$ & cell & $p(\ell)$\\ \hline
$(1/2,7/12]$&$89/390$&$(1/7,1/6]$&$1/78$\\
$(7/12,5/8]$&$31/130$&$(1/8,1/7]$&$3/364$\\
$(5/8,17/24]$&$19/78$&$(1/9,1/8]$&$5/936$\\
$(1/3,3/8]$&$103/1560$&$(1/10,1/9]$&$2/585$\\
$(1/4,7/24]$&$16/351$&$(1/11,1/10]$&$3/1430$\\
$(1/5,1/4]$&$9/260$&$(1/12,1/11]$&$1/858$\\
$(1/6,1/5]$&$4/195$&$(1/13,1/12]$&$1/2028$
\end{tabular}
\end{center}

If the chosen cells have lower endpoints $\ell_i$, the total bin charge
is at most
\[
 \frac{14}{13}+\sum_i p(\ell_i),
 \qquad \sum_i\ell_i<1.                         \tag{*}
\]
This follows by charging $14/13$ times the unused capacity and then the
listed extras.  At most one of the first three cells can occur.

For completeness, the remaining arithmetic is shown in the next table.
For each possible heavy cell, list the lower endpoints that maximize the
extra sum in~$(*)$.  Trying the possible counts in decreasing order of
$p(\ell)/\ell$ gives these four rows; all numbers are exact.
\[
\begin{array}{c|c|c}
\text{heavy cell used}&\text{maximizing lower endpoints}
 &\frac{14}{13}+\max\sum p(\ell)\\ \hline
\text{none}&1/3,\ 1/4,\ 1/5,\ 1/5&17659/14040\\
(1/2,7/12]&1/2,\ 1/4,\ 1/5&1945/1404\\
(7/12,5/8]&7/12,\ 1/5,\ 1/5&18/13\\
(5/8,17/24]&5/8,\ 1/3&721/520
\end{array}
\]
To check the four rows without decimals, take the common denominator
$D=360360$ of the lower endpoints.  Index the eleven nonheavy cells in
any order, set $B_{12}(c)=0$, and use the integer-capacity recurrence
\[
 B_i(c)=\max_{0\le m\le\lfloor c/(D\ell_i)\rfloor}
       \{m p(\ell_i)+B_{i+1}(c-mD\ell_i)\}.
\]
For the first row evaluate $B_1(D-1)$.  For each other row, force the
indicated heavy cell once, exclude all three heavy cells from the
recurrence, and evaluate
$p(\ell_H)+B_1(D-1-D\ell_H)$.  Substitution gives exactly the four
fractions in the last column.  The displayed lower endpoints give the
corresponding suprema in this relaxed calculation; actual weights
approach them from above.
Every one is at most $104/75$.  This proves the lemma.
\end{proof}

\begin{proof}[Proof of Theorem~\ref{thm:bp-unrestricted}]
Apply Lemma~\ref{lem:compatible-modes} with
Lemma~\ref{lem:three-mode-charge}.  Since
\[
 \frac7{30}+\frac{104}{75}=\frac{81}{50},
\]
the theorem follows.  The error and running-time statements are those
of Lemma~\ref{lem:compatible-modes}.
\end{proof}


\subsection{Why the first coloring helps}
The first coloring mixes edges that may have different neighbours.
This step is essential to the improvement: grouping only parallel edges
and properly coloring the groups has a larger limiting factor.
To state this precisely, use the constant $h_\infty$ defined in
\eqref{eq:harmonic-constant}.

\begin{proposition}[A limit of grouping parallel edges]\label{prop:packet-method-limit}
Let $r_i$ and $h_\infty$ be as in~\eqref{eq:harmonic-constant}.
Consider a method with the following two requirements. Each packet
contains original edges with the same two endpoints and has total weight
at most one. For each fixed list of edge weights and endpoint pairs,
the number of original edges left outside packets on each pair is bounded
independently of their multiplicities. The method properly colors the
packet graph before finishing any edges outside packets.

For every $c<h_\infty$, there is a fixed finite list of weights and a family of
bipartite multigraphs $G_N$ such that
\[
 \bp(G_N)=\chi'_w(G_N)=N,
\]
but this method uses more than $cN$ colors for all sufficiently large
$N$. Thus the method cannot have an asymptotic coefficient below $h_\infty$.
\end{proposition}
\begin{proof}
Choose $j\ge1$ such that $\sum_{i=1}^{j}1/r_i>c$. The identity
\[
 \frac1{r_i+1}=\frac1{r_i}-\frac1{r_{i+1}}
\]
gives
\[
 \sum_{i=1}^{j}\frac1{r_i+1}=1-\frac1{r_{j+1}}<1.
\]
We may therefore choose fixed weights
\[
 \frac1{r_i+1}<w_i\le\frac1{r_i}
 \quad(1\le i\le j),\qquad
 \sum_{i=1}^{j}w_i\le1.
\]
For example, $w_i=1/(r_i+1)+1/(2jr_{j+1})$ has these properties.

Take a star with $j$ leaves, and replace its $i$th edge by $N$ parallel
edges of weight $w_i$. Using one edge of each type in every color gives
a feasible coloring with $N$ colors. Since $w_1>1/2$, the $N$ edges of
type one need different colors. Hence $\bp(G_N)=\chi'_w(G_N)=N$.

A packet on pair $i$ contains at most $r_i$ original edges, because
$(r_i+1)w_i>1$. Let $B_i$ bound the number of original edges left outside
packets on that pair; $B_i$ is independent of $N$. The packet graph
therefore has at least $(N-B_i)/r_i$ edges on pair $i$. All packets meet
at the centre of the star, so their proper coloring uses at least
\[
 \sum_{i=1}^{j}\frac{N-B_i}{r_i}
 =N\sum_{i=1}^{j}\frac1{r_i}-\sum_{i=1}^{j}\frac{B_i}{r_i}
 >cN
\]
colors for all sufficiently large $N$. This proves the assertion.
\end{proof}

The original graph in this example needs exactly $\bp(G_N)$ colors.
The proposition limits the stated grouping method alone. Our first
coloring avoids this restriction: it colors a number of edges on
each pair that can grow with the multiplicity, and allows edges from
different pairs to share a color.

\section{Why the exact total-weight bound fails}\label{sec:counterexample}
We give a simple bipartite counterexample to the exact total-weight
conjecture. The construction also makes the distinction between the
two conjectures explicit.

\begin{proposition}[A simple bipartite counterexample]\label{prop:counterexample}
There is a simple bipartite weighted graph with
\[
 \Delta_w(G)=39/20<2,\qquad \bp(G)=3,\qquad \chi'_w(G)=4.
\]
\end{proposition}
\begin{proof}
Take the two vertex sets
\[
 \{P,a_1,a_2,a_3,v\},\qquad \{b_1,b_2,b_3,u\}.
\]
The edges and weights are
\[
 \begin{array}{c|c|c}
 \text{edge}&\text{indices}&\text{weight}\\ \hline
 p_i=Pb_i&i\in\{1,2,3\}&13/20\\
 q_{ij}=a_jb_i&i,j\in\{1,2,3\},\ i\ne j&3/5\\
 x_j=a_ju&j\in\{1,2,3\}&9/20\\
 e=vu&&3/5
 \end{array}
\]
See Figure~\ref{fig:counterexample}. The weighted degrees at
$P,u,b_i,a_j,v$, respectively, are
\[
 39/20,\quad39/20,\quad37/20,\quad33/20,\quad3/5.
\]
Thus $\Delta_w(G)=39/20$. Every vertex has a packing into three bins:
all but $u$ have degree at most three, and at $u$ two of the three
$9/20$ edges share a bin while the other two edges get one bin each.
The three edges at $P$ have pairwise sums greater than one, so they
require three bins. Hence $\bp(G)=3$.

\begin{figure}[htbp]
\centering
\begin{tikzpicture}[x=1cm,y=1cm,font=\small]
\node[vertex,fill=softblue,label=left:$P$] (P) at (0,0) {};
\foreach \ii/\yy in {1/1.4,2/0,3/-1.4}{
 \node[vertex,fill=softgold,label=above:$b_{\ii}$] (b\ii) at (2.6,\yy) {};
 \node[vertex,fill=softblue,label=above:$a_{\ii}$] (a\ii) at (6.2,\yy) {};
}
\node[vertex,fill=softgold,label=right:$u$] (u) at (8.8,0) {};
\node[vertex,fill=softblue,label=below:$v$] (v) at (8.8,-2.4) {};
\foreach \ii in {1,2,3}{
 \draw[inkblue,thick] (P)--(b\ii);
 \draw[warmred,thick,dashed] (a\ii)--(u);
 \foreach \jj in {1,2,3}{
  \ifnum\ii=\jj\else
   \draw[black!65,thick] (b\ii)--(a\jj);
  \fi
 }
}
\draw[black!65,thick] (u)--node[right] {$3/5$}(v);
\node[inkblue] at (1.3,2.1) {$p_i:\ 13/20$};
\node at (4.4,2.1) {$q_{ij}:\ 3/5$};
\node[warmred] at (7.5,2.1) {$x_j:\ 9/20$};
\node[align=center] at (4.4,-2.2) {Every middle pair is joined\\except $a_i b_i$.};
\end{tikzpicture}
\caption{A simple bipartite graph with maximum weighted degree below
two that needs four colors. Blue-filled vertices form one side and
yellow-filled vertices the other. Edge styles identify weights, not
an edge coloring. Crossings of drawn edges are not vertices.}
\label{fig:counterexample}
\end{figure}

Suppose three colors suffice. At each of $P,a_1,a_2,a_3$, every pair
of incident edges has total weight greater than one. Thus all three
colors occur once at each of these four vertices. Counting the twelve
edges other than $e$ at their endpoints in this set, each color occurs
exactly four times.

At each of $b_1,b_2,b_3$, the three incident weights also have pairwise
sums greater than one, so all three colors occur once there. These
vertices account for three occurrences of each color. The remaining
edges $x_1,x_2,x_3$ must therefore account for one occurrence of each
color. At $u$, the edge $e$ cannot share a color with any $x_j$,
because $3/5+9/20>1$. All three colors are blocked, a contradiction.

Finally, the ordinary maximum degree is four. K\H{o}nig's theorem
gives a proper four-edge-coloring, which is feasible since all weights
are at most one. Thus $\chi'_w(G)=4$.
\end{proof}

For this graph the least integer greater than $\Delta_w(G)$ is $N=2$.
The total-weight conjecture predicts $2N-1=3$ colors, so it is false.
The Chung--Ross bound uses $\bp(G)=3$ and permits five colors; this
example does not disprove it. The example is also consistent with an
asymptotic upper bound $2\Delta_w(G)+o(\Delta_w(G))$, which concerns
large values of $\Delta_w(G)$.

\section{Implementing the coloring algorithms}\label{sec:algorithms}
We use rational weights written in binary, with every physical edge
listed in the input. The running-time claims below are for each fixed
accuracy; they do not concern binary-compressed lists of multiplicities.
All returned colorings are checked by summing their loads.

\subsection{Producing the local inputs in polynomial time}
\label{sec:local-input-algorithms}
All weights and proposed fractions in this subsection are rational.
The accuracy parameters are fixed. We explain how to produce the local
packings and distributions used by the coloring algorithms.

\paragraph{Approximate local packings.}
For every fixed $\xi>0$, a finite item multiset $I$ can be packed in
polynomial time into
\[
 (1+\xi)\operatorname{bp}(I)+O_\xi(1)
\]
bins. Choose fixed $\nu,\tau>0$ so that
$1+\nu\le1+\xi$ and $(1-\tau)^{-1}\le1+\xi$.
Apply Lemma~\ref{lem:bin-rounding} to the items of weight at least $\tau$.
Their LP value is at most $\operatorname{bp}(I)$, so this gives at most
$(1+\nu)\operatorname{bp}(I)+O_{\nu,\tau}(1)$ bins.
Insert the lighter items greedily. Whenever a new bin is needed,
every previous bin has load greater than $1-\tau$. Hence the final
number is at most
\[
 \max\left\{(1+\nu)\operatorname{bp}(I)+O_{\nu,\tau}(1),
                   \frac{W(I)}{1-\tau}+1\right\}.
\]
Since $W(I)\le\operatorname{bp}(I)$, this proves the claim.
A uniform choice among the resulting bins is an explicit distribution
on actual sets, giving each item probability exactly the reciprocal
of the number of bins. Empty bins may be added to equalize this number
between vertices.

\begin{lemma}[An explicit distribution for a bounded load]
\label{lem:algorithmic-local-load}
Suppose $0<w_i\le b\le r$, $0\le x_i\le1$, and
$\sum_{i=1}^m w_ix_i\le r-b$.
In polynomial time one can express $x$ as an average of at most
$2m+2$ zero--one vectors, each of weight at most $r$.
Similarly, if $\sum_i x_i\le k$ for an integer $k$, one can express
$x$ as an average of at most $m+1$ sets of at most $k$ items.
\end{lemma}
\begin{proof}
Consider the bounded polytope
\[
 P=\{y\in[0,1]^m:\textstyle\sum_i w_iy_i\le r-b\}.
\]
We first describe a polynomial procedure expressing $x$ as an average
of at most $m+1$ vertices of $P$. At the current point, keep all
inequalities that hold with equality; they describe its smallest face.
Find a vertex $z$ of this face by successive linear optimization of
the coordinates. If the current point is not already $z$, extend the
ray from $z$ through that point until another inequality becomes tight.
The current point lies between $z$ and this new boundary point.
Record that convex combination and continue at the boundary point.
The dimension of the current face drops at every step, so there are
at most $m$ steps. The polytope has $2m+1$ inequalities, and all the
linear programs and ray intersections use polynomially many rational
operations with polynomial encoding length.

A vertex of $P$ has at most one coordinate strictly between zero and
one. Otherwise two such coordinates could be moved slightly in opposite
weighted directions while staying in $P$, contradicting that it is a
vertex. For each vertex in the decomposition, replace its possible
fractional coordinate $y_i$ by one with probability $y_i$ and by zero
otherwise. This preserves its mean. Rounding up increases its weight
by at most $b$, so both resulting sets have weight at most $r$.
Each vertex gives at most two sets, proving the first claim.

For the second claim use
$\{y\in[0,1]^m:\sum_i y_i\le k\}$ instead.
The same two-coordinate argument leaves at most one fractional
coordinate at a vertex. If the sum constraint is not tight, that
coordinate can be moved alone; if it is tight, its value is an integer
because $k$ and the other coordinates are integers. Thus every vertex
is already the required zero--one vector.
\end{proof}

\paragraph{The distributions in the harmonic-load proof.}
Fix a vertex and a color, and abbreviate its proposed fractions by
$x_e$. The non-tiny classes satisfy $k=\lfloor1/w_e\rfloor$; their
charges are $q_e=1/k$. Tiny items have $w_e<b$ and charge
$q_e=w_e/(1-b)$. Define
\[
 p_k=\max\left\{\frac1k\sum_{e\in R_k}x_e,\ \max_{e\in R_k}x_e\right\},
 \qquad
 p_0=\max\left\{\sum_{e\text{ tiny}}\frac{w_e x_e}{1-b},
                          \ \max_{e\text{ tiny}}x_e\right\},
\]
with maxima over empty sets equal to zero.
The harmonic-load argument proves $\sum_k p_k+p_0\le1$.
Choose class $k$ with probability $p_k$, the tiny class with probability
$p_0$, and the empty set with the remaining probability.

For $p_k>0$, the vector $(x_e/p_k:e\in R_k)$ has coordinates at most
one and sum at most $k$. The second part of
Lemma~\ref{lem:algorithmic-local-load} supplies the required distribution
on sets of at most $k$ items; their weights are at most $1/k$, so they fit.
For $p_0>0$, the conditional tiny-item vector has weight at most $1-b$.
Use the first part of the lemma with $r=1$. In each case the unconditional
inclusion probability of $e$ is exactly $x_e$. There are only a fixed
number of classes, and the total support size is polynomial.
The same bounded-load construction applies whenever the proof uses
fractions of load at most $r-b$ for items of weight at most $b$.
Thus these local inputs require no exact bin-packing oracle and no
implicit distribution of exponential support.
\subsection{Computing the progress bounds and sampling vertex lists}
\label{sec:computable-progress}

Here we give the implementation details for the progress bounds in
\eqref{eq:computable-progress}. All probabilities, weights, and assigned
fractions used by the algorithm are rational numbers. The number of
colors is $O(n)$, and the number $M$ of profile labels is fixed once
the error parameters have been chosen.

Fix an edge $e=uv$ and a color $c$. For a profile $P$, define
\[
 A_{v,c}(P)=\sum_i a_{v,c}(e;P,i),\qquad
 b_{v,c}(P)=\frac{A_{v,c}(P)}{\lambda'_{v,c,P}}.
\]
The sum includes all ordinary slots, or the reserved space when $e$ is
small. If the list at $v$ has already been fixed, put
\[
 m_{v,c}=b_{v,c}(P_{v,c}),\qquad
 q_{v,c}=b_{v,c}(P_{v,c})^2.
\]
If that list is still unchosen, its original distribution instead gives
\[
 m_{v,c}=\sum_P A_{v,c}(P)=y_{e,c},\qquad
 q_{v,c}=\sum_P\frac{A_{v,c}(P)^2}{\lambda'_{v,c,P}}.
\]
Define $m_{u,c},q_{u,c}$ in the same way. The contribution of this color
to $S_e$ is
$\rho b_{u,c}(P_{u,c})b_{v,c}(P_{v,c})/y_{e,c}$.
Thus its first and second moments are
\[
 a_c=\frac{\rho m_{u,c}m_{v,c}}{y_{e,c}},\qquad
 d_c=\frac{\rho^2q_{u,c}q_{v,c}}{y_{e,c}^2}.
\]
When $y_{e,c}=0$, set $a_c=d_c=0$. The unchosen endpoint lists are
independent, and their choices at different colors are independent.
Fixing any complete endpoint lists preserves the latter independence.
It follows that
\begin{equation}\label{eq:progress-moment-formula}
 U_e^{(k)}
 =\left(\sum_c a_c-\mu_e\right)^2
   +\sum_c(d_c-a_c^2).
\end{equation}
This formula evaluates the conditional second moment using only
$O(M|\mathcal C|)$ operations per edge. In particular, it does not sum
over all possible complete lists.

For rational arithmetic, choose the largest power of two $\kappa$ not
exceeding $n^{-2}$. Then $(2n^2)^{-1}<\kappa\le n^{-2}$, so the number of
bits needed to write $\kappa$ is $O(\log(n+2))$. For a nonnegative rational
number $z$, the value
\[
 \phi(z)=\kappa\left\lceil\frac{\sqrt z}{\kappa}\right\rceil
\]
is computed by finding the least nonnegative integer $b$ with
$b^2\kappa^2\ge z$. Integer comparisons and binary search compute it in
polynomial time. The same method computes the integer threshold
$t=\lceil n^{3/4}\rceil$, using $t^4\ge n^3$. Every progress value and
every test that accepts or rejects a list is therefore rational and
polynomially computable.

The upward rounding does not change the error bound. Conditional on the
lists fixed before a group is chosen, the tower property gives
$\E U_e^{(k)}=U_e^{(k-1)}$. The Cauchy--Schwarz inequality gives
$\E\sqrt{U_e^{(k)}}\le\sqrt{U_e^{(k-1)}}$. Hence
\[
 \E\phi(U_e^{(k)})\le\phi(U_e^{(k-1)})+\kappa,
 \qquad
 \E Q_v^{(k)}\le Q_v^{(k-1)}+W_H(v)\kappa
                      \le Q_v^{(k-1)}+1/n.
\]
All these expectations use the original distributions of the lists being
chosen. At the end, $U_e^{(r)}=(S_e-\mu_e)^2$, so
$Z_v\le Q_v^{(r)}$. Thus the original $O(n^{3/4})$ absolute-error bound
is retained.

At group $V_k$, first compute
\[
 p_k=\max_{u\in V_k}\frac{Q_u^{(k-1)}+1/n}{t},
\]
with value zero for an empty group. To draw a retained list at $u$, draw
its $O(n)$ profile labels independently from the original rational
probabilities and test $Q_u^{(k)}\le t$. Repeat on rejection. Markov's
inequality bounds the rejection probability by $p_k$. In the range where
the proof is applied, $p_k=o(1)$ and in particular $p_k\le1/2$, so the
expected number of attempts is at most two. There is no edge within
$V_k$, so the acceptance test for $u$ depends only on its own list and
on earlier groups. It remains the same throughout this group's
resampling procedure. Consequently this procedure samples exactly the
required independent retained distributions, even though their
normalizing probabilities are not computed.

We finally explain why the use of Harris--Srinivasan does not require
listing every possible complete vertex list. In the proof of their
Theorem~4.7, a fractional product distribution is first found and then
rounded by their Algorithm~2. Here that distribution is already supplied
by the sampling procedure just described, so the first step is
unnecessary. More explicitly, their Theorem~3.8 uses resampling
probabilities $p_{i,j}=\lambda_{i,j}/\sum_j\lambda_{i,j}$. The proof
of their Theorem~4.7 takes $\lambda_{i,j}=(1+\varepsilon)z_{i,j}$,
where $z$ is the supplied distribution, so these resampling
probabilities are exactly $z_{i,j}$. No new distribution has to be
computed. The rounding algorithm only needs to draw independent
values of its variables, find a violated row, and sample a subset of the
assignments currently contributing to that row. The row family here has
polynomial size, and its coefficients are computable by
\eqref{eq:progress-moment-formula} and the displayed formula for $\phi$.

For clarity, the subset sampling can also be implemented directly. If a
violated row currently has coefficients $a_1,\ldots,a_N\in[0,1]$, the
Harris--Srinivasan rule chooses a set $I$ of a specified size $d$ with
probability proportional to $\prod_{i\in I}a_i$. Put
\[
 E(j,\ell)=\sum_{\substack{I\subseteq\{1,\ldots,j\}\\|I|=\ell}}
                  \prod_{i\in I}a_i.
\]
The recurrence
\[
 E(j,\ell)=E(j-1,\ell)+a_jE(j-1,\ell-1),\qquad E(0,0)=1
\]
computes the table in $O(Nd)$ operations, taking impossible entries to
be zero. Starting with $j=N$ and $\ell=d$, include $j$ with probability
$a_jE(j-1,\ell-1)/E(j,\ell)$ and decrease $\ell$ when it is included;
then decrease $j$ and continue. This samples the required distribution.
Only the current assignments are used, not the set of all possible
values of any variable. This is the implementation principle in
Appendix~C of~\cite{HS}.

The proof of Theorem~4.7 of~\cite{HS} bounds the expected number of
resampling steps by a constant times the number of variables. It also
allows supplied upper bounds on the row means. For our load rows use
$m_k/(1-p_k)$; for an error row with unchanged contribution $C$ use
$(Q_v^{(k-1)}-C+1/n)/(1-p_k)$, with the scaling from
\eqref{eq:APtable}. Thus no expectation under a retained distribution is
evaluated. Sampling a list, evaluating rows, and the subset selection
all take expected polynomial time. The entire choice of retained lists
therefore does too, for fixed error parameters.
\subsection*{Finding the mixing operations}
The list of color pairings in Lemma~\ref{lem:weighted-moment-mix}
can also be found explicitly. For each candidate list, form its
single-color and unordered-pair averaging matrices. If $J_d$ denotes
the $d$ by $d$ matrix all of whose entries are $1/d$, accept a list
when
\[
 \left\|A_s\cdots A_1-J_K\right\|_F^2
 +\left\|B_s\cdots B_1-J_{\binom K2}\right\|_F^2
 \le K^{-8}.
\]
Here the squared matrix norm means the sum of its squared entries,
so this is an exact rational test. It implies both operator bounds
used in the proof. Enumerating lists eventually finds one, by the
averaging argument already proved.

Alternatively, generate independent uniformly random pairings and
repeat this test. The expected squared error is less than $K^{-10}$,
so Markov's inequality gives success probability at least $1-K^{-2}$
on each trial. The matrices have polynomial dimension; each row of a
pair-averaging matrix has at most four possible images. Thus this
particular search has an expected polynomial-time implementation.

For a fixed pairing, the two-color paths and cycles are found by
following edges. The possible path switches are binary choices. Their
moment rows are explicit finite sums, so the Harris--Srinivasan resampling procedure chooses the switches. Here the binary state spaces and
the moment coefficients can also be used directly with the
Harris--Srinivasan algorithm; no conditioned whole-list space occurs.


\subsection{Putting the procedures together}
Use the approximate local packings above in the simple-graph algorithm.
The two multigraph algorithms construct their feasible local mixtures
by the bounded-load and size-class procedures. The compression in
Lemma~\ref{lem:profiles} scans explicitly given slot sizes and assigns
fractions greedily, so it takes polynomial time for fixed accuracy.

The profile choices use the resampling procedure described above.
The matching of large-edge options uses
Proposition~\ref{prop:algorithmic-rank-three-matching}: its job degrees
exceed its slot degrees by a fixed fraction, exactly as in
\eqref{eq:degrees}. The list of option triples has polynomial size.
Indeed its multiplicity scale is $\lceil n\rceil^6$, and the parameters
$n=\lceil\Delta_w(G)\rceil$ or an approximate local bin count are
bounded by a constant times the number of input edges, apart from a
constant depending on accuracy.

The small-edge choices use Lemma~\ref{lem:two}, with exact rational
linear algebra. Sorting, grouping parallel edges, finding ordinary
matchings, proper bipartite edge coloring, and the Euler-tour
construction also take polynomial time. Each random procedure has an
expected polynomial number of resampling steps and can check its
output. Thus the complete algorithms have expected polynomial running
time for every fixed accuracy.

Choose a fixed threshold large enough for the inequalities in these
proofs. Below that threshold use decreasing-weight greedy cleanup
directly, before invoking any resampling procedure. Its color count is
then bounded by a constant depending only on accuracy. This gives an
algorithm on every input and accounts for the additive constants in
the fixed-accuracy statements.

The asymptotic $o(n)$ notation describes the family of fixed-accuracy
guarantees. We do not claim that the running time is polynomial jointly
in the input length and the reciprocal of the accuracy.

\Needspace{10\baselineskip}
\section{Finding disjoint options efficiently}\label{sec:matching-algorithm}
The following implementation is for the particular matching problem
used here. A job belongs to $A$ and a slot belongs to $B$; every option
contains one job and two distinct slots. A matching is a collection of
options that use no element twice. Degrees count the options containing
an element, including repetitions. The pair count $\Delta_2$ is the
largest number of options containing two fixed elements.

\begin{lemma}[A small regular extension for three-element edges]
\label{lem:rank-three-regularization}
Let $G=(A,B)$ be a family of triples: every
edge consists of one vertex in $A$ and two distinct vertices in $B$.
Suppose every vertex of $A$ has degree $a$, every vertex of $B$ has
degree at most $b$, and $a\ge b\ge1$ are integers. In polynomial
time, one can construct a family of triples
$G'=(A',B')$ containing $G$ as an induced subhypergraph such that
\begin{itemize}
\item every vertex of $A'$ has degree $a$ and every vertex of $B'$
has degree $b$;
\item no new edge meets an old vertex of $A$;
\item the maximum pair count is at most
$\max\{1,\Delta_2(G)\}$;
\item fewer than $4a$ vertices are added to $B$, and at most
$|B|/2+2b$ vertices are added to $A$.
\end{itemize}
The polynomial bound refers to the explicitly listed hypergraph.
\end{lemma}

\begin{proof}
The total degree deficit in $B$ is
\[
 T=\sum_{v\in B}(b-d_G(v))=b|B|-2a|A|.
\]
Let $g=\gcd(b,2a)$. Choose an integer $N$ with
\[
 2a\le N<4a,\qquad N\equiv-|B|\pmod{2a/g}.
\]
Such an $N$ exists because $2a/g\le2a$. Add $N$ new vertices to $B$,
each with deficit $b$. Then the total deficit is divisible by $2a$.
Put
\[
 m=\frac{T+Nb}{2a}.
\]
We have $b\le m\le |B|/2+2b$. Add $m$ new vertices to $A$.

First find a simple bipartite graph $J$ between the new $A$ vertices
and all the $B$ vertices. Every new $A$ vertex must have degree $2a$;
each $B$ vertex must have degree equal to its deficit, denoted $d_v$.
These requirements are feasible: $0\le d_v\le b\le m$ and
$\sum_vd_v=2am$. For completeness, use a network with capacity $2a$
from the source to each new $A$ vertex, capacity one on every possible
new-$A$--$B$ edge, and capacity $d_v$ from $v$ to the sink. For
$0\le k\le m$,
\[
 \sum_v\min\{d_v,k\}\ge\frac{k}{m}\sum_vd_v=2ak.
\]
These inequalities verify every cut condition for flow $2am$.
An integer maximum flow therefore supplies $J$ in polynomial time.

Process the new $A$ vertices one at a time. A new vertex $u$ has
$2a$ distinct neighbours in $J$. Pair those neighbours, forbidding
any pair already contained in an edge of the growing hypergraph.
Before processing $u$, each of its neighbours still needs the one
incident edge supplied at this step. Its present degree is therefore
at most $b-1$, so it has at most $b-1$ forbidden partners. The allowed
graph on the $2a$ neighbours has minimum degree at least
\[
 2a-1-(b-1)=2a-b\ge a.
\]
It has a perfect matching. One elementary proof takes a largest
matching: if $x,y$ are unmatched, they are nonadjacent, and the sum
of their degrees is at least $2a$. But unless a three-edge augmenting
path exists, each matched pair accounts for at most two of these
incidences, giving a sum at most $2a-2$, a contradiction. A standard
polynomial matching algorithm finds the required pairing.

For every paired pair $v,w$, add the hyperedge $\{u,v,w\}$. This
gives $u$ exactly $a$ hyperedges and uses each of its incidence edges
in $J$ once. After all new vertices have been processed, every deficit
is filled. A pair consisting of a new $A$ vertex and a $B$ vertex
occurs at most once, since $J$ is simple. A pair of $B$ vertices is
never reused. Pairs meeting old $A$ vertices are unchanged. This proves
the pair count claim, and the construction adds only polynomially many
vertices and edges.
\end{proof}

% Uses rank-three-polynomial-regularization.tex and the usual
% Moser--Tardos variable-model local lemma (arXiv:0903.0544, Theorem 1.2).
% Independently audited by bp199_tiny against the original DP text.
\begin{proposition}[Finding the matching with fixed degree slack]
\label{prop:algorithmic-rank-three-matching}
Fix $\eta,\zeta>0$. For sufficiently large integer $D$, an explicitly
listed family of triples $G=(A,B)$ satisfying
\[
 d(a)\ge(1+\eta)D\quad(a\in A),\qquad
 d(b)\le D\quad(b\in B),\qquad
 \Delta_2(G)\le D^{1-\zeta}
\]
has a matching covering every element of $A$ that can be found by a randomized algorithm in expected polynomial time. The degree threshold and running-time
constants may depend on $\eta,\zeta$. Each edge has exactly one vertex
in $A$ and two distinct vertices in $B$.
\end{proposition}

\begin{proof}
If $A$ is empty, return the empty matching. We henceforth assume it is
nonempty, so $D$ is at most the number of explicitly listed input edges.
We describe how to make the no-forbidden-configuration case of the
Delcourt--Postle argument algorithmic under the fixed slack used here.
The only concentration input is their iterative bipartite lemma and
its variable-model local-lemma proof~\cite[Lemma~3.3 and Sections~5.5--6]{DP}.
We use the standard Moser--Tardos algorithm: sample all independent
variables and, while a bad event holds, resample its variables. Under
the symmetric local-lemma bound $ep(d+1)<1$, its expected number of
resamplings is polynomial in the number of explicitly checkable events.

Choose fixed rational $\theta$ with
$0<\theta<\min\{\eta/4,1/8\}$. Delete edges at vertices of $A$
until each has degree $a_*:=\lceil(1+4\theta)D\rceil$. This is
possible for large $D$, and all degrees are then at most $2D$.

\paragraph{Reducing the degrees once.}
Choose a fixed rational $s$ with $0<s<\min\{\zeta/2,1/2\}$ and
put $p=2^{-\lceil(1-s)\log_2D\rceil}$. This is an exactly
computable dyadic probability; equivalently its exponent can be found
using integer powers, since $s$ is rational. Write $d=pD$.
Independently retain each edge with probability $p$.

Let $\delta=\theta/10$. The bad events are a retained degree below
$(1-\delta)pa_*$ at a vertex of $A$, a retained degree above
$(1+\delta)d$ at a vertex of $B$, and a pair count above
\[
 t=\left\lfloor\frac{\lfloor\log_2d\rfloor^2}{4}\right\rfloor.
\]
For large $D$, $t\le0.6(\log d)^2$. Chernoff bounds give degree-event
probabilities $\exp(-\Omega_{\theta}(d))$. A pair has expected
pair count at most $pD^{1-\zeta}\le D^{s-\zeta}$, which tends to zero;
the standard binomial upper-tail bound gives probability at most
$\exp(-\Omega_{s,\zeta}((\log D)^3))$ for a pair event.

Each event involves at most $2D$ edge variables. One edge variable
belongs to at most three degree events and three pair events.
Consequently the dependency degree is $O(D)$, and the symmetric
Moser--Tardos condition holds for large $D$. All events are checked
by counting retained incidences. This finds a retained graph in
expected polynomial time. Set
\[
 a_0=\lfloor(1-\delta)pa_*\rfloor,
 \qquad b_0=\lceil(1+\delta)d\rceil.
\]
For large $D$, its $A$ degrees are at least $a_0$, its $B$ degrees
are at most $b_0$, and $a_0/b_0\ge1+2\theta$. Its pair counts
are at most $0.6(\log d)^2$.

\paragraph{Choosing a small random matching.}
Suppose the current degree bounds are integers $a\ge b$, with
$1+\theta\le R:=a/b\le2000$, and pair counts at most
$(\log b)^2$. Fix rational constants
\[
 0<\gamma\le\min\{\theta/8,1/1000\},\qquad
 \varepsilon=\gamma/10^{12},\qquad w=1.
\]
These satisfy the inequalities in the Delcourt--Postle iterative
lemma for $r=3$, $g=2$, and $R\le2000$, once $b$ is large enough.
In particular,
\[
 \gamma\ge 96R(15R+3)\varepsilon,
 \qquad \varepsilon\ge b^{-1/18},\qquad R\le\log b.
\]
The last two conditions hold for all sufficiently large $b$ because
the parameters are fixed.

Trim the degrees in $A$ to $a$, and apply
Lemma~\ref{lem:rank-three-regularization} to make the $B$ degrees
exactly $b$. This extension has polynomial size and preserves the
pair-count bound. Independently select each edge with probability
$\varepsilon/b$, and independently mark it for deletion with probability
$\varepsilon$; these are the equalizing coins of the small random matching step when the
forbidden family is empty and $w=1$. Keep the selected edges that
meet no other selected edge and were not marked for deletion. They
form a matching. Delete its vertices and the edges marked for deletion.

The degree conclusions of the cited iterative lemma are
\[
 a'=a e^{-\varepsilon(3+\gamma)},\qquad
 b'=b e^{-\varepsilon(R+2-\gamma)}.
\]
We briefly specify why the probabilistic step is algorithmic, rather
than relying only on its positive probability. Its bad events are
the failed residual-degree bounds at the surviving vertices.
They are explicitly checkable functions of independent edge coins
within a fixed distance in the incidence graph. The concentration
calculation in the cited proof bounds their probabilities by
$\exp(- (\log b)^{3/2})$, while their dependency degrees are bounded
by a fixed power of $b$. With the forbidden family empty, no other
configuration events or probability-evaluation subroutines are needed.
Moser--Tardos therefore finds acceptable coins in expected polynomial
time. We may use weaker, rational degree thresholds, as specified next;
their bad events are subsets of these events and satisfy the same bounds.

\paragraph{Exact thresholds and growth of the degree ratio.}
Put $\tau=\varepsilon\theta/100$. Compute rational numbers $L,U$
such that
\begin{align*}
 (1-2\tau)e^{-\varepsilon(3+\gamma)}
 &\le L\le(1-\tau)e^{-\varepsilon(3+\gamma)},\\
 (1+\tau)e^{-\varepsilon(R+2-\gamma)}
 &\le U\le(1+2\tau)e^{-\varepsilon(R+2-\gamma)}.
\end{align*}
Use the integer thresholds $a_+=\lfloor aL\rfloor$ and
$b_+=\lceil bU\rceil$ in the event tests. These relax the preceding
degree bounds. Taylor expansion with a certified remainder computes
the required exponential intervals in polynomial time; only a fixed
relative precision depending on $\theta$ is required. All actual
coin probabilities and event comparisons remain rational.

For large $b$, the integer rounding changes the relative bounds by
at most another $\tau$. Hence
\[
 \frac{a_+}{b_+}
 \ge R e^{\varepsilon(R-1-2\gamma)}
          \frac{1-3\tau}{1+3\tau}
 \ge R e^{\varepsilon\theta/2}.
\]
The ratio thus reaches $1000$ after at most
$\lceil2\log(1000)/(\varepsilon\theta)\rceil+1$ steps, a constant
depending only on $\eta$. Before and immediately after the last step
it remains below $2000$. Throughout these constantly many steps,
$b\ge c_{\eta}d$ for a positive constant $c_{\eta}$. Thus the
original pair-count bound $0.6(\log d)^2$ is at most
$(\log b)^2$ for large $D$, and every step's hypotheses hold.

\paragraph{Discarding the regularization after each step.}
This is essential for the size bound. After each step retain only
the matching edges and residual vertices belonging to the current
original graph; discard all dummy vertices. No regularization edge
meets an old $A$ vertex. Consequently an old $A$ vertex is removed
only when an original edge matches it, and the degree of each surviving
old $A$ vertex is unchanged by discarding dummy vertices. Degrees
in old $B$ can only decrease. Old $B$ vertices used by a discarded
dummy matching edge may simply be left unavailable in later steps.
Therefore the next current graph is always a subhypergraph of the
initial reduced graph. The extensions never accumulate or multiply.
Each individual extension has $O(|V(G)|+D)$ vertices and polynomially
many edges, and there are only constantly many extensions.

\paragraph{Finishing.}
Once $a\ge1000b$, trim each remaining job to exactly $900b$ options.
Independently choose one uniformly for each job. A bad event specifies
two options from different jobs that share a slot and are both chosen.
Its probability is $(900b)^{-2}$. A job belongs to at most
$(900b)(2b)$ such events, so an event has dependency degree at most
$3600b^2$. The symmetric local-lemma inequality holds with ample slack:
\[
 e(900b)^{-2}(3600b^2+1)<1.
\]
Another Moser--Tardos run finds a collision-free choice. Together with
the original matching edges retained from the previous steps, it
covers every original job exactly once.

All input hyperedges, including repeated copies, are explicitly listed.
After the first trimming their number, the numbers of bad events,
the sizes of the regularized graphs, and the number of coin variables
are polynomial in that input size. Each event test and resampling uses
polynomial time. The parameters are fixed, all actual coin probabilities
are rational with polynomial encoding length, and the number of stages
is independent of input size. The standard expected resampling bounds
therefore give expected polynomial total time. Every returned object is
checked to be a matching covering every element of $A$, so every output is correct.
\end{proof}

% Complete inline bibliography for the pasted full manuscript.
% Fifteen entries recovered verbatim from resident prior TeX sources (citation keys remapped).
% Five missing entries verified against publisher/original-paper DOI pages.
% This file is a fragment for insertion, not a standalone document.
\begin{thebibliography}{99}

\bibitem{CG07}
J.~R. Correa and M.~X. Goemans,
\emph{Improved bounds on nonblocking 3-stage Clos networks},
SIAM J. Comput. \textbf{37} (2007), no.~3, 870--894.
\href{https://doi.org/10.1137/060656413}{doi:10.1137/060656413}.

\bibitem{CR91}
S.-P. Chung and K.~W. Ross,
\emph{On nonblocking multirate interconnection networks},
SIAM J. Comput. \textbf{20} (1991), no.~4, 726--736.
\href{https://doi.org/10.1137/0220046}{doi:10.1137/0220046}.

% Verified: Wiley original publication.
\bibitem{Clos53} Charles Clos.
A study of non-blocking switching networks.
\emph{Bell System Technical Journal} 32(2), pages 406--424, 1953.
\href{https://doi.org/10.1002/j.1538-7305.1953.tb01433.x}{doi:10.1002/j.1538-7305.1953.tb01433.x}.

% Verified: SIAM original publication.
\bibitem{DGHK98} D. Z. Du, B. Gao, F. K. Hwang, and J. H. Kim.
On multirate rearrangeable Clos networks.
\emph{SIAM Journal on Computing} 28(2), pages 463--470, 1998.
\href{https://doi.org/10.1137/S0097539795284716}{doi:10.1137/S0097539795284716}.

\bibitem{DP} Michelle Delcourt and Luke Postle.
Finding an almost perfect matching in a hypergraph avoiding forbidden
submatchings.
\href{https://arxiv.org/html/2204.08981v3}{arXiv:2204.08981v3}, 2024.

\bibitem{Edmonds65} Jack Edmonds.
Maximum matching and a polyhedron with 0,1-vertices.
\emph{Journal of Research of the National Bureau of Standards,
Section B} 69B, pages 125--130, 1965.
\href{https://doi.org/10.6028/jres.069b.013}{doi:10.6028/jres.069b.013}.

\bibitem{Epstein} Leah Epstein.
Tighter bounds for the harmonic bin packing algorithm.
\emph{European Journal of Operational Research} 316(1), pages 72--84, 2024.
\href{https://doi.org/10.1016/j.ejor.2024.01.051}{doi:10.1016/j.ejor.2024.01.051}.

\bibitem{FS} Uriel Feige and Mohit Singh.
Edge coloring and decompositions of weighted graphs.
In \emph{Algorithms--ESA 2008}, Lecture Notes in Computer Science 5193,
pages 405--416, Springer, 2008.
\href{https://doi.org/10.1007/978-3-540-87744-8_34}{doi:10.1007/978-3-540-87744-8\_34}.

\bibitem{GG61} Paul C. Gilmore and Ralph E. Gomory.
A linear programming approach to the cutting-stock problem.
\emph{Operations Research} 9(6), pages 849--859, 1961.
\href{https://doi.org/10.1287/opre.9.6.849}{doi:10.1287/opre.9.6.849}.

\bibitem{HS} David G. Harris and Aravind Srinivasan.
The Moser--Tardos framework with partial resampling.
\emph{Journal of the ACM} 66(5), Article 36, 2019.
\href{https://arxiv.org/html/1406.5943v5}{arXiv:1406.5943v5}.

% Verified: publisher DOI with indexed journal metadata.
\bibitem{Huc11} Florian Huc.
Weighted-edge-coloring of $k$-degenerate graphs and bin-packing.
\emph{Journal of Interconnection Networks} 12(1--2), pages 109--124, 2011.
\href{https://doi.org/10.1142/S0219265911002861}{doi:10.1142/S0219265911002861}.

\bibitem{KS15} Arindam Khan and Mohit Singh.
On weighted bipartite edge coloring.
In \emph{35th IARCS Annual Conference on Foundations of Software Technology
and Theoretical Computer Science (FSTTCS 2015)}, LIPIcs 45, pages 136--150,
Schloss Dagstuhl, 2015.
\href{https://doi.org/10.4230/LIPIcs.FSTTCS.2015.136}{doi:10.4230/LIPIcs.FSTTCS.2015.136}.

% Verified: Springer original publication.
\bibitem{Konig16} D\'enes K\H{o}nig.
\"Uber Graphen und ihre Anwendung auf Determinantentheorie und Mengenlehre.
\emph{Mathematische Annalen} 77, pages 453--465, 1916.
\href{https://doi.org/10.1007/BF01456961}{doi:10.1007/BF01456961}.

\bibitem{MT} Robin A. Moser and G\'abor Tardos.
A constructive proof of the general Lov\'asz Local Lemma.
\emph{Journal of the ACM} 57(2), Article 11, 2010.
\href{https://doi.org/10.1145/1667053.1667060}{doi:10.1145/1667053.1667060}.

% Verified: SIAM original publication.
\bibitem{MT89} Riccardo Melen and Jonathan S. Turner.
Nonblocking multirate networks.
\emph{SIAM Journal on Computing} 18(2), pages 301--313, 1989.
\href{https://doi.org/10.1137/0218021}{doi:10.1137/0218021}.

\bibitem{NV} Hung Q. Ngo and Van H. Vu.
Multirate rearrangeable Clos networks and a generalized edge-coloring
problem on bipartite graphs.
\emph{SIAM Journal on Computing} 32(4), pages 1040--1049, 2003.
\href{https://doi.org/10.1137/S0097539702408235}{doi:10.1137/S0097539702408235}.

\bibitem{SS} Peter Sanders and David Steurer.
An asymptotic approximation scheme for multigraph edge coloring.
\emph{ACM Transactions on Algorithms} 4(2), Article 21, 2008.
\href{https://doi.org/10.1145/1361192.1361198}{doi:10.1145/1361192.1361198}.

\bibitem{Sannyasi} Debarsho Sannyasi.
Improved approximation algorithms for weighted edge coloring of graphs.
Preprint, 2020.
\href{https://arxiv.org/abs/2012.15056}{arXiv:2012.15056}.

\bibitem{Shannon49} Claude E. Shannon.
A theorem on coloring the lines of a network.
\emph{Journal of Mathematics and Physics} 28, pages 148--151, 1949.
\href{https://doi.org/10.1002/sapm1949281148}{doi:10.1002/sapm1949281148}.

\bibitem{Vizing64}
V.~G. Vizing,
\emph{On an estimate of the chromatic class of a \(p\)-graph},
Diskret. Analiz \textbf{3} (1964), 25--30 (in Russian).

\end{thebibliography}
\end{document}

